% Generated mechanically from frozen input p07/ko/models.tex.
% Do not edit this generated file; edit the bound locale source instead.

% OK, start here.
%

\setcounter{chapter}{54}
\chapter{반안정 축약}



\phantomsection
\label{models-section-phantom}


\section{서론}
\label{models-section-introduction}

\noindent
이 장에서는 곡선에 대한 반안정 축약 정리를 증명한다.
Artin과 Winters가 논문 \cite{Artin-Winters}에서 제시한 방법을
사용할 것이다.

\medskip\noindent
실제로 비특이화에 관한 어떤 결과도 사용하지 않고 곡선의 반안정 축약 정리를
증명할 수 있다. 즉 Mumford가 발전시킨 기하적 불변식론(GIT)을 사용하여
반안정 곡선 모듈라이의 존재성과 사영성을 확립하는 방법이 있다
(\cite{GIT} 참조). 이 방법을 적극적으로 발전시킨 Gieseker는 강의록
\cite{Gieseker}\footnote{Gieseker의 강의록은 대수적으로 닫힌 체 위에서
쓰였지만, 같은 방법은 $\mathbf{Z}$ 위에서도 작동한다.}에서 완전한 결과를
증명하였다. 이는 상당히 놀라운 성과이다. 양의 차원 기저 위의 곡선족을
실제로 연구하지 않고도 이런 결과를 증명할 수 있다는 점은 다소 직관에
어긋나 보이기 때문이다.

\medskip\noindent
역사적으로 곡선에 대한 반안정 축약 정리의 최초 증명은 Deligne과 Mumford의
논문 \cite{DM}에서 찾을 수 있다. 이 증명은 문제를 아벨 다양체의 경우로
환원하는데, 그 경우는 Grothendieck 등의 연구 덕분에 당시 이미 알려져
있었다(\cite{SGA7-I}, \cite{SGA7-II} 참조). 아벨 다양체의 반안정 축약
정리는 N\'eron 모형 이론을 사용하는데, 이 이론은 다시 기저 위의 쌍유리
군 법칙에 대한 연구에 의존한다.

\medskip\noindent
Artin과 Winters 논문의 방법은 곡면 특이점의 비특이화를 이용하여 정칙
모형을 얻는다. 정칙 모형의 존재가 주어지면, 증명은 특수 올의 가능한
형태를 분석한 뒤 체 위의 $1$차원 스킴의 Picard 군에서 꼬임에 관한
부등식을 사용하여 결론을 내리는 것으로 이루어진다. 비슷한 논증은
Saito의 논문 \cite{Saito}에서도 볼 수 있다. 그는 에탈 코호몰로지를
직접 사용하며, 관성군이 $\ell$-진 에탈 코호몰로지에 작용하는 방식으로
반안정 축약을 특징지을 수 있다는 더 강한 결과를 얻는다.

\medskip\noindent
이 정리를 증명하는 또 다른 접근은 강직 해석기하의 기법을 사용하는 것이다.
이에 관해서는 \cite{vanderPut} 및 \cite{Arzdorf-Wewers}를 참조하라.

\medskip\noindent
Temkin의 논문 \cite{Temkin}은 값매김 이론의 기법을 사용하며 그 밖에도
훨씬 많은 내용을 증명한다. 또한 이 논문의 부록 A에는 여러 증명과
$2$차원 스킴의 비특이화 사이의 관계가 잘 정리되어 있다.

\medskip\noindent
독자가 참고할 만한 또 다른 개관 논문으로 Ahmed Abbes의
\cite{Abbes-ssr}가 있다.






\section{선형대수}
\label{models-section-linear-algebra}

\noindent
뒤에서 사용할 몇 가지 보조정리를 제시한다.

\begin{lemma}
\label{models-lemma-recurring}
\begin{reference}
\cite[정리 I]{Taussky}
\end{reference}
$A = (a_{ij})$를 복소 $n \times n$ 행렬이라 하자.
\begin{enumerate}
\item 모든 $i$에 대해 $|a_{ii}| > \sum_{j \not = i} |a_{ij}|$이면
$\det(A)$는 영이 아니다.
\item $m_i > 0$이고 모든 $i$에 대해
$|a_{ii} m_i| > \sum_{j \not = i} |a_{ij}m_j|$를 만족하는 실벡터
$m = (m_1, \ldots, m_n)$가 존재하면 $\det(A)$는 영이 아니다.
\end{enumerate}
\end{lemma}

\begin{proof}
$A$가 (1)의 조건을 만족하고 $\det(A) = 0$이라 하자. 그러면 $Az=0$인
영이 아닌 벡터 $z$가 존재한다. $|z_r|$가 최대가 되도록 $r$을 택하면
$$
|a_{rr} z_r| = |\sum\nolimits_{k \not = r} a_{rk}z_k| \leq
\sum\nolimits_{k \not = r} |a_{rk}||z_k| \leq
|z_r| \sum\nolimits_{k \not = r} |a_{rk}| < |a_{rr}||z_r|
$$
이고 이는 모순이다. (2)를 증명하려면 행렬식이
$m_1 \ldots m_n \det(A)$인 행렬 $(a_{ij}m_j)$에 (1)을 적용하면 된다.
\end{proof}

\begin{lemma}
\label{models-lemma-recurring-real}
$A = (a_{ij})$를 $i \not = j$일 때 $a_{ij} \geq 0$인 실
$n \times n$ 행렬이라 하고, $m_i>0$인 실벡터
$m = (m_1, \ldots, m_n)$를 택하자. $I \subset \{1, \ldots, n\}$에 대해
$x_I \in \mathbf{R}^n$를 $i \in I$이면 $i$번째 좌표가 $m_i$이고
그렇지 않으면 $0$인 벡터라 하자. 만일
\begin{equation}
\label{models-equation-ineq}
-a_{ii}m_i \geq \sum\nolimits_{j \not = i} a_{ij}m_j
\end{equation}
가 모든 $i$에 대해 성립하면, $\Ker(A)$는 다음 조건을 만족하는 벡터
$x_I$들이 생성하는 벡터공간이다.
\begin{enumerate}
\item $i \in I$, $j \not \in I$이면 $a_{ij}=0$이고,
\item $i \in I$이면 (\ref{models-equation-ineq})에서 등호가 성립한다.
\end{enumerate}
\end{lemma}

\begin{proof}
$a_{ij}$를 $a_{ij}m_j$로 바꾸어 모든 $i$에 대해 $m_i=1$이라 가정해도
된다. (1), (2)를 만족하는 $I \subset \{1, \ldots, n\}$에 대해서는
간단한 계산으로 $x_I$가 $A$의 핵에 속함을 알 수 있다. 반대로
$x = (x_1, \ldots, x_n) \in \mathbf{R}^n$를 $A$의 핵에 속하는 영이 아닌
벡터라 하자. $x$의 영이 아닌 좌표 수에 대한 귀납법으로 $x$가 (1),
(2)를 만족하는 $x_I$들의 생성공간에 속함을 보이겠다.
$I \subset \{1, \ldots, n\}$를 $|x_r|$가 최대인 지표 $r$들의 집합이라
하자. $r \in I$에 대해
$$
|a_{rr} x_r| = |\sum\nolimits_{k \not = r} a_{rk}x_k| \leq
\sum\nolimits_{k \not = r} a_{rk}|x_k| \leq
|x_r| \sum\nolimits_{k \not = r} a_{rk} \leq |a_{rr}||x_r|
$$
이다. 따라서 모든 곳에서 등호가 성립한다. 특히 $r \in I$이고
$k \not \in I$이면 $a_{rk}=0$이며, $r \in I$에 대해
(\ref{models-equation-ineq})에서 등호가 성립한다. 이제 $x$에서 $x_I$의
적절한 배수를 빼면 영이 아닌 좌표의 수를 줄일 수 있다.
\end{proof}

\begin{lemma}
\label{models-lemma-recurring-symmetric-real}
$A = (a_{ij})$를 $i \not = j$일 때 $a_{ij}\geq0$인 대칭 실
$n\times n$ 행렬이라 하자. $m_i>0$인 실벡터
$m=(m_1,\ldots,m_n)$를 택하고 다음을 가정하자.
\begin{enumerate}
\item $Am = 0$,
\item $i \in I$, $j \not \in I$이면 $a_{ij}=0$이 되는 공집합이 아닌
진부분집합 $I \subset \{1, \ldots, n\}$는 존재하지 않는다.
\end{enumerate}
그러면 $x^tAx\leq0$이고, 등호가 성립할 필요충분조건은 어떤
$q \in \mathbf{R}$에 대해 $x=qm$인 것이다.
\end{lemma}

\begin{proof}[첫째 증명]
$a_{ij}$를 $a_{ij}m_im_j$로 바꾸어 모든 $i$에 대해 $m_i=1$이라
가정해도 된다. 조건 (1)은 모든 $i$에 대해
$-a_{ii}=\sum_{j\not=i}a_{ij}$라는 뜻이다.
$x^tAx=\sum_{i,j}x_ia_{ij}x_j$임을 상기하면
\begin{align*}
\sum\nolimits_{i \not = j} -a_{ij}(x_j - x_i)^2 & =
\sum\nolimits_{i \not = j} -a_{ij}x_j^2 + 2a_{ij}x_ix_i - a_{ij}x_i^2 \\
& =
\sum\nolimits_j a_{jj} x_j^2 +
\sum\nolimits_{i \not = j} 2a_{ij}x_ix_i +
\sum\nolimits_j a_{jj} x_i^2 \\
& = 2x^tAx
\end{align*}
이고, 이는 분명 $\leq0$이다. 등호가 성립한다고 하자. $I$를
% P07-KO55-ERR-001: frozen source appears to require I empty here, but says I is the full set; translated canonically without silent repair.
$x_i \not = x_1$인 지표 $i$들의 집합이라 하면 $i \in I$,
$j \not \in I$일 때 $a_{ij}=0$이다. 따라서 조건 (2)에 의해
$I = \{1, \ldots, n\}$이어야 하므로
$x$는 $m=(1,\ldots,1)$의 배수이다.
\end{proof}

\begin{proof}[둘째 증명]
스펙트럼 정리에 의해 행렬 $A$의 고윳값은 실수이다. 모든 고윳값이
$\leq0$임을 보이자. 조건 (1)은 모든 $i$에 대해
$-a_{ii}m_i=\sum_{j\not=i}a_{ij}m_j$라는 뜻이므로, $\lambda>0$일 때
행렬 $A'=A-\lambda I$는 모든 $i$에 대해
$|a'_{ii}m_i|>\sum a'_{ij}m_j=\sum|a'_{ij}m_j|$를 만족한다.
따라서 보조정리 \ref{models-lemma-recurring}에 의해 $A'$는 가역이다.
그러므로 대칭 쌍선형 형식 $x^tAy$는 반음의 정부호, 즉 모든 $x$에 대해
$x^tAx\leq0$이다. 이에 따라 $A$의 핵은 $x^tAx=0$인 벡터 $x$들의
집합과 같다. 보조정리 \ref{models-lemma-recurring-real}의 핵에 대한 기술로
마지막 주장을 얻는다.
\end{proof}

\begin{lemma}
\label{models-lemma-orthogonal-direct-sum}
정수값 대칭 쌍선형 양의 정부호 형식
$\langle\ ,\ \rangle:L\times L\to\mathbf Z$가 주어진 유한 자유
$\mathbf Z$-가군 $L$을 생각하자. $L/A$가 꼬임 없는 부분가군
$A\subset L$을 택하고
$B = \{b \in L \mid \langle a, b\rangle = 0,\ \forall a \in A\}$.
라 두자. 그러면 단사 사상
$$
A^\#/A \leftarrow L/(A \oplus B) \rightarrow B^\#/B
$$
이 있고, 그 여핵들은 $L^\#/L$의 몫이다. 여기서
$A^\# = \{a' \in A \otimes \mathbf{Q} \mid
\langle a, a'\rangle \in \mathbf{Z},\ \forall a \in A\}$
이며 $B$와 $L$에 대해서도 마찬가지로 정의한다.
\end{lemma}

\begin{proof}
$A$ 위에서 형식이 비퇴화이므로(양의 정부호성에 의해)
$L \otimes \mathbf{Q} = A \otimes \mathbf{Q} \oplus B \otimes \mathbf{Q}$
이다. 사영을 $\pi_B:L\otimes\mathbf Q\to B\otimes\mathbf Q$라 쓰자.
형식이 정수값이므로 $x\in L$에 대해 $\pi_B(x)\in B^\#$이다.
따라서 완전열
$$
0 \to A \to L \xrightarrow{\pi_B} B^\# \to Q \to 0
$$
을 얻는다. 여기서 $Q$는 $L\to B^\#$의 여핵이다. 사상
$L^\#\to B^\#$는 $\mathbf Z$-가군 사상으로 분할되는 $B\to L$의
$\mathbf Z$-선형 쌍대이므로 전사이다. 따라서 $Q$는 $L^\#/L$의 몫이다.
$A\oplus B$로 나누면 짧은 완전열
$$
0 \to L/(A \oplus B) \to B^\#/B \to Q \to 0
$$
을 얻으며, 이것으로 보조정리가 증명된다.
\end{proof}

\begin{lemma}
\label{models-lemma-coker}
정수값 대칭 쌍선형 양의 정부호 형식
$\langle\ ,\ \rangle:L_i\times L_i\to\mathbf Z$가 주어진 유한 자유
$\mathbf Z$-가군 $L_0$, $L_1$을 생각하자. $\langle\ ,\ \rangle$에 관해
서로 수반인 사상
$\text d:L_0\to L_1$, $\text d^*:L_1\to L_0$이 주어졌다고 하자.
$L_0$ 위의 형식이 유니모듈러이면 동형
$$
\Phi :
\Coker(\text{d}^*\text{d})_{torsion}
\longrightarrow
\Im(\text{d})^\#/\Im(\text{d})
$$
이 존재한다. 표기는 보조정리 \ref{models-lemma-orthogonal-direct-sum}와 같다.
\end{lemma}

\begin{proof}
$x\in L_0$가 $\Coker(\text d^*\text d)$의 꼬임류를 나타내는 원소라
하자. 어떤 $a>0$에 대해 $ax=\text d^*\text d(y)$로 쓸 수 있다.
$z=\text d(y')$인 임의의 $z\in\Im(\text d)$에 대해
$$
\langle (1/a)\text{d}(y), z \rangle =
\langle (1/a)\text{d}(y), \text{d}(y') \rangle =
\langle x, y' \rangle \in \mathbf{Z}
$$
이다. 따라서 $(1/a)\text d(y)\in\Im(\text d)^\#$이다.
$\Phi(x)=(1/a)\text d(y)\bmod\Im(\text d)$로 정의한다.
$\Phi$가 잘 정의되고 가법적이며 단사라는 증명은 생략한다.

\medskip\noindent
$\Phi$의 전사성을 보이기 위해 $z\in\Im(\text d)^\#$를 택하자.
$z$는 $x\mapsto\langle z,\text d(x)\rangle$에 의해 선형 사상
$L_0\to\mathbf Z$를 정의한다. 가정에 의해 $L_0$의 짝이
유니모듈러이므로, 모든 $x\in L_0$에 대해
$\langle x',x\rangle=\langle z,\text d(x)\rangle$인 $x'\in L_0$를
찾을 수 있다. 특히 $x'$는 $\Ker(\text d)$와의 짝이 영이다.
$\Im(\text d^*\text d)\otimes\mathbf Q$가
$\Ker(\text d)\otimes\mathbf Q$의 직교여공간이므로, $x'$는
$\Coker(\text d^*\text d)$의 꼬임류를 정의한다. $\Phi(x')=z$임을
보이자. 어떤 $y\in L_0$, $a>0$에 대해
$ax'=\text d^*\text d(y)$로 쓰면, 임의의 $x\in L_0$에 대해
$$
\langle z, \text{d}(x)\rangle =
\langle x', x \rangle =
\langle (1/a)\text{d}^*\text{d}(y), x \rangle =
\langle (1/a)\text{d}(y),\text{d}(x) \rangle
$$
이다. 따라서 $z=\Phi(x')$이고 증명이 끝난다.
\end{proof}

\begin{lemma}
\label{models-lemma-recurring-symmetric-integer}
$A=(a_{ij})$를 $i\ne j$일 때 $a_{ij}\geq0$인 대칭 정수
$n\times n$ 행렬이라 하자. $m_i>0$인 정수벡터
$m=(m_1,\ldots,m_n)$를 택하고 다음을 가정하자.
\begin{enumerate}
\item $Am = 0$,
\item $i\in I$, $j\notin I$이면 $a_{ij}=0$이 되는 공집합이 아닌
진부분집합 $I \subset \{1, \ldots, n\}$는 없다.
\end{enumerate}
$e$를 $i<j$이고 $a_{ij}>0$인 순서쌍 $(i,j)$의 수라 하자.
모든 $a_{ij}$ 및 $m_i$와 서로소인 소수 $\ell$에 대해
$$
\dim_{\mathbf{F}_\ell}(\Coker(A)[\ell]) \leq 1 - n + e
$$
\end{lemma}

\begin{proof}
보조정리 \ref{models-lemma-recurring-symmetric-real}에 의해 $A$의 계수는
$n-1$이다. 합성
$$
\mathbf{Z}^{\oplus n} \xrightarrow{\text{diag}(m_1, \ldots, m_n)}
\mathbf{Z}^{\oplus n} \xrightarrow{(a_{ij})}
\mathbf{Z}^{\oplus n} \xrightarrow{\text{diag}(m_1, \ldots, m_n)}
\mathbf{Z}^{\oplus n}
$$
의 행렬은 $a_{ij}m_im_j$이다. $\ell$에 대한 가정에 의해 첫째와 마지막
사상의 여핵은 그 위수가 $\ell$과 서로소인 꼬임군이므로, 성분이
$a_{ij}m_im_j$인 행렬에 대해 보조정리를 증명하면 충분하다. 따라서
$m=(1,\ldots,1)$이라 가정해도 된다.

\medskip\noindent
$m=(1,\ldots,1)$이라 가정하자. $V=\{1,\ldots,n\}$ 및
$E=\{(i,j)\mid i<j\text{이고 }a_{ij}>0\}$라 두자.
$e=(i,j)\in E$에 대해 $a_e=a_{ij}$라 하고,
$s(i,j)=i$, $t(i,j)=j$로 사상 $s,t:E\to V$를 정의한다. 또한
$\mathbf{Z}(V) = \bigoplus_{i \in V} \mathbf{Z}i$이고
$\mathbf{Z}(E) = \bigoplus_{e \in E} \mathbf{Z}e$.
$\mathbf Z(V)$와 $\mathbf Z(E)$ 위의 대칭 양의 정부호 정수값 짝을
$$
\langle i, i \rangle = 1\text{ (모든 }i \in V\text{에 대해)}, \quad
\langle e, e \rangle = a_e\text{ (모든 }e \in E\text{에 대해)}
$$
로 정하고, 나머지 모든 짝은 영으로 둔다. 사상
$$
\text{d} : \mathbf{Z}(V) \to \mathbf{Z}(E), \quad
i \longmapsto
\sum\nolimits_{e \in E,\ s(e) = i} e - \sum\nolimits_{e \in E,\ t(e) = i} e
$$
및
$$
\text{d}^*(e) = a_e(s(e) - t(e))
$$
을 생각하자. 계산하면
$$
\langle d(x), y\rangle = \langle x, \text{d}^*(y) \rangle
$$
이다. 즉 $\text d$와 $\text d^*$는 서로 수반이다. 다음으로 계산하면
\begin{align*}
\text{d}^*\text{d}(i)
& = 
\text{d}^*(
\sum\nolimits_{e \in E,\ s(e) = i} e - \sum\nolimits_{e \in E,\ t(e) = i} e) \\
& =
\sum\nolimits_{e \in E,\ s(e) = i} a_e(s(e) - t(e)) -
\sum\nolimits_{e \in E,\ t(e) = i} a_e(s(e) - t(e))
\end{align*}
$\text d^*\text d(i)$에서 $i$의 계수는
$$
\sum\nolimits_{e \in E,\ s(e) = i} a_e +
\sum\nolimits_{e \in E,\ t(e) = i} a_e = - a_{ii}
$$
이다. 이는 $\sum_j a_{ij}=0$이기 때문이고,
$\text d^*\text d(i)$에서 $j\ne i$의 계수는 $-a_{ij}$이다.
따라서 $\Coker(A)=\Coker(\text d^*\text d)$이다.

\medskip\noindent
포함
$$
\Im(\text{d}) \oplus \Ker(\text{d}^*) \subset \mathbf{Z}(E)
$$
을 생각하자. 왼쪽은 직교 직합이다.
$\mathbf Z(E)/\Ker(\text d^*)$는 분명 꼬임이 없다.
$\mathbf Z(E)/\Im(\text d)$도 꼬임이 없음을 보이자.
$x=\sum x_e e\in\mathbf Z(E)$와 $a>1$이 있어 어떤
$y=\sum y_i i\in\mathbf Z(V)$에 대해 $ax=\text d y$라 하자.
그러면 $ax_e=y_{s(e)}-y_{t(e)}$이다. 조건 (2)에 의해 모든 $y_i$는
법 $a$에서 같은 합동류를 갖는다. 따라서
$y=ay'+(y_1,y_1,\ldots,y_1)$로 쓸 수 있다.
$\text d(y_1,y_1,\ldots,y_1)=0$이므로 $x=\text d(y')$이고, 이것이
보일 바였다.

\medskip\noindent
따라서 보조정리 \ref{models-lemma-orthogonal-direct-sum}을 적용하여 단사 사상
$$
\Im(\text{d})^\#/\Im(\text{d}) \leftarrow
\mathbf{Z}(E)/(\Im(\text{d}) \oplus \Ker(\text{d}^*)) \rightarrow
\Ker(\text{d}^*)^\#/\Ker(\text{d}^*)
$$
을 얻는다. 그 여핵들은 $a_e$들의 곱에 의해 소멸하며, 이 곱은
$\ell$과 서로소이다. $\Ker(\text d^*)$는 계수 $1-n+e$인 격자이므로
동형
$$
\Phi : M_{torsion} \longrightarrow \Im(\text{d})^\#/\Im(\text{d})
$$
의 존재만 보이면 증명이 끝난다. 이는 보조정리 \ref{models-lemma-coker}에서
증명하였다.
\end{proof}









\section{수치형}
\label{models-section-numerical-types}

\noindent
논증의 일부에서는 다음 자료구조의 조합론을 사용한다.

\begin{definition}
\label{models-definition-type}
{\it 수치형} $T$는 자료
$$
n, m_i, a_{ij}, w_i, g_i
$$
로 주어진다. 여기서 $n \geq 1$은 정수이고 $1 \leq i, j \leq n$에 대한
$m_i$, $a_{ij}$, $w_i$, $g_i$는 다음 조건을 만족하는 정수이다.
\begin{enumerate}
\item $m_i > 0$, $w_i > 0$, $g_i \geq 0$,
\item 행렬 $A = (a_{ij})$는 대칭이고 $i \not = j$이면 $a_{ij} \geq 0$이다.
\item $i \in I$, $j \not \in I$이면 $a_{ij} = 0$이 되는 진부분집합
$I \subset \{1, \ldots, n\}$는 존재하지 않는다.
\item 각 $i$에 대해 $\sum_j a_{ij}m_j = 0$이고,
\item $w_i | a_{ij}$.
\end{enumerate}
\end{definition}

\noindent
이는 다루기 다소 번거로운 구조임이 분명하지만, 매끄럽고 기하적으로 연결된
곡선의 고유 정칙 모형의 특수 올에서 정확히 이런 구조가 나타난다. 물론
지표의 순서를 바꾸는 것까지 동일시하여 이 형들을 다룬다.

\begin{definition}
\label{models-definition-type-equivalent}
두 수치형 $n, m_i, a_{ij}, w_i, g_i$와
$n', m'_i, a'_{ij}, w'_i, g'_i$에 대해 $\{1, \ldots, n\}$의 순열
$\sigma$가 존재하여 $m_i = m'_{\sigma(i)}$,
$a_{ij} = a'_{\sigma(i)\sigma(j)}$, $w_i = w'_{\sigma(i)}$,
$g_i = g'_{\sigma(i)}$이면 두 형이 {\it 동치}라고 한다.
\end{definition}

\noindent
수치형에는 종수가 있다.

\begin{lemma}
\label{models-lemma-genus}
$n, m_i, a_{ij}, w_i, g_i$를 수치형이라 하자. 그러면 식
$$
g = 1 + \sum m_i(w_i(g_i - 1) - \frac{1}{2} a_{ii})
$$
은 정수이다.
\end{lemma}

\begin{proof}
$g$가 정수임을 보이려면 $\sum a_{ii}m_i$가 짝수임을 보이면 된다.
이는 다음과 같이 법 $2$로 계산하면 알 수 있다.
\begin{align*}
\sum\nolimits_i a_{ii} m_i 
& \equiv
\sum\nolimits_{i,\ m_i\text{가 홀수}} a_{ii}m_i \\
& \equiv
\sum\nolimits_{i,\ m_i\text{가 홀수}} \sum\nolimits_{j \not = i} a_{ij}m_j \\
& \equiv
\sum\nolimits_{i,\ m_i\text{가 홀수}}
\sum\nolimits_{j \not = i,\ m_j\text{가 홀수}} a_{ij}m_j \\
& \equiv
\sum\nolimits_{i < j,\ m_i\text{와 }m_j\text{가 홀수}} a_{ij}(m_i + m_j) \\
& \equiv
0
\end{align*}
여기서는 $a_{ij} = a_{ji}$라는 사실과 모든 $i$에 대해
$\sum_j a_{ij}m_j = 0$이라는 사실을 사용하였다.
\end{proof}

\begin{definition}
\label{models-definition-genus}
$g = 1 + \sum m_i(w_i(g_i - 1) - \frac{1}{2} a_{ii})$가 보조정리
\ref{models-lemma-genus}의 정수이면 $n, m_i, a_{ij}, w_i, g_i$를
{\it 종수 $g$인 수치형}이라 한다.
\end{definition}

\noindent
아래 보조정리 \ref{models-lemma-genus-nonnegative}에서 종수가 거의 항상
$\geq 0$임을 증명할 것이다. 그러나 음의 종수를 갖는 수치형도 존재한다.

\begin{lemma}
\label{models-lemma-irreducible}
$n, m_i, a_{ij}, w_i, g_i$를 종수 $g$인 수치형이라 하자.
$n = 1$이면 $a_{11} = 0$이고 $g = 1 + m_1w_1(g_1 - 1)$이다.
또한 이런 모든 수치형은 다음과 같이 분류된다.
\begin{enumerate}
\item $g < 0$이면 $g_1 = 0$이고, $m_1w_1 = 1 - g$의 인수분해에
대응하는 $n = 1$, 종수 $g$인 수치형은 유한 개뿐이다.
\item $g = 0$이면 보조정리 \ref{models-lemma-genus-zero}에서처럼
$m_1 = 1$, $w_1 = 1$, $g_1 = 0$이다.
\item $g = 1$이면 $g_1 = 1$이지만 $m_1, w_1$은 임의의 양의 정수일 수
있다. 이는 보조정리 \ref{models-lemma-genus-one}의 경우 (\ref{models-item-one})이다.
\item $g > 1$이면 $g_1 > 1$이고, $m_1w_1(g_1 - 1) = g - 1$의
인수분해에 대응하는 $n = 1$, 종수 $g$인 수치형은 유한 개뿐이다.
\end{enumerate}
\end{lemma}

\begin{proof}
주장에서 바로 알 수 있다.
\end{proof}

\begin{lemma}
\label{models-lemma-diagonal-negative}
$n, m_i, a_{ij}, w_i, g_i$를 종수 $g$인 수치형이라 하자.
$n > 1$이면 모든 $i$에 대해 $a_{ii} < 0$이다.
\end{lemma}

\begin{proof}
행렬 $A$에 보조정리 \ref{models-lemma-recurring-symmetric-real}을 적용하면 된다.
\end{proof}

\begin{lemma}
\label{models-lemma-minus-one}
$n, m_i, a_{ij}, w_i, g_i$를 종수 $g$인 수치형이라 하고
$n > 1$이라 가정하자. 지표 $i$가 종수 $g$에 기여하는 항
$m_i(w_i(g_i - 1) - \frac{1}{2} a_{ii})$가 $< 0$이면
$g_i = 0$이고 $a_{ii} = -w_i$이다.
\end{lemma}

\begin{proof}
보조정리 \ref{models-lemma-diagonal-negative}과 $w_i > 0$, $g_i \geq 0$,
$w_i | a_{ii}$에서 즉시 따른다.
\end{proof}

\begin{definition}
\label{models-definition-type-minus-one}
$n, m_i, a_{ij}, w_i, g_i$를 수치형이라 하자. $g_i = 0$이고
$a_{ii} = -w_i$이면 $i$를 {\it $(-1)$-지표}라 한다.
\end{definition}

\noindent
$(-1)$-지표를 ``수축''할 수 있다.

\begin{lemma}
\label{models-lemma-contract}
$n, m_i, a_{ij}, w_i, g_i$를 수치형 $T$라 하자. $n$이
$(-1)$-지표라고 가정하자. 그러면 $n', m'_i, a'_{ij}, w'_i, g'_i$로
주어지는 수치형 $T'$가 존재하여 다음을 만족한다.
\begin{enumerate}
\item $n' = n - 1$,
\item $m'_i = m_i$,
\item $a'_{ij} = a_{ij} - a_{in}a_{jn}/a_{nn}$,
\item $a_{in}/w_n$이 짝수이고 $a_{in}/w_i$가 홀수이면
$w'_i = w_i/2$이고, 그렇지 않으면 $w'_i = w_i$이다.
\item $g'_i =
\frac{w_i}{w'_i}(g_i - 1) + 1 + \frac{a_{in}^2 - w_na_{in}}{2w'_iw_n}$.
\end{enumerate}
더욱이 $g = g'$이다.
\end{lemma}

\begin{proof}
예를 들어 보조정리 \ref{models-lemma-irreducible}에 의해 $n > 1$이고 따라서
$n' \geq 1$이다. $n', m'_i, a'_{ij}, w'_i, g'_i$에 대해 정의
\ref{models-definition-type}의 조건 (1)--(5)를 확인하자.

\medskip\noindent
조건 (1)은 자명하다.

\medskip\noindent
조건 (2). $A' = (a'_{ij})$의 대칭성은 자명하다. 보조정리
\ref{models-lemma-diagonal-negative}에 의해 $a_{nn} < 0$이므로
$i \not = j$이면 $a'_{ij} \geq a_{ij} \geq 0$이다.

\medskip\noindent
조건 (3). $i \in I$이고 $i' \in \{1, \ldots, n - 1\} \setminus I$이면
$a'_{ii'} = 0$이 되는 $I \subset \{1, \ldots, n - 1\}$가 있다고 하자.
그러면 각 $i \in I$, $i' \in I'$에 대해 $a_{in}a_{i'n} = 0$이다.
따라서 모든 $i \in I$에 대해 $a_{in} = 0$이어서
$I \subset \{1, \ldots, n\}$가 $T$의 조건 (3)에 모순되거나, 모든
$i' \in \{1, \ldots, n - 1\} \setminus I$에 대해 $a_{i'n} = 0$이어서
$I \cup \{n\} \subset \{1, \ldots, n\}$가 $T$의 조건 (3)에
모순된다. 따라서 $T'$에 대해 (3)이 성립한다.

\medskip\noindent
조건 (4). 계산하면
$$
\sum\nolimits_{j = 1}^{n - 1} a'_{ij}m_j =
\sum\nolimits_{j = 1}^{n - 1}
(a_{ij}m_j - \frac{a_{in}a_{jn}m_j}{a_{nn}}) =
- a_{in}m_n - \frac{a_{in}}{a_{nn}}(-a_{nn}m_n) = 0
$$
이며 원하는 결과이다.

\medskip\noindent
조건 (5). $w'_i$가 $a_{in}a_{jn}/a_{nn}$을 나눔을 보이면 된다.
$a_{nn} = -w_n$, $w_n | a_{jn}$, $w_i | a_{in}$이므로 이는 분명하다.

\medskip\noindent
$g = g'$임을 보이기 위해 먼저
\begin{align*}
g
& =
1 + \sum\nolimits_{i = 1}^n m_i(w_i(g_i - 1) - \frac{1}{2}a_{ii}) \\
& =
1 + \sum\nolimits_{i = 1}^{n - 1} m_i(w_i(g_i - 1) - \frac{1}{2}a_{ii})
-\frac{1}{2}m_nw_n \\
& =
1 +  \sum\nolimits_{i = 1}^{n - 1}
m_i(w_i(g_i - 1) - \frac{1}{2}a_{ii} - \frac{1}{2}a_{in})
\end{align*}
$g'$의 식과 비교하면
$$
w'_i(g'_i - 1) - \frac{1}{2}a'_{ii} =
w_i(g_i - 1) - \frac{1}{2}a_{in} - \frac{1}{2}a_{ii}
$$
가 $i \leq n - 1$에 대해 성립하면 충분하다. 다시 말해
$$
g'_i = \frac{2w_i(g_i - 1) - a_{in} - a_{ii} + a'_{ii} + 2w'_i}{2w'_i} =
\frac{w_i}{w'_i}(g_i - 1) + 1 + \frac{a_{in}^2 - w_na_{in}}{2w'_iw_n}
$$
(4)처럼 $w'_i$를 택하면 이것이 정수 $\geq 0$임은 초등적으로 확인된다.
\end{proof}

\begin{lemma}
\label{models-lemma-top-genus}
$n, m_i, a_{ij}, w_i, g_i$를 수치형이라 하자. $e$를 $i < j$이고
$a_{ij} > 0$인 순서쌍 $(i, j)$의 수라 하자. 그러면
$g_{top} = 1 - n + e$는 $\geq 0$이다.
\end{lemma}

\begin{proof}
그렇지 않으면 $e < n - 1$이고, 이는 모든 $j \not = i$에 대해
$a_{ij} = 0$인 $i$가 존재한다는 뜻이다. 이는 정의
\ref{models-definition-type}의 가정 (3)에 모순된다.
\end{proof}

\begin{definition}
\label{models-definition-top-genus}
$n, m_i, a_{ij}, w_i, g_i$를 수치형 $T$라 하자. $T$의
{\it 위상적 종수}는 보조정리 \ref{models-lemma-top-genus}의 음이 아닌 정수
$g_{top} = 1 - n + e$이다.
\end{definition}

\noindent
종수를 위상적 종수로 한정하고자 한다. 그러나 보조정리
\ref{models-lemma-irreducible}의 $n = 1$인 수치형처럼 이것이 항상 가능한 것은
아니다. 다음과 같이 정의되는 극소 수치형에 대해서는 가능하다.

\begin{definition}
\label{models-definition-type-minimal}
종수 $g$인 수치형 $n, m_i, a_{ij}, w_i, g_i$에 $g_i = 0$이고
$a_{ii} = -w_i$인 $i$가 존재하지 않으면, 즉 $(-1)$-지표가 존재하지
않으면 이 수치형을 {\it 극소}라 한다.
\end{definition}

\noindent
$n > 1$인 극소형의 종수 $g$가 $\max(1, g_{top})$ 이상임을 증명할 것이다.

\begin{lemma}
\label{models-lemma-non-irreducible-minimal-type-genus-at-least-one}
$n, m_i, a_{ij}, w_i, g_i$가 $n > 1$인 극소 수치형이면 $g \geq 1$이다.
\end{lemma}

\begin{proof}
보조정리 \ref{models-lemma-minus-one}과 극소형의 정의에 의해
$\Phi_i = m_i(w_i(g_i - 1) - \frac{1}{2} a_{ii})$가 음이 아니므로
$g = 1 + \sum \Phi_i$이다.
\end{proof}

\begin{lemma}
\label{models-lemma-genus-nonnegative}
$n, m_i, a_{ij}, w_i, g_i$가 $n > 1$인 극소 수치형이면
$g \geq g_{top}$이다.
\end{lemma}

\begin{proof}
고유 정칙 모형에 결부된 수치형의 경우에만 관심 있는 독자는 이 증명을
건너뛰어도 된다. 뒤에서 기하적 상황에 대해 다시 증명할 것이다. 다음과
같이 쓸 수 있다.
$$
g_{top} = 1 - n + \frac{1}{2}\sum\nolimits_{a_{ij} > 0} 1 =
1 + \sum\nolimits_i (-1 +
\frac{1}{2}\sum\nolimits_{j \not = i,\ a_{ij} > 0} 1) 
$$
한편
\begin{align*}
g & =
1 + \sum m_i(w_i(g_i - 1) - \frac{1}{2} a_{ii}) \\
& =
1 + \sum m_iw_ig_i - \sum m_iw_i +
\frac{1}{2} \sum\nolimits_{i \not = j} a_{ij}m_j \\
& =
1 + \sum\nolimits_i
m_iw_i(-1 + g_i + \frac{1}{2} \sum\nolimits_{j \not = i} \frac{a_{ij}}{w_i})
\end{align*}
첫째 등식은 정의이고, 둘째 등식에는 $\sum a_{ij}m_j = 0$을 사용했으며,
마지막 등식에는 $a_{ij} = a_{ji}$와 합의 순서 교환을 사용했다.
$g_{top}$의 식과 비교하면
$$
\Psi_i =
m_iw_i(-1 + g_i + \frac{1}{2} \sum\nolimits_{j \not = i} \frac{a_{ij}}{w_i})
- (-1 + \frac{1}{2}\sum\nolimits_{j \not = i,\ a_{ij} > 0} 1)
$$
가 각 $i$에 대해 $\geq 0$이면 보조정리가 성립한다. 그러나 이것이 항상
성립하지는 않는다. 어떤 지표에 대해 $\Psi_i < 0$이 가능한지 분석하자.
먼저
$$
(-1 + g_i + \frac{1}{2}\sum\nolimits_{j \not = i} \frac{a_{ij}}{w_i}) \geq
(-1 + \frac{1}{2}\sum\nolimits_{j \not = i,\ a_{ij} > 0} 1)
$$
이다. 이는 $a_{ij}/w_i$가 음이 아닌 정수이기 때문이다. $m_iw_i$가
양의 정수이므로, $m_iw_i = 1$이거나 위 부등식의 왼쪽이 $\geq 0$이면
$\Psi_i \geq 0$이다. 후자는 $g_i > 0$이거나 적어도 두 지표 $j$에 대해
$a_{ij} > 0$이거나 $a_{ij} > w_i$인 $j$가 있을 때 성립한다. 따라서
$$
P = \{i : \Psi_i < 0\}
$$
는 $m_iw_i > 1$, $g_i = 0$이고, 유일한 $j$에 대해 $a_{ij} > 0$이며
그 $j$에 대해 $a_{ij} = w_i$인 지표 $i$들의 집합이다. 더욱이
$$
i \in P \Rightarrow \Psi_i = \frac{1}{2}(-m_iw_i + 1)
$$
증명의 전략은 $i \in P$가 주어졌을 때 $i$의 이웃, 즉 $a_{ij} > 0$인
$j$의 $\Psi_j$에서 조금을 빌릴 수 있음을 보이는 것이다. 다만
$\Psi_j = 0$인 지표 $j$일 수도 있어 이 전략이 그대로 작동하지는 않는다.

\medskip\noindent
집합
$$
Z = \{j : g_j = 0\text{이고 }
j\text{는 정확히 두 이웃 }i, k\text{를 가지며 }
a_{ij} = w_j = a_{jk}\}
$$
을 생각하자. $j \in Z$에 대해 $\Psi_j = 0$이다. $s \geq 0$,
$i \in P$, $j_1, \ldots, j_s \in Z$이고
$a_{ij_1} > 0, a_{j_1j_2} > 0, \ldots, a_{j_{s - 1}j_s} > 0$인 수열
$M = (i, j_1, \ldots, j_s)$을 생각한다. 수치형이 $P$에 속한 두 지표로만
이루어지거나, 더 일반적으로 $P$에 속한 두 지표와 $Z$에 속한 나머지 모든
지표로 이루어지면 $g_{top} = 0$이고 보조정리
\ref{models-lemma-non-irreducible-minimal-type-genus-at-least-one}로 결론을 얻는다.
따라서 이런 경우는 제외해도 되고 실제로 제외한다.

\medskip\noindent
극대 수열 $M = (i, j_1, \ldots, j_s)$을 택하고 $k$를 $j_s$의 둘째
이웃이라 하자. ($s = 0$이면 $k$는 $i$의 유일한 이웃이다.) 극대성에
의해 $k \not \in Z$이고, 방금 논의에 의해 $k \not \in P$이다. 또한
$w_i = a_{ij_1} = w_{j_1} = a_{j_1j_2} = \ldots = w_{j_s} =
a_{j_sk}$이다. 수치형의 정의를 보면
\begin{align*}
m_ia_{ii} + m_{j_1}w_i & = 0,\\
m_iw_i + m_{j_1}a_{j_1j_1} + m_{j_2}w_i & = 0,\\
\ldots & \ldots \\
m_{j_{s - 1}}w_i + m_{j_s}a_{j_sj_s} + m_kw_i & = 0
\end{align*}
첫째 등식과 수치형의 극소성으로부터 $m_{j_1} \geq 2m_i$를 얻는다.
이어서 둘째 등식으로부터 $m_{j_2} \geq 3m_i$를 얻으며 같은 방식으로
계속한다. 어느 경우든 ($s = 0$인 경우도 포함하여) $m_k \geq 2m_i$이다.

\medskip\noindent
$t > 0$과 위와 같은 서로 다른 극대 수열 $M_1, \ldots, M_t$가 있어,
$M_b = (i_b, j_{b, 1}, \ldots, j_{b, s_b})$이고 각
$b = 1, \ldots, t$에 대해 $k$가 $j_{b, s_b}$의 이웃이 되는 지표 $k$를
택하자. $\Phi_j + \sum_{b = 1, \ldots, t} \Phi_{i_b} \geq 0$임을 보이겠다.
그러면 앞의 논의로 보조정리의 증명이 끝난다. $M$을
$b = 1, \ldots, t$에 대해 $M_b$에 나타나는 지표들의 합집합이라 하자.
다음과 같이 쓴다.
$$
\Psi_k =
-\sum\nolimits_{b = 1, \ldots, t} \Psi_{i_b} + \Psi_k'
$$
여기서
\begin{align*}
\Psi_k' & =
m_kw_k\left(-1 + g_k +
\frac{1}{2} \sum\nolimits_{b = 1, \ldots t}
(\frac{a_{kj_{b, s_b}}}{w_k} - \frac{m_{i_b}w_{i_b}}{m_kw_k}) +
\frac{1}{2} \sum\nolimits_{l \not = k,\ l \not \in M}
\frac{a_{kl}}{w_k}
\right) \\
&
-\left(
-1 + \frac{1}{2}\sum\nolimits_{l \not = k,\ l \not \in M,\ a_{kl} > 0} 1
\right)
\end{align*}
모순을 얻기 위해 $\Psi_k' < 0$이라 가정하자. 집합
$\{l : l \not = k,\ l \not \in M,\ a_{kl} > 0\}$이 비어 있으면
$\{1, \ldots, n\} = M \cup \{k\}$이고, 이 경우 $e = n - 1$이므로
$g_{top} = 0$이다. 따라서 보조정리
\ref{models-lemma-non-irreducible-minimal-type-genus-at-least-one}로 결과가
따른다. 그러므로 이런 $l$이 적어도 하나 있으며, 첫째 괄호 안의 합에
$(1/2)a_{kl}/w_k \geq 1/2$를 기여한다고 가정해도 된다. 각
$b = 1, \ldots, t$에 대해
$$
\frac{a_{kj_{b, s_b}}}{w_k} - \frac{m_{i_b}w_{i_b}}{m_kw_k} =
\frac{w_{i_b}}{w_k}(1 - \frac{m_{i_b}}{m_k})
$$
이다. 앞 문단의 $m_k \geq 2m_{i_b}$ 때문에 이 식은
$\geq \frac{1}{2}$이고, $w_k < w_{i_b}$이면 $\geq 1$이다. 따라서
$\Psi_k' < 0$이면 $g_k = 0$이다. $t \geq 2$이거나 $t = 1$이고
$w_k < w_{i_1}$이면, 위에서 보인 $l$의 존재를 사용하여
$\Psi_k' \geq 0$을 얻는데 이 역시 모순이다. 따라서 $t = 1$이고
$w_k = w_{i_1}$이다. $l$에 대한 합에 영이 아닌 항이 적어도 두 개 있거나
그런 $k$가 하나이고 $a_{kl} > w_k$이면 역시 $\Psi_k' \geq 0$이다.
남은 가능성은 $t = 1$이고 $a_{kl} = w_k$인 $l$이 하나뿐인 경우이다.
그러나 이는 $k \in Z$를 뜻하여 $M_1$의 극대성과 모순되므로 허용되지 않는다.
\end{proof}

\begin{lemma}
\label{models-lemma-minus-two}
$n, m_i, a_{ij}, w_i, g_i$를 종수 $g$인 수치형이라 하고
$n > 1$이라 가정하자. 지표 $i$가 종수 $g$에 기여하는 항
$m_i(w_i(g_i - 1) - \frac{1}{2} a_{ii})$가 $0$이면
$g_i = 0$이고 $a_{ii} = -2w_i$이다.
\end{lemma}

\begin{proof}
보조정리 \ref{models-lemma-diagonal-negative}과 $w_i > 0$, $g_i \geq 0$,
$w_i | a_{ii}$에서 즉시 따른다.
\end{proof}

\noindent
이 관계를 만족하는 지표는 극소 수치형의 구조에서 중요한 역할을 한다.
따라서 이들에게 이름을 붙인다.

\begin{definition}
\label{models-definition-type-minus-two}
$n, m_i, a_{ij}, w_i, g_i$를 종수 $g$인 수치형이라 하자.
$g_i = 0$이고 $a_{ii} = -2w_i$이면 $i$를 {\it $(-2)$-지표}라 한다.
\end{definition}

\noindent
종수 $g$인 극소 수치형이 주어졌을 때 $(-2)$-지표는 정확히 식
$$
g = 1 + \sum m_i(w_i(g_i - 1) - \frac{1}{2} a_{ii})
$$
에서 종수에 양수를 기여하지 않는 지표들이다. 따라서 뒤에서 보듯이
$(-2)$-지표에 결부된 양들을 한정하는 일은 다소 까다롭다.

\begin{remark}
\label{models-remark-genus-equality}
$n, m_i, a_{ij}, w_i, g_i$를 $n > 1$인 극소 수치형이라 하자.
보조정리 \ref{models-lemma-genus-nonnegative}에서 등식 $g = g_{top}$이 성립할
수 있다. 예를 들어 모든 $i$에 대해 $m_i = w_i = 1$, $g_i = 0$이고
$i < j$에 대해 $a_{ij} \in \{0, 1\}$인 경우가 그렇다.
\end{remark}


\section{수치형의 Picard 군}
\label{models-section-picard-group}

\noindent
정의는 다음과 같다.

\begin{definition}
\label{models-definition-picard-group}
$n, m_i, a_{ij}, w_i, g_i$를 수치형 $T$라 하자. $T$의
{\it Picard 군}은 행렬 $(a_{ij}/w_i)$의 여핵, 더 정확히는
$$
\Pic(T) =
\Coker\left(
\mathbf{Z}^{\oplus n} \to \mathbf{Z}^{\oplus n},\quad
e_i
\mapsto
\sum \frac{a_{ij}}{w_j}e_j
\right)
$$
이다. 여기서 $e_i$는 $\mathbf{Z}^{\oplus n}$의 $i$번째 표준기저 벡터이다.
\end{definition}

\begin{lemma}
\label{models-lemma-picard-rank-1}
$n, m_i, a_{ij}, w_i, g_i$를 수치형 $T$라 하자. $T$의 Picard 군은
계수 $1$인 유한 생성 아벨군이다.
\end{lemma}

\begin{proof}
$n = 1$이면 $A = (a_{ij})$는 영행렬이므로 결과는 분명하다.
$n > 1$이면 보조정리 \ref{models-lemma-recurring-real} 또는 보조정리
\ref{models-lemma-recurring-symmetric-real}에 의해 행렬 $A$의 계수는 $n - 1$이다.
행들을 $1/w_i$배 하여도 계수는 변하지 않는다. 이것으로 증명된다.
\end{proof}

\begin{lemma}
\label{models-lemma-picard-T-and-A}
$n, m_i, a_{ij}, w_i, g_i$를 수치형 $T$라 하자. $A = (a_{ij})$라 하면
$\Pic(T) \subset \Coker(A)$이다.
\end{lemma}

\begin{proof}
$\Pic(T)$가 $(a_{ij}/w_i)$의 여핵이므로 가환 그림
$$
\xymatrix{
0 \ar[r] &
\mathbf{Z}^{\oplus n} \ar[rr]_A & &
\mathbf{Z}^{\oplus n} \ar[rr] & &
\Coker(A) \ar[r] & 0 \\
0 \ar[r] &
\mathbf{Z}^{\oplus n} \ar[rr]^{(a_{ij}/w_i)} \ar[u]_{\text{id}} & &
\mathbf{Z}^{\oplus n} \ar[rr] \ar[u]_{\text{diag}(w_1, \ldots, w_n)} & &
\Pic(T) \ar[r] \ar[u] & 0
}
$$
이 있고 각 행은 완전하다. 뱀 보조정리에 의해
$\Pic(T) \subset \Coker(A)$이다.
\end{proof}

\begin{lemma}
\label{models-lemma-contract-picard-group}
$n, m_i, a_{ij}, w_i, g_i$를 수치형 $T$라 하고 $n$이 $(-1)$-지표라고
가정하자. $T'$을 보조정리 \ref{models-lemma-contract}에서 구성한 수치형이라
하자. 단사 사상
$$
\Pic(T) \to \Pic(T')
$$
이 존재하고, 그 여핵은 초등 아벨 $2$-군이다.
\end{lemma}

\begin{proof}
$n' = n - 1$임을 상기하자. $e_i$와 $e'_i$를 각각
$\mathbf{Z}^{\oplus n}$과 $\mathbf{Z}^{\oplus n - 1}$의 $i$번째 기저
벡터라 하자. 먼저
$$
q : \mathbf{Z}^{\oplus n} \to \mathbf{Z}^{\oplus n - 1},
\quad e_n \mapsto 0\text{이고 }e_i \mapsto e'_i\text{ (}i \leq n - 1\text{)}
$$
라 쓰고
$$
p : \mathbf{Z}^{\oplus n} \to \mathbf{Z}^{\oplus n - 1},\quad
e_n \mapsto \sum\nolimits_{j = 1}^{n - 1} \frac{a_{nj}}{w'_j} e'_j
\text{이고 }
e_i \mapsto \frac{w_i}{w'_i} e'_i\text{ (}i \leq n - 1\text{)}
$$
라 둔다. 생략하는 계산으로부터 가환 그림
$$
\xymatrix{
\mathbf{Z}^{\oplus n} \ar[rr]_{(a_{ij}/w_i)} \ar[d]_q & &
\mathbf{Z}^{\oplus n} \ar[d]^p \\
\mathbf{Z}^{\oplus n'} \ar[rr]^{(a'_{ij}/w'_i)} & &
\mathbf{Z}^{\oplus n'}
}
$$
이 있음을 알 수 있다. 위 화살표의 여핵이 $\Pic(T)$이고 아래 화살표의
여핵이 $\Pic(T')$이므로 원하는 Picard 군 준동형을 얻는다.
$\frac{w_i}{w'_i} \in \{1, 2\}$이므로 $\Pic(T) \to \Pic(T')$의 여핵은
$2$에 의해 소멸한다. 실제로 모든 $i \leq n - 1$에 대해 $2e'_i$는
$p$의 상에 속한다. 마지막으로 $\Pic(T) \to \Pic(T')$가 단사임을 보이자.
$L = (l_1, \ldots, l_n)$을 $\Pic(T')$에서 영으로 가는 $\Pic(T)$ 원소의
대표라 하자. $q$가 전사이므로 그림 추적으로 $L$이 $p$의 핵에 속한다고
가정할 수 있다. 이는 $l_na_{ni}/w'_i + l_iw_i/w'_i = 0$, 즉
$l_i = - a_{ni}/w_i l_n$이라는 뜻이다. 따라서 $L$은 사상
$(a_{ij}/w_j)$ 아래에서 $-l_ne_n$의 상이고 보조정리가 증명된다.
\end{proof}

\begin{lemma}
\label{models-lemma-picard-group-genus-nonpositive}
$n, m_i, a_{ij}, w_i, g_i$를 수치형 $T$라 하자. $T$의 종수 $g$가
$\leq 0$이면 $\Pic(T) = \mathbf{Z}$이다.
\end{lemma}

\begin{proof}
$n$에 대한 귀납법을 쓴다. $n = 1$이면 주장은 분명하다. $n > 1$이면
보조정리 \ref{models-lemma-non-irreducible-minimal-type-genus-at-least-one}에
의해 $T$는 극소가 아니다. $T$를 동치인 형으로 바꾸어 $n$이
$(-1)$-지표라고 가정할 수 있다. 보조정리
\ref{models-lemma-contract-picard-group}에 의해 $\Pic(T) \subset \Pic(T')$이다.
보조정리 \ref{models-lemma-contract}에 의해 $T'$의 종수는 $T$의 종수와 같으므로
귀납가정으로 결론을 얻는다.
\end{proof}








\section{진부분그래프의 분류}
\label{models-section-classify-proper-subgraphs}

\noindent
이 절에서는 종수 $g$인 수치형 $n, m_i, a_{ij}, w_i, g_i$가 주어졌다고
가정한다. $(-2)$-지표(정의 \ref{models-definition-type-minus-two})들로만 이루어진
가능한 ``부분그래프''의 완전한 목록을 찾는 동시에 종수 $1$인 모든 극소
수치형을 분류한다. 다시 말해 이 절에서는 명제
\ref{models-proposition-classify-subgraphs}와 보조정리 \ref{models-lemma-genus-one}을
증명한다.

\medskip\noindent
전략은 다음과 같다. 종수 $g$인 수치형 $n, m_i, a_{ij}, w_i, g_i$를
택하자. $I \subset \{1, \ldots, n\}$를 $(-2)$-지표들로 이루어진
부분집합으로서, $j \in J$, $j' \in I \setminus J$일 때
$a_{jj'} = 0$이 되는 공집합이 아닌 진부분집합 $J \subset I$가 없다고 하자.
$I$의 크기 $|I|$에 대해 귀납한다. $I = \{i\}$가 $1$개의 지표로 이루어진
경우 $m_i$, $a_{ii}$, $w_i$에 대한 제약은 정의
\ref{models-definition-type}의 $w_i | a_{ii}$와 보조정리
\ref{models-lemma-diagonal-negative}의 $a_{ii} < 0$뿐이며, 이를 시작 경우로
삼는다. 귀납 단계에서는 먼저 $|I'| < |I|$인 부분집합 $I' \subset I$에
귀납가정을 적용한다. 이로부터 $i, j \in I$에 대한 가능한
$m_i, a_{ij}, w_i$에 제약이 생긴다. 특히 $|I'| < |I| \leq n$이므로
$\sum a_{ij}m_j = 0$과 보조정리 \ref{models-lemma-recurring-symmetric-real}에서
부분행렬 $(a_{ij})_{i, j \in I'}$가 음의 정부호이고 그 행렬식의 부호가
$(-1)^m$임이 따른다. 남은 각 가능성에 대해 $(a_{ij})_{i, j \in I}$의
행렬식을 계산한다. 그 부호가 $-(-1)^{|I|}$이면 Sylvester 정리에 의해
$(a_{ij})_{i, j \in I}$가 반음의 정부호가 아니므로 이 경우를 버릴 수
있다. 행렬식의 부호가 $(-1)^{|I|}$이면 $|I| < n$이고, 이 경우가 가능한
진부분그래프로 나타날 수 있다고 잠정적으로 결론내려 이 절의 보조정리 중
하나에 기록한다. 행렬식이 $0$이면 보조정리
\ref{models-lemma-recurring-symmetric-real}에 의해 $|I| = n$이어야 하고
$g = 0$이다. 이 경우에는 $i, j \in I$에 대한 가능한 모든
$m_i, a_{ij}, w_i$를 실제로 찾아 보조정리 \ref{models-lemma-genus-one}에
열거한다. 논증을 마치면 $n > 1$이고 종수 $1$인 모든 극소 수치형을
얻는다. 종수 공식과 앞 절의 논의에 의해 각각이 반드시 $(-2)$-지표로만
이루어져 귀납 과정에 나타나기 때문이다. 마지막으로 보조정리
\ref{models-lemma-genus-one}의 가능성 목록을 살펴보면 이 절에서 가능하다고
열거한 진부분그래프의 모든 배치가 실제로 종수 $1$인 수치형에서 이미
나타남을 확인할 수 있다.

\medskip\noindent
$i$와 $j$가 $a_{ij} > 0$인 $(-2)$-지표라고 하자. 보조정리
\ref{models-lemma-recurring-symmetric-real}에 의해 행렬 $A = (a_{ij})$가
반음의 정부호이므로 행렬
$$
\left(
\begin{matrix}
-2w_i & a_{ij} \\
a_{ij} & -2w_j
\end{matrix}
\right)
$$
는 $n = 2$가 아닌 한 음의 정부호이다. $n = 2$인 경우도 가능하며, 이때
행렬식 $4w_1w_2 - a_{12}^2$는 영이다. $\text{lcm}(w_1, w_2)$이
$a_{12}$를 나눈다는 사실을 사용하면 가능한 경우가
$$
(w_1, w_2, a_{12}) = (w, w, 2w), (w, 4w, 4w), \text{ 또는 }(4w, w, 4w)
$$
뿐임을 쉽게 알 수 있다. $(4w, w, 4w)$인 경우는 $(w, 4w, 4w)$인
경우에서 지표 $i, j$를 서로 바꾸어 얻는다. 이 경우들에는 $g = 1$이고,
보조정리 \ref{models-lemma-genus-one}의 경우 (\ref{models-item-two-cycle})와
(\ref{models-item-up4})를 얻는다. $n > 2$라 가정하면 표시된 행렬의 행렬식
$4w_iw_j - a_{ij}^2$가 $> 0$이므로 $a_{ij}^2/w_iw_j < 4$이다. 한편
$\text{lcm}(w_i, w_j) | a_{ij}$임을 알고 있으므로
$a_{ij}^2/w_iw_j$는 정수이다. 따라서 $a_{ij}^2/w_iw_j \in \{1, 2, 3\}$이고
$w_i | w_j$이거나 그 반대이다. 이에 따라 가능한 경우는
$$
(w_1, w_2, a_{12}) = (w, w, w), (w, 2w, 2w), (w, 3w, 3w),
(2w, w, 2w), \text{ 또는 }(3w, w, 3w)
$$
이다. $(2w, w, 2w)$인 경우는 $(w, 2w, 2w)$인 경우에서 지표 $i,j$를
바꾸어 얻으며, $(3w, w, 3w)$와 $(w, 3w, 3w)$도 마찬가지이다.
처음 세 해는 보조정리 \ref{models-lemma-two-by-two}의 경우 (\ref{models-item-A2}),
(\ref{models-item-B2}), (\ref{models-item-G2})를 준다. 이 보조정리에서는 각 $k$에
대한 $\sum_l a_{kl}m_l = 0$이 특히 $k = i$에 대해
$a_{ii}m_i + a_{ij}m_j \leq 0$, $k = j$에 대해
$a_{ij}m_i + a_{jj}m_j \leq 0$을 함의한다는 사실로부터 정수
$m_i$와 $m_j$에 대한 결과도 명시하였다.

\begin{lemma}
\label{models-lemma-two-by-two}
다음 형태의 진부분그래프를 분류한다.
$$
\xymatrix{
\bullet \ar@{-}[r] & \bullet
}
$$
$n > 2$이고 $a_{ij} > 0$인 $(-2)$-지표의 쌍 $i,j$가 주어지면,
순서를 바꾸는 것까지 동일시하여 $m$들, $a$들, $w$들은 다음과 같다.
\begin{enumerate}
\item
\label{models-item-A2}
다음으로 주어지고,
$$
\left(
\begin{matrix}
m_1 \\
m_2
\end{matrix}
\right),
\quad
\left(
\begin{matrix}
-2w & w \\
w & -2w
\end{matrix}
\right),
\quad
\left(
\begin{matrix}
w \\
w
\end{matrix}
\right)
$$
$w$는 임의이며 $2m_1 \geq m_2$이고 $2m_2 \geq m_1$이다. 또는
\item
\label{models-item-B2}
다음으로 주어지고,
$$
\left(
\begin{matrix}
m_1 \\
m_2
\end{matrix}
\right),
\quad
\left(
\begin{matrix}
-2w & 2w \\
2w & -4w
\end{matrix}
\right),
\quad
\left(
\begin{matrix}
w \\
2w
\end{matrix}
\right)
$$
$w$는 임의이며 $m_1 \geq m_2$이고 $2m_2 \geq m_1$이다. 또는
\item
\label{models-item-G2}
다음으로 주어지고,
$$
\left(
\begin{matrix}
m_1 \\
m_2
\end{matrix}
\right),
\quad
\left(
\begin{matrix}
-2w & 3w \\
3w & -6w
\end{matrix}
\right),
\quad
\left(
\begin{matrix}
w \\
3w
\end{matrix}
\right)
$$
$w$는 임의이며 $2m_1 \geq 3m_2$이고 $2m_2 \geq m_1$이다.
\end{enumerate}
\end{lemma}

\begin{proof}
위의 논의를 참조하라.
\end{proof}

\noindent
$i$, $j$, $k$가 $a_{ij} > 0$, $a_{jk} > 0$인 세 $(-2)$-지표라고 하자.
다시 말해 지표 $i$는 $j$와 ``만나고'' $j$는 $k$와 ``만난다''.
각 쌍 $(i, j)$, $(i, k)$, $(j, k)$가 보조정리 \ref{models-lemma-two-by-two}의
목록 중 하나라는 사실을 이후 별도 언급 없이 사용한다. 보조정리
\ref{models-lemma-recurring-symmetric-real}에 의해 행렬 $A = (a_{ij})$가
반음의 정부호이므로 행렬
$$
\left(
\begin{matrix}
-2w_i & a_{ij} & a_{ik} \\
a_{ij} & -2w_j & a_{jk} \\
a_{ik} & a_{jk} & -2w_k
\end{matrix}
\right)
$$
은 $n = 3$이 아닌 한 음의 정부호이다. $n = 3$인 경우도 가능하다.
이때 행렬식\footnote{그 값은
$-8w_iw_jw_k + 2a_{ij}^2w_k + 2a_{jk}^2w_i + 2a_{ik}^2w_j +
2a_{ij}a_{jk}a_{ik}$이다.}은 영이고 방정식
$$
4 = \frac{a_{ij}^2}{w_iw_j} +
\frac{a_{jk}^2}{w_jw_k} +
\frac{a_{ik}^2}{w_iw_k} +
\frac{a_{ij}a_{ik}a_{jk}}{w_iw_jw_k}
$$
을 얻는다. 이 방정식 오른쪽의 마지막 항은 다음 등식 때문에 나머지
항들에 의해 결정된다.
$$
\left(\frac{a_{ij}a_{ik}a_{jk}}{w_iw_jw_k}\right)^2 =
\frac{a_{ij}^2}{w_iw_j} \frac{a_{jk}^2}{w_jw_k} \frac{a_{ik}^2}{w_iw_k}
$$
위에서
$\frac{a_{ij}^2}{w_iw_j}, \frac{a_{jk}^2}{w_jw_k}$는
$\{1, 2, 3\}$에 속하고 $\frac{a_{ik}^2}{w_iw_k}$는
$\{0, 1, 2, 3\}$에 속함을 보았으므로 가능한 경우는
$$
(\frac{a_{ij}^2}{w_iw_j}, \frac{a_{jk}^2}{w_jw_k}, \frac{a_{ik}^2}{w_iw_k}) =
(1, 1, 1), (1, 3, 0), (2, 2, 0),\text{ 또는 } (3, 1, 0)
$$
뿐이다. $(3, 1, 0)$인 경우는 $(1, 3, 0)$인 경우에서 지표 $i,j,k$의
순서를 뒤집어 얻는다. 각 경우에 $g = 1$이며, 보조정리
\ref{models-lemma-genus-one}의 경우 (\ref{models-item-three-cycle}),
(\ref{models-item-equal-up3}), (\ref{models-item-equal-down3}), (\ref{models-item-up-up}),
(\ref{models-item-up-down}), (\ref{models-item-down-up})에서 찾을 수 있다. 이 가운데
하나는 $(1, 1, 1)$에, 둘은 $(1, 3, 0)$에, 셋은 $(2, 2, 0)$에 대응한다.
$n > 3$이라 가정하면 부등식
$$
4 > \frac{a_{ij}^2}{w_iw_j} + \frac{a_{ik}^2}{w_iw_k} +
\frac{a_{jk}^2}{w_jw_k} + \frac{a_{ij}a_{ik}a_{jk}}{w_iw_jw_k}
$$
을 얻는다. 위에서 얻은 수들에 대한 제약을 사용하면 가능한 경우는
$$
(\frac{a_{ij}^2}{w_iw_j}, \frac{a_{jk}^2}{w_jw_k}, \frac{a_{ik}^2}{w_iw_k}) =
(1, 1, 0), (1, 2, 0),\text{ 또는 }(2, 1, 0)
$$
뿐이고, 특히 $a_{ik} = 0$이다. 여기서 $a_{ij} > 0$ 및 $a_{jk} > 0$을
가정했음을 상기하라. $(2, 1, 0)$인 경우는 $(1, 2, 0)$인 경우에서 지표
$i,j,k$의 순서를 뒤집어 얻는다. 처음 두 해는 보조정리
\ref{models-lemma-three-by-three}의 경우 (\ref{models-item-A3}), (\ref{models-item-C3}),
(\ref{models-item-B3})를 준다. 이 보조정리에서는 정수 $m_i$, $m_j$, $m_k$에 대한
결과도 명시하였다.

\begin{lemma}
\label{models-lemma-three-by-three}
다음 형태의 진부분그래프를 분류한다.
$$
\xymatrix{
\bullet \ar@{-}[r] &
\bullet \ar@{-}[r] &
\bullet
}
$$
$n > 3$이고 $(-2)$-지표의 삼중항 $i,j,k$가 주어져
$a_{ij}, a_{ik}, a_{jk}$ 중 적어도 둘이 영이 아니면, 순서를 바꾸는 것까지
동일시하여 $m$들, $a$들, $w$들은 다음과 같다.
\begin{enumerate}
\item
\label{models-item-A3}
다음으로 주어지고,
$$
\left(
\begin{matrix}
m_1 \\
m_2 \\
m_3
\end{matrix}
\right),
\quad
\left(
\begin{matrix}
-2w & w & 0 \\
w & -2w & w \\
0 & w & -2w
\end{matrix}
\right),
\quad
\left(
\begin{matrix}
w \\
w \\
w
\end{matrix}
\right)
$$
$2m_1 \geq m_2$, $2m_2 \geq m_1 + m_3$, $2m_3 \geq m_2$이다. 또는
\item
\label{models-item-C3}
다음으로 주어지고,
$$
\left(
\begin{matrix}
m_1 \\
m_2 \\
m_3
\end{matrix}
\right),
\quad
\left(
\begin{matrix}
-2w & w & 0 \\
w & -2w & 2w \\
0 & 2w & -4w
\end{matrix}
\right),
\quad
\left(
\begin{matrix}
w \\
w \\
2w
\end{matrix}
\right)
$$
$2m_1 \geq m_2$, $2m_2 \geq m_1 + 2m_3$, $2m_3 \geq m_2$이다. 또는
\item
\label{models-item-B3}
다음으로 주어지고,
$$
\left(
\begin{matrix}
m_1 \\
m_2 \\
m_3
\end{matrix}
\right),
\quad
\left(
\begin{matrix}
-4w & 2w & 0 \\
2w & -4w & 2w \\
0 & 2w & -2w
\end{matrix}
\right),
\quad
\left(
\begin{matrix}
2w \\
2w \\
w
\end{matrix}
\right)
$$
$2m_1 \geq m_2$, $2m_2 \geq m_1 + m_3$, $m_3 \geq m_2$이다.
\end{enumerate}
\end{lemma}

\begin{proof}
위의 논의를 참조하라.
\end{proof}

\noindent
$i$, $j$, $k$, $l$이 $a_{ij} > 0$, $a_{jk} > 0$, $a_{kl} > 0$인 네
$(-2)$-지표라고 하자. 즉 지표 $i$는 $j$와, $j$는 $k$와, $k$는 $l$과
``만난다''. 보조정리 \ref{models-lemma-three-by-three}에서
$a_{ik} = a_{jl} = 0$임을 알 수 있다. 행렬 $A = (a_{ij})$가 반음의
정부호이므로 행렬
$$
\left(
\begin{matrix}
-2w_i & a_{ij} & 0 & a_{il} \\
a_{ij} & -2w_j & a_{jk} & 0 \\
0 & a_{jk} & -2w_k & a_{kl} \\
a_{il} & 0 & a_{kl} & -2w_l
\end{matrix}
\right)
$$
은 $n = 4$가 아닌 한 음의 정부호이다. $n = 4$인 경우도 가능하며,
이때 행렬식\footnote{그 값은
$16w_iw_jw_kw_l - 4a_{ij}^2w_kw_l - 4a_{jk}^2w_iw_l - 4a_{kl}^2w_iw_j -
4a_{il}^2w_jw_k + a_{ij}^2a_{kl}^2 + a_{jk}^2a_{il}^2 -
2a_{ij}a_{il}a_{jk}a_{kl}$이다.}은 영이어서 방정식
$$
16 +
\frac{a_{ij}^2}{w_iw_j}\frac{a_{kl}^2}{w_kw_l} +
\frac{a_{jk}^2}{w_jw_k}\frac{a_{il}^2}{w_iw_l} =
4\frac{a_{ij}^2}{w_iw_j} + 4\frac{a_{jk}^2}{w_jw_k} + 4\frac{a_{kl}^2}{w_kw_l}
+ 4\frac{a_{il}^2}{w_iw_l} + 2\frac{a_{ij}a_{il}a_{jk}a_{kl}}{w_iw_jw_kw_l}
$$
을 얻는다. 이 방정식 오른쪽의 마지막 항은 다음 등식 때문에 나머지 항들에
의해 결정된다.
$$
\left(\frac{a_{ij}a_{il}a_{jk}a_{kl}}{w_iw_jw_kw_l}\right)^2 =
\frac{a_{ij}^2}{w_iw_j} \frac{a_{jk}^2}{w_jw_k}
\frac{a_{kl}^2}{w_kw_l} \frac{a_{il}^2}{w_iw_l}
$$
위에서
$\frac{a_{ij}^2}{w_iw_j}, \frac{a_{jk}^2}{w_jw_k}, \frac{a_{kl}^2}{w_kw_l}$은
$\{1, 2\}$에 속하고 $\frac{a_{il}^2}{w_iw_l}$은
$\{0, 1, 2\}$에 속함을 보았으므로 가능한 해는
$$
(\frac{a_{ij}^2}{w_iw_j}, \frac{a_{jk}^2}{w_jw_k}, \frac{a_{kl}^2}{w_kw_l},
\frac{a_{il}^2}{w_iw_l}) =
(1, 1, 1, 1) \text{ 또는 } (2, 1, 2, 0)
$$
뿐이고 $g = 1$인 경우이다. 독자는 이를 보조정리 \ref{models-lemma-genus-one}의
경우 (\ref{models-item-four-cycle}), (\ref{models-item-up-equal-up}),
(\ref{models-item-up-equal-down}), (\ref{models-item-down-equal-up})에서 찾을 수 있다.
$n > 4$라 가정하면 부등식
$$
16 +
\frac{a_{ij}^2}{w_iw_j}\frac{a_{kl}^2}{w_kw_l} +
\frac{a_{jk}^2}{w_jw_k}\frac{a_{il}^2}{w_iw_l} >
4\frac{a_{ij}^2}{w_iw_j} + 4\frac{a_{jk}^2}{w_jw_k} + 4\frac{a_{kl}^2}{w_kw_l}
+ 4\frac{a_{il}^2}{w_iw_l} + 2\frac{a_{ij}a_{il}a_{jk}a_{kl}}{w_iw_jw_kw_l}
$$
을 얻는다. 위 수들에 대한 제약을 사용하면 가능한 경우는
$$
(\frac{a_{ij}^2}{w_iw_j}, \frac{a_{jk}^2}{w_jw_k}, \frac{a_{kl}^2}{w_kw_l},
\frac{a_{il}^2}{w_iw_l}) =
(1, 1, 1, 0), (1, 1, 2, 0), (1, 2, 1, 0), \text{ 또는 }(2, 1, 1, 0)
$$
뿐이고, 특히 $a_{il} = 0$이다. 나머지 셋은 영이 아니라고 가정했음을
상기하라. $(2, 1, 1, 0)$은 $(1, 1, 2, 0)$에서 지표 $i,j,k,l$의 순서를
뒤집어 얻는다. 처음 세 해는 보조정리 \ref{models-lemma-four-by-four}의 경우
(\ref{models-item-A4}), (\ref{models-item-C4}), (\ref{models-item-B4}), (\ref{models-item-F4})를
준다. 이 보조정리에서는 정수 $m_i$, $m_j$, $m_k$, $m_l$에 대한 결과도 명시하였다.

\begin{lemma}
\label{models-lemma-four-by-four}
다음 형태의 진부분그래프를 분류한다.
$$
\xymatrix{
\bullet \ar@{-}[r] &
\bullet \ar@{-}[r] &
\bullet \ar@{-}[r] &
\bullet
}
$$
$n > 4$이고 $a_{ij}, a_{jk}, a_{kl}$이 영이 아닌 네 $(-2)$-지표
$i,j,k,l$이 주어지면, 순서를 바꾸는 것까지 동일시하여 $m$들, $a$들,
$w$들은 다음과 같다.
\begin{enumerate}
\item
\label{models-item-A4}
다음으로 주어지고,
$$
\left(
\begin{matrix}
m_1 \\
m_2 \\
m_3 \\
m_4
\end{matrix}
\right),
\quad
\left(
\begin{matrix}
-2w & w & 0 & 0 \\
w & -2w & w & 0 \\
0 & w & -2w & w \\
0 & 0 & w & -2w 
\end{matrix}
\right),
\quad
\left(
\begin{matrix}
w \\
w \\
w \\
w
\end{matrix}
\right)
$$
$2m_1 \geq m_2$, $2m_2 \geq m_1 + m_3$, $2m_3 \geq m_2 + m_4$이고
$2m_4 \geq m_3$이다. 또는
\item
\label{models-item-C4}
다음으로 주어지고,
$$
\left(
\begin{matrix}
m_1 \\
m_2 \\
m_3 \\
m_4
\end{matrix}
\right),
\quad
\left(
\begin{matrix}
-2w & w & 0 & 0 \\
w & -2w & w & 0 \\
0 & w & -2w & 2w \\
0 & 0 & 2w & -4w 
\end{matrix}
\right),
\quad
\left(
\begin{matrix}
w \\
w \\
w \\
2w
\end{matrix}
\right)
$$
$2m_1 \geq m_2$, $2m_2 \geq m_1 + m_3$, $2m_3 \geq m_2 + 2m_4$이고
$2m_4 \geq m_3$이다. 또는
\item
\label{models-item-B4}
다음으로 주어지고,
$$
\left(
\begin{matrix}
m_1 \\
m_2 \\
m_3 \\
m_4
\end{matrix}
\right),
\quad
\left(
\begin{matrix}
-4w & 2w & 0 & 0 \\
2w & -4w & 2w & 0 \\
0 & 2w & -4w & 2w \\
0 & 0 & 2w & -2w 
\end{matrix}
\right),
\quad
\left(
\begin{matrix}
2w \\
2w \\
2w \\
w
\end{matrix}
\right)
$$
$2m_1 \geq m_2$, $2m_2 \geq m_1 + m_3$, $2m_3 \geq m_2 + m_4$이고
$m_4 \geq m_3$이다. 또는
\item
\label{models-item-F4}
다음으로 주어지고,
$$
\left(
\begin{matrix}
m_1 \\
m_2 \\
m_3 \\
m_4
\end{matrix}
\right),
\quad
\left(
\begin{matrix}
-2w & w & 0 & 0 \\
w & -2w & 2w & 0 \\
0 & 2w & -4w & 2w \\
0 & 0 & 2w & -4w 
\end{matrix}
\right),
\quad
\left(
\begin{matrix}
w \\
w \\
2w \\
2w
\end{matrix}
\right)
$$
$2m_1 \geq m_2$, $2m_2 \geq m_1 + 2m_3$, $2m_3 \geq m_2 + m_4$이고
$2m_4 \geq m_3$이다.
\end{enumerate}
\end{lemma}

\begin{proof}
위의 논의를 참조하라.
\end{proof}

\noindent
$i$, $j$, $k$, $l$이 $a_{ij} > 0$, $a_{ij} > 0$, $a_{il} > 0$인 네
$(-2)$-지표라고 하자. 즉 지표 $i$는 지표 $j$, $k$, $l$과 ``만난다''.
보조정리 \ref{models-lemma-three-by-three}에서
$a_{jk} = a_{jl} = a_{kl} = 0$임을 알 수 있다. 행렬 $A = (a_{ij})$가
반음의 정부호이므로 행렬
$$
\left(
\begin{matrix}
-2w_i & a_{ij} & a_{ik} & a_{il} \\
a_{ij} & -2w_j & 0 & 0 \\
a_{ik} & 0 & -2w_k & 0 \\
a_{il} & 0 & 0 & -2w_l
\end{matrix}
\right)
$$
은 $n = 4$가 아닌 한 음의 정부호이다. $n = 4$인 경우도 가능하며,
이때 행렬식\footnote{그 값은
$16w_iw_jw_kw_l - 4a_{ij}^2w_kw_l - 4a_{ik}^2w_jw_l - 4a_{il}^2w_jw_k$.}
은 영이어서 방정식
$$
4 = \frac{a_{ij}^2}{w_iw_j} + \frac{a_{ik}^2}{w_iw_k} + \frac{a_{il}^2}{w_jw_l}
$$
을 얻는다. 위에서
$\frac{a_{ij}^2}{w_iw_j}, \frac{a_{ik}^2}{w_iw_k}, \frac{a_{il}^2}{w_iw_l}$
가 $\{1, 2\}$에 속함을 보았으므로, 순서를 바꾸는 것까지 동일시하여
유일한 가능성은 $4 = 1 + 1 + 2$이다. 각 경우에 $g = 1$이고, 보조정리
\ref{models-lemma-genus-one}의 경우 (\ref{models-item-triple-with-up}) 및
(\ref{models-item-triple-with-down})에서 찾을 수 있다. $n > 4$라 가정하면 부등식
$$
4 > \frac{a_{ij}^2}{w_iw_j} + \frac{a_{ik}^2}{w_iw_k} + \frac{a_{il}^2}{w_jw_l}
$$
을 얻는다. 이는 $\frac{a_{ij}^2}{w_iw_j} =
\frac{a_{ik}^2}{w_iw_k} = \frac{a_{il}^2}{w_jw_l} = 1$ 및
$w_i = w_j = w_k = w_l$을 함의한다. 이에 따라 보조정리
\ref{models-lemma-D4}의 경우 (\ref{models-item-D4})를 얻는다. 그 보조정리에서는 정수
$m_i$, $m_j$, $m_k$, $m_l$에 대한 결과도 명시하였다.

\begin{lemma}
\label{models-lemma-D4}
다음 형태의 진부분그래프를 분류한다.
$$
\xymatrix{
\bullet \ar@{-}[r] & \bullet \ar@{-}[r] \ar@{-}[d] & \bullet \\
& \bullet
}
$$
$n > 4$이고 $a_{ij}, a_{ik}, a_{il}$이 영이 아닌 네 $(-2)$-지표
$i,j,k,l$이 주어지면, 순서를 바꾸는 것까지 동일시하여 $m$들, $a$들,
$w$들은 다음과 같다.
\begin{enumerate}
\item
\label{models-item-D4}
다음으로 주어지고,
$$
\left(
\begin{matrix}
m_1 \\
m_2 \\
m_3 \\
m_4
\end{matrix}
\right),
\quad
\left(
\begin{matrix}
-2w & w & w & w \\
w & -2w & 0 & 0 \\
w & 0 & -2w & 0 \\
w & 0 & 0 & -2w 
\end{matrix}
\right),
\quad
\left(
\begin{matrix}
w \\
w \\
w \\
w
\end{matrix}
\right)
$$
$2m_1 \geq m_2 + m_3 + m_4$, $2m_2 \geq m_1$, $2m_3 \geq m_1$,
$2m_4 \geq m_1$이다. 이것이 $m_1 \geq \max(m_2, m_3, m_4)$를
함의함에 유의하라.
\end{enumerate}
\end{lemma}

\begin{proof}
위의 논의를 참조하라.
\end{proof}

\noindent
$h$, $i$, $j$, $k$, $l$이 $a_{hi} > 0$, $a_{ij} > 0$, $a_{jk} > 0$,
$a_{kl} > 0$인 다섯 $(-2)$-지표라고 하자. 즉 $h$는 $i$와, $i$는
$j$와, $j$는 $k$와, $k$는 $l$과 ``만난다''. 보조정리
\ref{models-lemma-three-by-three}과 \ref{models-lemma-four-by-four}를 적용하면
$a_{hj} = a_{hk} = a_{ik} = a_{il} = a_{jl} = 0$
이고, 분수
$\frac{a_{hi}^2}{w_hw_i}, \frac{a_{ij}^2}{w_iw_j}, \frac{a_{jk}^2}{w_jw_k},
\frac{a_{kl}^2}{w_kw_l}$는 $\{1, 2\}$에 속하고 분수
$\frac{a_{hl}^2}{w_hw_l} \in \{0, 1, 2\}$.
임을 알 수 있다. 행렬 $A = (a_{ij})$가 반음의 정부호이므로 행렬
$$
\left(
\begin{matrix}
-2w_h & a_{hi} & 0 & 0 & a_{hl} \\
a_{hi} & -2w_i & a_{ij} & 0 & 0 \\
0 & a_{ij} & -2w_j & a_{jk} & 0 \\
0 & 0 & a_{jk} & -2w_k & a_{kl} \\
a_{hl} & 0 & 0 & a_{kl} & -2w_l
\end{matrix}
\right)
$$
은 $n = 5$가 아닌 한 음의 정부호이다. $n = 5$인 경우도 가능하며,
이때 행렬식\footnote{그 값은
$-32w_hw_iw_jw_kw_l +
8a_{hi}^2w_jw_kw_l +
8a_{ij}^2w_hw_kw_l +
8a_{jk}^2w_hw_iw_l +
8a_{kl}^2w_hw_iw_j +
8a_{hl}^2w_iw_jw_k -
2a_{hi}^2a_{jk}^2w_l -
2a_{hi}^2a_{kl}^2w_j -
2a_{ij}^2a_{kl}^2w_h -
2a_{hl}^2a_{ij}^2w_k -
2a_{hl}^2a_{jk}^2w_i +
2a_{hi}a_{ij}a_{jk}a_{kl}a_{hl}
$
.}
행렬의 행렬식은 영이어서 방정식
\begin{align*}
16 + 
\frac{a_{hi}^2}{w_hw_i}\frac{a_{jk}^2}{w_jw_k} +
\frac{a_{hi}^2}{w_hw_i}\frac{a_{kl}^2}{w_kw_l} +
\frac{a_{ij}^2}{w_iw_j}\frac{a_{kl}^2}{w_kw_l} +
\frac{a_{hl}^2}{w_hw_l}\frac{a_{ij}^2}{w_iw_j} +
\frac{a_{hl}^2}{w_hw_l}\frac{a_{jk}^2}{w_jw_k} \\
= 4\frac{a_{hi}^2}{w_hw_i} +
4\frac{a_{ij}^2}{w_iw_j} +
4\frac{a_{jk}^2}{w_jw_k} +
4\frac{a_{kl}^2}{w_kw_l} +
4\frac{a_{hl}^2}{w_hw_l} +
\frac{a_{hi}a_{ij}a_{jk}a_{kl}a_{hl}}{w_hw_iw_jw_kw_l}
\end{align*}
을 얻는다. 이 방정식 오른쪽의 마지막 항은 다음 등식 때문에 나머지 항들에
의해 결정된다.
$$
\left(\frac{a_{hi}a_{ij}a_{jk}a_{kl}a_{hl}}{w_hw_iw_jw_kw_l} \right)^2 =
\frac{a_{hi}^2}{w_hw_i} \frac{a_{ij}^2}{w_iw_j}
\frac{a_{jk}^2}{w_jw_k} \frac{a_{kl}^2}{w_kw_l} \frac{a_{hl}^2}{w_hw_l}
$$
따라서 가능한 해는
$$
(\frac{a_{hi}^2}{w_hw_i}, \frac{a_{ij}^2}{w_iw_j},
\frac{a_{jk}^2}{w_jw_k}, \frac{a_{kl}^2}{w_kw_l}, \frac{a_{hl}^2}{w_hw_l}) =
(1, 1, 1, 1, 1), (1, 1, 2, 1, 0), (1, 2, 1, 1, 0), \text{ 또는 }(2, 1, 1, 2, 0)
$$
$(1, 2, 1, 1, 0)$은 $(1, 1, 2, 1, 0)$에서 지표 $h,i,j,k,l$의 순서를
뒤집어 얻는다. 이 경우들에는 $g = 1$이고, 보조정리 \ref{models-lemma-genus-one}의
경우 (\ref{models-item-five-cycle}), (\ref{models-item-equal-equal-up-equal}),
(\ref{models-item-equal-equal-down-equal}), (\ref{models-item-up-equal-equal-up}),
(\ref{models-item-up-equal-equal-down}), (\ref{models-item-down-equal-equal-up})에서
찾을 수 있다. 이 가운데 하나는 $(1, 1, 1, 1, 1)$에, 둘은
$(1, 1, 2, 1, 0)$에, 셋은 $(2, 1, 1, 2, 0)$에 대응한다.
$n > 5$라 가정하면 부등식
\begin{align*}
16 + 
\frac{a_{hi}^2}{w_hw_i}\frac{a_{jk}^2}{w_jw_k} +
\frac{a_{hi}^2}{w_hw_i}\frac{a_{kl}^2}{w_kw_l} +
\frac{a_{ij}^2}{w_iw_j}\frac{a_{kl}^2}{w_kw_l} +
\frac{a_{hl}^2}{w_hw_l}\frac{a_{ij}^2}{w_iw_j} +
\frac{a_{hl}^2}{w_hw_l}\frac{a_{jk}^2}{w_jw_k} \\
> 4\frac{a_{hi}^2}{w_hw_i} +
4\frac{a_{ij}^2}{w_iw_j} +
4\frac{a_{jk}^2}{w_jw_k} +
4\frac{a_{kl}^2}{w_kw_l} +
4\frac{a_{hl}^2}{w_hw_l} +
\frac{a_{hi}a_{ij}a_{jk}a_{kl}a_{hl}}{w_hw_iw_jw_kw_l}
\end{align*}
을 얻는다. 위 수들에 대한 제약을 사용하면 가능한 경우는
$$
(\frac{a_{hi}^2}{w_hw_i}, \frac{a_{ij}^2}{w_iw_j},
\frac{a_{jk}^2}{w_jw_k}, \frac{a_{kl}^2}{w_kw_l}, \frac{a_{hl}^2}{w_hw_l}) =
(1, 1, 1, 1, 0), (1, 1, 1, 2, 0), \text{ 또는 } (2, 1, 1, 1, 0)
$$
뿐이고, 특히 $a_{hl} = 0$이다. 나머지 넷은 영이 아니라고 가정했음을
상기하라. $(1, 1, 1, 2, 0)$은 $(2, 1, 1, 1, 0)$에서 지표
$h,i,j,k,l$의 순서를 뒤집어 얻는다. 처음 두 해는 보조정리
\ref{models-lemma-five-by-five}의 경우 (\ref{models-item-A5}), (\ref{models-item-C5}),
(\ref{models-item-B5})를 준다. 그 보조정리에서는 정수 $m_h$, $m_i$, $m_j$, $m_k$, $m_l$에
대한 결과도 명시하였다.

\begin{lemma}
\label{models-lemma-five-by-five}
다음 형태의 진부분그래프를 분류한다.
$$
\xymatrix{
\bullet \ar@{-}[r] &
\bullet \ar@{-}[r] &
\bullet \ar@{-}[r] &
\bullet \ar@{-}[r] &
\bullet
}
$$
$n > 5$이고 $a_{hi}, a_{ij}, a_{jk}, a_{kl}$이 영이 아닌 다섯
$(-2)$-지표 $h,i,j,k,l$이 주어지면, 순서를 바꾸는 것까지 동일시하여
$m$들, $a$들, $w$들은 다음과 같다.
\begin{enumerate}
\item
\label{models-item-A5}
다음으로 주어지고,
$$
\left(
\begin{matrix}
m_1 \\
m_2 \\
m_3 \\
m_4 \\
m_5
\end{matrix}
\right),
\quad
\left(
\begin{matrix}
-2w & w & 0 & 0 & 0 \\
w & -2w & w & 0 & 0 \\
0 & w & -2w & w & 0 \\
0 & 0 & w & -2w & w \\
0 & 0 & 0 & w & -2w
\end{matrix}
\right),
\quad
\left(
\begin{matrix}
w \\
w \\
w \\
w \\
w
\end{matrix}
\right)
$$
$2m_1 \geq m_2$, $2m_2 \geq m_1 + m_3$, $2m_3 \geq m_2 + m_4$,
$2m_4 \geq m_3 + m_5$, $2m_5 \geq m_4$이다. 또는
\item
\label{models-item-C5}
다음으로 주어지고,
$$
\left(
\begin{matrix}
m_1 \\
m_2 \\
m_3 \\
m_4 \\
m_5
\end{matrix}
\right),
\quad
\left(
\begin{matrix}
-2w & w & 0 & 0 & 0 \\
w & -2w & w & 0 & 0 \\
0 & w & -2w & w & 0 \\
0 & 0 & w & -2w & 2w \\
0 & 0 & 0 & 2w & -4w
\end{matrix}
\right),
\quad
\left(
\begin{matrix}
w \\
w \\
w \\
w \\
2w
\end{matrix}
\right)
$$
$2m_1 \geq m_2$, $2m_2 \geq m_1 + m_3$, $2m_3 \geq m_2 + 2m_4$,
$2m_4 \geq m_3 + m_5$, $2m_5 \geq m_4$이다. 또는
\item
\label{models-item-B5}
다음으로 주어지고,
$$
\left(
\begin{matrix}
m_1 \\
m_2 \\
m_3 \\
m_4 \\
m_5
\end{matrix}
\right),
\quad
\left(
\begin{matrix}
-4w & 2w & 0 & 0 & 0 \\
2w & -4w & 2w & 0 & 0 \\
0 & 2w & -4w & 2w & 0 \\
0 & 0 & 2w & -4w & 2w \\
0 & 0 & 0 & 2w & -2w
\end{matrix}
\right),
\quad
\left(
\begin{matrix}
2w \\
2w \\
2w \\
2w \\
w
\end{matrix}
\right)
$$
$2m_1 \geq m_2$, $2m_2 \geq m_1 + m_3$, $2m_3 \geq m_2 + m_4$,
$2m_4 \geq m_3 + m_5$, $m_4 \geq m_3$이다.
\end{enumerate}
\end{lemma}

\begin{proof}
위의 논의를 참조하라.
\end{proof}

\noindent
$h$, $i$, $j$, $k$, $l$이 $a_{hi} > 0$, $a_{hj} > 0$, $a_{hk} > 0$,
$a_{hl} > 0$인 다섯 $(-2)$-지표라고 하자. 즉 지표 $h$는
$i$, $j$, $k$, $l$과 ``만난다''. 보조정리 \ref{models-lemma-three-by-three}에서
$a_{ij} = a_{ik} = a_{il} = a_{jk} = a_{jl} = a_{kl} = 0$이고,
보조정리 \ref{models-lemma-D4}에서
$w_h = w_i = w_j = w_k = w_l = w$인 어떤 정수 $w > 0$가 있으며
$a_{hi} = a_{hj} = a_{hk} = a_{hl} = -2w$.
임을 안다. 이에 대응하는 행렬
$$
\left(
\begin{matrix}
-2w & w & w & w & w \\
w & -2w & 0 & 0 & 0 \\
w & 0 & -2w & 0 & 0 \\
w & 0 & 0 & -2w & 0 \\
w & 0 & 0 & 0 & -2w
\end{matrix}
\right)
$$
는 특이하다. 따라서 이는 $n = 5$이고 $g = 1$일 때만 가능하다.
보조정리 \ref{models-lemma-genus-one}의 경우 (\ref{models-item-quadruple})에서
찾을 수 있다.

\begin{lemma}
\label{models-lemma-fourfold}
다음 형태의 진부분그래프는 존재하지 않는다.
$$
\xymatrix{
\bullet \ar@{-}[r] & \bullet \ar@{-}[ld] \ar@{-}[r] \ar@{-}[d] & \bullet \\
\bullet & \bullet
}
$$
$n > 5$이면 $a_{hi} > 0$, $a_{hj} > 0$, $a_{hk} > 0$,
$a_{hl} > 0$인 다섯 $(-2)$-지표 $h$, $i$, $j$, $k$는 {\bf 존재하지 않는다}.
\end{lemma}

\begin{proof}
위의 논의를 참조하라.
\end{proof}

\noindent
$h$, $i$, $j$, $k$, $l$이 $a_{hi} > 0$, $a_{ij} > 0$, $a_{jk} > 0$,
$a_{jl} > 0$인 다섯 $(-2)$-지표라고 하자. 즉 $h$는 $i$와 만나고
$j$는 $i$, $k$, $l$과 만난다. 보조정리 \ref{models-lemma-D4}에 의해
$a_{ik} = a_{il} = a_{kl} = 0$, $w_i = w_j = w_k = w_l = w$이고,
어떤 정수 $w$에 대해 $a_{ij} = a_{jk} = a_{jl} = w$이다.
네 원소열 $h, i, j, k$와 $h, i, j, l$에 보조정리
\ref{models-lemma-four-by-four}를 적용하면 $a_{hj} = a_{hk} = a_{hl} = 0$이고,
$w_h = \frac{1}{2}w$, $w$, 또는 $2w$이며 이에 따라
$a_{hi} = w$, $w$, 또는 $2w$이다. $A$가 반음의 정부호이므로 행렬
$$
\left(
\begin{matrix}
-2w_h & a_{hi} & 0 & 0 & 0 \\
a_{hi} & -2w & w & 0 & 0 \\
0 & w & -2w & w & w \\
0 & 0 & w & -2w & 0 \\
0 & 0 & w & 0 & -2w
\end{matrix}
\right)
$$
은 $n = 5$가 아닌 한 음의 정부호이다. $w_h = \frac{1}{2}w$ 또는
$2w$일 때 이 행렬의 행렬식이 $0$임을 계산할 수 있다. 이에 따라
보조정리 \ref{models-lemma-genus-one}의 경우 (\ref{models-item-triple-extended-up})과
(\ref{models-item-triple-extended-down})를 얻는다. $w_h = w$이면 보조정리
\ref{models-lemma-D5}의 경우 (\ref{models-item-D5})를 얻는다.

\begin{lemma}
\label{models-lemma-D5}
다음 형태의 진부분그래프를 분류한다.
$$
\xymatrix{
\bullet \ar@{-}[r] & \bullet \ar@{-}[r] &
\bullet \ar@{-}[r] \ar@{-}[d] & \bullet \\
& & \bullet
}
$$
$n > 5$이고 $a_{hi}, a_{ij}, a_{jk}, a_{jl}$이 영이 아닌 다섯
$(-2)$-지표 $h,i,j,k,l$이 주어지면, 순서를 바꾸는 것까지 동일시하여
$m$들, $a$들, $w$들은 다음과 같다.
\begin{enumerate}
\item
\label{models-item-D5}
다음으로 주어지고,
$$
\left(
\begin{matrix}
m_1 \\
m_2 \\
m_3 \\
m_4 \\
m_5
\end{matrix}
\right),
\quad
\left(
\begin{matrix}
-2w & w & 0 & 0 & 0 \\
w & -2w & w & 0 & 0 \\
0 & w & -2w & w & w \\
0 & 0 & w & -2w & 0 \\
0 & 0 & w & 0 & -2w
\end{matrix}
\right),
\quad
\left(
\begin{matrix}
w \\
w \\
w \\
w \\
w
\end{matrix}
\right)
$$
$2m_1 \geq m_2$, $2m_2 \geq m_1 + m_3$, $2m_3 \geq m_2 + m_4 + m_5$,
$2m_4 \geq m_3$, $2m_5 \geq m_3$이다.
\end{enumerate}
\end{lemma}

\begin{proof}
위의 논의를 참조하라.
\end{proof}

\noindent
$t > 5$이고 $i_1, \ldots, i_t$가 서로 다른 $t$개의 $(-2)$-지표이며
$j = 1, \ldots, t - 1$에 대해 $a_{i_ji_{j + 1}}$가 영이 아니라고 하자.
$t$에 대한 귀납법으로 $n = t$이면 보조정리 \ref{models-lemma-genus-one}의
경우 (\ref{models-item-n-cycle}), (\ref{models-item-up-chain-equal-up}),
(\ref{models-item-up-chain-equal-down}), (\ref{models-item-down-chain-equal-up})를,
$n > t$이면 보조정리 \ref{models-lemma-long}의 경우 (\ref{models-item-An}),
(\ref{models-item-Cn}), (\ref{models-item-Bn})를 얻음을 증명한다. 먼저
$a_{i_1i_t}$가 영이 아니면 보조정리 \ref{models-lemma-five-by-five}에서
$w_{i_1} = \ldots = w_{i_t} = w$이고 $j = 1, \ldots, t - 1$에 대해
$a_{i_ji_{j + 1}} = w$, 또 $a_{i_1i_t} = w$임이 분명하다. 그러면 벡터
$(1, \ldots, 1)$은 대응하는 $t \times t$ 행렬의 핵에 속한다. 따라서
$n = t$이어야 하고, 종수는 $1$이며 보조정리 \ref{models-lemma-genus-one}의
경우 (\ref{models-item-n-cycle})이다. 그러므로 $a_{i_1i_t} = 0$이라 가정할 수
있다. 귀납가정에 의해($t = 6$이면 보조정리
\ref{models-lemma-five-by-five}에 의해) $k > j + 1$이면
$a_{i_ji_k} = 0$이다. 더욱이 어떤 정수 $w$에 대해
$w_{i_1} = \ldots = w_{i_{t - 1}} = w$이고
$w_{i_1}, w_{i_t} \in \{\frac{1}{2}w, w, 2w\}$이다. 또한 $w_{i_1}$과
$w_{i_t}$의 값이 각각 $\frac{1}{2}w$, $w$, $2w$인 데 따라
$a_{i_1i_2}$와 $a_{i_{t - 1}i_t}$의 값은 각각 $w$, $w$, $2w$이다.
따라서 $9$가지 가능성이 있고, 각 경우에 어떤 일이 일어나는지 쉽게 판정된다.
\begin{enumerate}
\item $(w_{i_1}, w_{i_t}) = (\frac{1}{2}w, \frac{1}{2}w)$이면 보조정리
\ref{models-lemma-genus-one}의 경우 (\ref{models-item-up-chain-equal-down})이다.
\item $(w_{i_1}, w_{i_t}) = (\frac{1}{2}w, w)$ 또는
$(w, \frac{1}{2}w)$이면 보조정리 \ref{models-lemma-long}의 경우
(\ref{models-item-Bn})이다.
\item $(w_{i_1}, w_{i_t}) = (\frac{1}{2}w, 2w)$ 또는
$(2w, \frac{1}{2}w)$이면 보조정리 \ref{models-lemma-genus-one}의 경우
(\ref{models-item-up-chain-equal-up})이다.
\item $(w_{i_1}, w_{i_t}) = (w, w)$이면 보조정리 \ref{models-lemma-long}의
경우 (\ref{models-item-An})이다.
\item $(w_{i_1}, w_{i_t}) = (w, 2w)$ 또는 $(2w, w)$이면 보조정리
\ref{models-lemma-long}의 경우 (\ref{models-item-Cn})이다.
\item $(w_{i_1}, w_{i_t}) = (2w, 2w)$이면 보조정리
\ref{models-lemma-genus-one}의 경우 (\ref{models-item-down-chain-equal-up})이다.
\end{enumerate}

\begin{lemma}
\label{models-lemma-long}
다음 형태의 진부분그래프를 분류한다.
$$
\xymatrix{
\bullet \ar@{-}[r] &
\bullet \ar@{-}[r] &
\bullet \ar@{..}[r] &
\bullet \ar@{-}[r] &
\bullet \ar@{-}[r] &
\bullet
}
$$
$t > 5$이고 $n > t$라 하자. $j = 1, \ldots, t - 1$에 대해
$a_{i_ji_{j + 1}}$가 영이 아닌 서로 다른 $t$개의 $(-2)$-지표
$i_1, \ldots, i_t$가 주어지면, 이 지표들의 순서를 뒤집는 것까지
동일시하여 $a$들과 $w$들은 다음과 같다.
\begin{enumerate}
\item
\label{models-item-An}
$w_{i_1} = w_{i_2} = \ldots = w_{i_t} = w$,
$a_{i_ji_{j + 1}} = w$이고 $k > j + 1$이면 $a_{i_ji_k} = 0$이다. 또는
\item
\label{models-item-Cn}
$w_{i_1} = w_{i_2} = \ldots = w_{i_{t - 1}} = w$, $w_{j_t} = 2w$이고,
$j < t - 1$에 대해 $a_{i_ji_{j + 1}} = w$, $a_{i_{t - 1}i_t} = 2w$이며
$k > j + 1$이면 $a_{i_ji_k} = 0$이다. 또는
\item
\label{models-item-Bn}
$w_{i_1} = w_{i_2} = \ldots = w_{i_{t - 1}} = 2w$, $w_{j_t} = w$,
$a_{i_ji_{j + 1}} = 2w$, $a_{i_{t - 1}i_t} = 2w$이고
$k > j + 1$이면 $a_{i_ji_k} = 0$이다.
\end{enumerate}
\end{lemma}

\begin{proof}
위의 논의를 참조하라.
\end{proof}

\noindent
$t > 4$이고 $i_1, \ldots, i_{t + 1}$이 서로 다른 $t + 1$개의
$(-2)$-지표이며, $j = 1, \ldots, t - 1$에 대해
$a_{i_ji_{j + 1}} > 0$이고 $a_{j_{t - 1}j_{t + 1}} > 0$이라 하자.
보조정리 \ref{models-lemma-Dn}의 그림을 참조하라. $t$에 대한 귀납법으로
$n = t + 1$이면 보조정리 \ref{models-lemma-genus-one}의 경우
(\ref{models-item-Dn-extended-up}), (\ref{models-item-Dn-extended-down})를,
$n > t + 1$이면 보조정리 \ref{models-lemma-Dn}의 경우 (\ref{models-item-Dn})를
얻음을 증명한다. 귀납가정에 의해($t = 5$이면 보조정리
\ref{models-lemma-D5}에 의해) $j, k \geq 2$에 대해 요구된 영 아닌 항을
제외한 $a_{i_ji_k}$는 모두 영이다. 또한 어떤 정수 $w$에 대해
$w_2 = \ldots = w_{t + 1} = w$이고, $j, k \geq 2$에 대해 영이 아닌
$a_{i_ji_k}$는 $w$와 같다. 수열 $i_1, \ldots, i_t$와
$i_1, \ldots, i_{t - 1}, i_{t + 1}$에 보조정리 \ref{models-lemma-long}을
적용하면($t = 5$이면 보조정리 \ref{models-lemma-five-by-five}를 적용하면),
$j \geq 3$에 대해 $a_{i_1 i_j} = 0$이고 $w_1$은
$\frac{1}{2}w$, $w$, 또는 $2w$이며 이에 따라 $a_{i_1i_2}$는
$w, w, 2w$이다. 따라서 $3$가지 가능성이 있고 각 경우는 쉽게 판정된다.
\begin{enumerate}
\item $w_1 = \frac{1}{2}w$이면 보조정리 \ref{models-lemma-genus-one}의 경우
(\ref{models-item-Dn-extended-down})이다.
\item $w_1 = w$이면 보조정리 \ref{models-lemma-Dn}의 경우
(\ref{models-item-Dn})이다.
\item $w_1 = 2w$이면 보조정리 \ref{models-lemma-genus-one}의 경우
(\ref{models-item-Dn-extended-up})이다.
\end{enumerate}

\begin{lemma}
\label{models-lemma-Dn}
다음 형태의 진부분그래프를 분류한다.
$$
\xymatrix{
\bullet \ar@{-}[r] & \bullet \ar@{..}[r] & \bullet \ar@{-}[r] &
\bullet \ar@{-}[r] \ar@{-}[d] & \bullet \\
& & & \bullet
}
$$
$t > 4$이고 $n > t + 1$이라 하자. $j = 1, \ldots, t - 1$에 대해
$a_{i_ji_{j + 1}}$가 영이 아니고 $a_{i_{t - 1}i_{t + 1}}$도 영이 아닌,
서로 다른 $t + 1$개의 $(-2)$-지표 $i_1, \ldots, i_{t + 1}$이 주어지면
$a$들과 $w$들은 다음과 같다.
\begin{enumerate}
\item
\label{models-item-Dn}
$w_{i_1} = w_{i_2} = \ldots = w_{i_{t + 1}} = w$이고,
$j = 1, \ldots, t - 1$에 대해 $a_{i_ji_{j + 1}} = w$이며,
$a_{i_{t - 1}i_{t + 1}} = w$이고 $j > k$인 나머지 쌍 $(j, k)$에 대해
$a_{i_ji_k} = 0$이다.
\end{enumerate}
\end{lemma}

\begin{proof}
위의 논의를 참조하라.
\end{proof}

\noindent
$a_{gh}, a_{hi}, a_{ij}, a_{jk}, a_{il}$이 영이 아닌 서로 다른 $6$개의
$(-2)$-지표 $g,h,i,j,k,l$이 주어졌다고 하자. 보조정리
\ref{models-lemma-E6}의 그림을 참조하라. 보조정리 \ref{models-lemma-D5}를 적용하면
보조정리 \ref{models-lemma-E6}의 상황이어야 한다. 행렬식이 $3w^6 > 0$이므로
이 경우에는 결코 $n = 6$일 수 없다.

\begin{lemma}
\label{models-lemma-E6}
다음 형태의 진부분그래프를 분류한다.
$$
\xymatrix{
\bullet \ar@{-}[r] & \bullet \ar@{-}[r] & \bullet \ar@{-}[r] \ar@{-}[d] &
\bullet \ar@{-}[r] & \bullet \\
& & \bullet
}
$$
$n > 6$이라 하자. $a_{12}, a_{23}, a_{34}, a_{45}, a_{36}$이 영이 아닌
서로 다른 $6$개의 $(-2)$-지표 $i_1, \ldots, i_6$이 주어지면 $m$들,
$a$들, $w$들은 다음과 같다.
\begin{enumerate}
\item
\label{models-item-E6}
다음으로 주어지고,
$$
\left(
\begin{matrix}
m_1 \\
m_2 \\
m_3 \\
m_4 \\
m_5 \\
m_6
\end{matrix}
\right),
\quad
\left(
\begin{matrix}
-2w & w & 0 & 0 & 0 & 0 \\
w & -2w & w & 0 & 0 & 0 \\
0 & w & -2w & w & 0 & w \\
0 & 0 & w & -2w & w & 0 \\
0 & 0 & 0 & w & -2w & 0 \\
0 & 0 & w & 0 & 0 & -2w
\end{matrix}
\right),
\quad
\left(
\begin{matrix}
w \\
w \\
w \\
w \\
w \\
w
\end{matrix}
\right)
$$
$2m_1 \geq m_2$, $2m_2 \geq m_1 + m_3$, $2m_3 \geq m_2 + m_4 + m_6$,
$2m_4 \geq m_3 + m_5$, $2m_5 \geq m_3$, $2m_6 \geq m_3$이다.
\end{enumerate}
\end{lemma}

\begin{proof}
위의 논의를 참조하라.
\end{proof}

\noindent
$t \geq 4$이고 $i_0, \ldots, i_{t + 1}$이 서로 다른 $t + 2$개의
$(-2)$-지표이며, $j = 1, \ldots, t - 1$에 대해
$a_{i_ji_{j + 1}} > 0$이고 $a_{i_0i_2} > 0$ 및
$a_{i_{t - 1}i_{t + 1}} > 0$이라 하자. 보조정리
\ref{models-lemma-double-triple}의 그림을 참조하라. 보조정리 \ref{models-lemma-D5}와
\ref{models-lemma-Dn}을 적용하면 $j < k$인 나머지 $a_{i_ji_k}$는 모두 영이고,
어떤 정수 $w$에 대해 $w_{i_0} = \ldots = w_{i_{t + 1}} = w$이며,
$A$에서 요구된 영 아닌 비대각 성분은 $w$와 같다. 계산하면 대응하는
행렬의 행렬식은 영이다. 따라서 $n = t + 2$이고 보조정리
\ref{models-lemma-genus-one}의 경우 (\ref{models-item-double-triple})이다.

\begin{lemma}
\label{models-lemma-double-triple}
다음 형태의 진부분그래프는 존재하지 않는다.
$$
\xymatrix{
\bullet \ar@{-}[r] & \bullet \ar@{..}[r] \ar@{-}[d] &
\bullet \ar@{-}[d] \ar@{-}[r] & \bullet \\
& \bullet & \bullet
}
$$
$t \geq 4$이고 $n > t + 2$라 가정하자. $j = 1, \ldots, t - 1$에 대해
$a_{i_ji_{j + 1}} > 0$이고 $a_{i_0i_2} > 0$ 및
$a_{i_{t - 1}i_{t + 1}} > 0$인 서로 다른 $t + 2$개의 $(-2)$-지표
$i_0, \ldots, i_{t + 1}$은 {\bf 존재하지 않는다}.
\end{lemma}

\begin{proof}
위의 논의를 참조하라.
\end{proof}

\noindent
$a_{fg}, a_{gh}, a_{ij}, a_{jh}, a_{kl}, a_{lh}$이 영이 아닌 서로 다른
$7$개의 $(-2)$-지표 $f,g,h,i,j,k,l$이 주어졌다고 하자. 보조정리
\ref{models-lemma-E6-completed}의 그림을 참조하라. 보조정리 \ref{models-lemma-D5}를
적용하면 대응하는 행렬은
$$
\left(
\begin{matrix}
-2w & w & 0 & 0 & 0 & 0 & 0 \\
w & -2w & w & 0 & 0 & 0 & 0 \\
0 & w & -2w & 0 & w & 0 & w \\
0 & 0 & 0 & -2w & w & 0 & 0 \\
0 & 0 & w & w & -2w & 0 & 0 \\
0 & 0 & 0 & 0 & 0 & -2w & w \\
0 & 0 & w & 0 & 0 & w & -2w
\end{matrix}
\right)
$$
이다. 행렬식이 $0$이므로 $n = 7$이고 $g = 1$이어야 하며, 보조정리
\ref{models-lemma-genus-one}의 경우 (\ref{models-item-E6-completed})를 얻는다.

\begin{lemma}
\label{models-lemma-E6-completed}
다음 형태의 진부분그래프는 존재하지 않는다.
$$
\xymatrix{
\bullet \ar@{-}[r] & \bullet \ar@{-}[r] & \bullet \ar@{-}[r] \ar@{-}[d] &
\bullet \ar@{-}[r] & \bullet \\
& & \bullet \ar@{-}[d] \\
& & \bullet
}
$$
$n > 7$이라 가정하자. $a_{fg}, a_{gh}, a_{ij}, a_{jh}, a_{kl}, a_{lh}$이
영이 아닌 서로 다른 $7$개의 $(-2)$-지표 $f,g,h,i,j,k,l$은
{\bf 존재하지 않는다}.
\end{lemma}

\begin{proof}
위의 논의를 참조하라.
\end{proof}

\noindent
$a_{fg}, a_{gh}, a_{hi}, a_{ij}, a_{jk}, a_{il}$이 영이 아닌 서로 다른
$7$개의 $(-2)$-지표 $f,g,h,i,j,k,l$이 주어졌다고 하자. 보조정리
\ref{models-lemma-E7}의 그림을 참조하라. 보조정리 \ref{models-lemma-D5}와
\ref{models-lemma-Dn}을 적용하면 보조정리 \ref{models-lemma-E7}의 상황이어야 한다.
행렬식이 $-8w^7 > 0$이므로 이 경우에는 결코 $n = 7$일 수 없다.

\begin{lemma}
\label{models-lemma-E7}
다음 형태의 진부분그래프를 분류한다.
$$
\xymatrix{
\bullet \ar@{-}[r] & \bullet \ar@{-}[r] & \bullet \ar@{-}[r] &
\bullet \ar@{-}[r] \ar@{-}[d] & \bullet \ar@{-}[r] & \bullet \\
& & & \bullet
}
$$
$n > 7$이라 하자. $a_{12}, a_{23}, a_{34}, a_{45}, a_{56}, a_{47}$이
영이 아닌 서로 다른 $7$개의 $(-2)$-지표 $i_1, \ldots, i_7$이 주어지면
$m$들, $a$들, $w$들은 다음과 같다.
\begin{enumerate}
\item
\label{models-item-E7}
다음으로 주어지고,
$$
\left(
\begin{matrix}
m_1 \\
m_2 \\
m_3 \\
m_4 \\
m_5 \\
m_6 \\
m_7
\end{matrix}
\right),
\quad
\left(
\begin{matrix}
-2w & w & 0 & 0 & 0 & 0 & 0 \\
w & -2w & w & 0 & 0 & 0 & 0 \\
0 & w & -2w & w & 0 & 0 & 0 \\
0 & 0 & w & -2w & w & 0 & w \\
0 & 0 & 0 & w & -2w & w & 0 \\
0 & 0 & 0 & 0 & w & -2w & 0 \\
0 & 0 & 0 & w & 0 & 0 & -2w
\end{matrix}
\right),
\quad
\left(
\begin{matrix}
w \\
w \\
w \\
w \\
w \\
w \\
w
\end{matrix}
\right)
$$
$2m_1 \geq m_2$, $2m_2 \geq m_1 + m_3$, $2m_3 \geq m_2 + m_4$,
$2m_4 \geq m_3 + m_5 + m_7$, $2m_5 \geq m_4 + m_6$,
$2m_6 \geq m_5$, $2m_7 \geq m_4$이다.
\end{enumerate}
\end{lemma}

\begin{proof}
위의 논의를 참조하라.
\end{proof}

\noindent
행렬 $A$의 영이 아닌 성분 $a_{ij}$의 무늬가
$$
\xymatrix{
\bullet \ar@{-}[r] & \bullet \ar@{-}[r] & \bullet \ar@{-}[r] &
\bullet \ar@{-}[r] & \bullet \ar@{-}[r] \ar@{-}[d] &
\bullet \ar@{-}[r] & \bullet \\
& & & & \bullet
}
$$
또는
$$
\xymatrix{
\bullet \ar@{-}[r] & \bullet \ar@{-}[r] & \bullet \ar@{-}[r] &
\bullet \ar@{-}[r] \ar@{-}[d] & \bullet \ar@{-}[r] &
\bullet \ar@{-}[r] & \bullet \\
& & & \bullet
}
$$
와 같은 서로 다른 $8$개의 $(-2)$-지표가 주어졌다고 하자. 보조정리
\ref{models-lemma-E7}의 증명과 똑같이 논증하면 첫째 무늬는 보조정리
\ref{models-lemma-E8}의 경우 (\ref{models-item-E8})를 주지만 보조정리
\ref{models-lemma-genus-one}에는 새 경우를 주지 않는다. 보조정리
\ref{models-lemma-E6-completed}의 증명과 똑같이 논증하면 둘째 무늬는
$n > 8$일 때 나타나지 않지만 $n = 8$일 때 보조정리
\ref{models-lemma-genus-one}의 경우 (\ref{models-item-E7-completed})를 준다.

\begin{lemma}
\label{models-lemma-E8}
다음 형태의 진부분그래프를 분류한다.
$$
\xymatrix{
\bullet \ar@{-}[r] & \bullet \ar@{-}[r] & \bullet \ar@{-}[r] &
\bullet \ar@{-}[r] & \bullet \ar@{-}[r] \ar@{-}[d] &
\bullet \ar@{-}[r] & \bullet \\
& & & & \bullet
}
$$
$n > 8$이라 하자. $a_{12}, a_{23}, a_{34}, a_{45}, a_{56}, a_{65}, a_{57}$이
영이 아닌 서로 다른 $8$개의 $(-2)$-지표 $i_1, \ldots, i_8$이 주어지면
$m$들, $a$들, $w$들은 다음과 같다.
\begin{enumerate}
\item
\label{models-item-E8}
다음으로 주어지고,
$$
\left(
\begin{matrix}
m_1 \\
m_2 \\
m_3 \\
m_4 \\
m_5 \\
m_6 \\
m_7 \\
m_8
\end{matrix}
\right),
\quad
\left(
\begin{matrix}
-2w & w & 0 & 0 & 0 & 0 & 0 & 0 \\
w & -2w & w & 0 & 0 & 0 & 0 & 0 \\
0 & w & -2w & w & 0 & 0 & 0 & 0 \\
0 & 0 & w & -2w & w & 0 & 0 & 0 \\
0 & 0 & 0 & w & -2w & w & 0 & w \\
0 & 0 & 0 & 0 & w & -2w & w & 0 \\
0 & 0 & 0 & 0 & 0 & w & -2w & 0 \\
0 & 0 & 0 & 0 & w & 0 & 0 & -2w
\end{matrix}
\right),
\quad
\left(
\begin{matrix}
w \\
w \\
w \\
w \\
w \\
w \\
w \\
w
\end{matrix}
\right)
$$
$2m_1 \geq m_2$, $2m_2 \geq m_1 + m_3$, $2m_3 \geq m_2 + m_4$,
$2m_4 \geq m_3 + m_5$, $2m_5 \geq m_4 + m_6 + m_8$,
$2m_6 \geq m_5 + m_7$, $2m_7 \geq m_6$, $2m_8 \geq m_5$이다.
\end{enumerate}
\end{lemma}

\begin{proof}
위의 논의를 참조하라.
\end{proof}

\begin{lemma}
\label{models-lemma-E7-completed}
다음 형태의 진부분그래프는 존재하지 않는다.
$$
\xymatrix{
\bullet \ar@{-}[r] &
\bullet \ar@{-}[r] &
\bullet \ar@{-}[r] & \bullet \ar@{-}[r] \ar@{-}[d] &
\bullet \ar@{-}[r] & \bullet \ar@{-}[r] & \bullet \\
& & & \bullet
}
$$
$n > 8$이라 가정하자. $a_{ef}, a_{fg}, a_{gh}, a_{hi}, a_{ij}, a_{jk}, a_{lh}$이
영이 아닌 서로 다른 $8$개의 $(-2)$-지표 $e,f,g,h,i,j,k,l$은
{\bf 존재하지 않는다}.
\end{lemma}

\begin{proof}
위의 논의를 참조하라.
\end{proof}

\noindent
행렬 $A$의 영이 아닌 성분 $a_{ij}$의 무늬가
$$
\xymatrix{
\bullet \ar@{-}[r] & \bullet \ar@{-}[r] &
\bullet \ar@{-}[r] & \bullet \ar@{-}[r] &
\bullet \ar@{-}[r] & \bullet \ar@{-}[r] \ar@{-}[d] &
\bullet \ar@{-}[r] & \bullet \\
& & & & & \bullet
}
$$
와 같은 서로 다른 $9$개의 $(-2)$-지표가 주어졌다고 하자. 보조정리
\ref{models-lemma-E6-completed}의 증명과 똑같이 논증하면 이 무늬는
$n > 9$일 때 나타나지 않지만 $n = 9$일 때 보조정리
\ref{models-lemma-genus-one}의 경우 (\ref{models-item-E8-completed})를 준다.

\begin{lemma}
\label{models-lemma-E8-completed}
다음 형태의 진부분그래프는 존재하지 않는다.
$$
\xymatrix{
\bullet \ar@{-}[r] & \bullet \ar@{-}[r] &
\bullet \ar@{-}[r] & \bullet \ar@{-}[r] &
\bullet \ar@{-}[r] & \bullet \ar@{-}[r] \ar@{-}[d] &
\bullet \ar@{-}[r] & \bullet \\
& & & & & \bullet
}
$$
$n > 9$이라 가정하자. $a_{de}, a_{ef}, a_{fg}, a_{gh}, a_{hi}, a_{ij}, a_{jk}, a_{lh}$이
영이 아닌 서로 다른 $9$개의 $(-2)$-지표 $d,e,f,g,h,i,j,k,l$은
{\bf 존재하지 않는다}.
\end{lemma}

\begin{proof}
위의 논의를 참조하라.
\end{proof}

\noindent
이 정보를 모두 모으면 다음을 얻는다.

\begin{proposition}
\label{models-proposition-classify-subgraphs}
$n, m_i, a_{ij}, w_i, g_i$를 종수 $g$인 수치형이라 하자.
$I \subset \{1, \ldots, n\}$를 크기가 $\geq 2$이고 $(-2)$-지표들로
이루어진 진부분집합으로서, $i' \in I$, $i \in I \setminus I'$일 때
$a_{i'i} = 0$이 되는 공집합이 아닌 진부분집합 $I' \subset I$가 없다고
하자. 그러면 순서를 바꾸는 것까지 동일시하여 $i, j \in I$에 대한
$m_i$들, $a_{ij}$들, $w_i$들은 보조정리 \ref{models-lemma-two-by-two},
\ref{models-lemma-three-by-three},
\ref{models-lemma-four-by-four},
\ref{models-lemma-D4},
\ref{models-lemma-five-by-five},
\ref{models-lemma-D5},
\ref{models-lemma-long},
\ref{models-lemma-Dn},
\ref{models-lemma-E6},
\ref{models-lemma-E7}, 또는
\ref{models-lemma-E8}.
\end{proposition}

\begin{proof}
이는 위의 논의에서 따른다. 절 \ref{models-section-classify-proper-subgraphs}의
첫머리 논의를 참조하라.
\end{proof}













\section{종수 0과 1인 극소형의 분류}
\label{models-section-classification-genus-one}

\noindent
절 제목이 내용을 모두 말해 준다.

\begin{lemma}[종수 0]
\label{models-lemma-genus-zero}
종수 0인 극소 수치형은 오직
$n = 1$, $m_1 = 1$, $a_{11} = 0$, $w_1 = 1$, $g_1 = 0$뿐이다.
\end{lemma}

\begin{proof}
보조정리 \ref{models-lemma-non-irreducible-minimal-type-genus-at-least-one}과
\ref{models-lemma-irreducible}에서 따른다.
\end{proof}

\begin{lemma}[종수 1]
\label{models-lemma-genus-one}
종수 1인 극소 수치형은 동치까지 다음과 같다.
\begin{enumerate}
\item
\label{models-item-one}
$n = 1$, $a_{11} = 0$, $g_1 = 1$이고 $m_1, w_1 \geq 1$은 임의이다.
\item
\label{models-item-two-cycle}
$n = 2$이고 $m_i, a_{ij}, w_i, g_i$는 다음과 같이 주어진다.
$$
\left(
\begin{matrix}
m \\
m
\end{matrix}
\right),
\quad
\left(
\begin{matrix}
-2w & 2w \\
2w & -2w
\end{matrix}
\right),
\quad
\left(
\begin{matrix}
w \\
w
\end{matrix}
\right),
\quad
\left(
\begin{matrix}
0 \\
0
\end{matrix}
\right)
$$
$w$와 $m$은 임의이다.
\item
\label{models-item-up4}
$n = 2$이고 $m_i, a_{ij}, w_i, g_i$는 다음과 같이 주어진다.
$$
\left(
\begin{matrix}
2m \\
m
\end{matrix}
\right),
\quad
\left(
\begin{matrix}
-2w & 4w \\
4w & -8w
\end{matrix}
\right),
\quad
\left(
\begin{matrix}
w \\
4w
\end{matrix}
\right),
\quad
\left(
\begin{matrix}
0 \\
0
\end{matrix}
\right)
$$
$w$와 $m$은 임의이다.
\item
\label{models-item-three-cycle}
$n = 3$이고 $m_i, a_{ij}, w_i, g_i$는 다음과 같이 주어진다.
$$
\left(
\begin{matrix}
m \\
m \\
m
\end{matrix}
\right),
\quad
\left(
\begin{matrix}
-2w & w & w \\
w & -2w & w \\
w & w & -2w
\end{matrix}
\right),
\quad
\left(
\begin{matrix}
w \\
w \\
w
\end{matrix}
\right),
\quad
\left(
\begin{matrix}
0 \\
0 \\
0
\end{matrix}
\right)
$$
$w$와 $m$은 임의이다.
\item
\label{models-item-equal-up3}
$n = 3$이고 $m_i, a_{ij}, w_i, g_i$는 다음과 같이 주어진다.
$$
\left(
\begin{matrix}
m \\
2m \\
m
\end{matrix}
\right),
\quad
\left(
\begin{matrix}
-2w & w & 0 \\
w & -2w & 3w \\
0 & 3w & -6w
\end{matrix}
\right),
\quad
\left(
\begin{matrix}
w \\
w \\
3w
\end{matrix}
\right),
\quad
\left(
\begin{matrix}
0 \\
0 \\
0
\end{matrix}
\right)
$$
$w$와 $m$은 임의이다.
\item
\label{models-item-equal-down3}
$n = 3$이고 $m_i, a_{ij}, w_i, g_i$는 다음과 같이 주어진다.
$$
\left(
\begin{matrix}
m \\
2m \\
3m
\end{matrix}
\right),
\quad
\left(
\begin{matrix}
-6w & 3w & 0 \\
3w & -6w & 3w \\
0 & 3w & -2w
\end{matrix}
\right),
\quad
\left(
\begin{matrix}
3w \\
3w \\
w
\end{matrix}
\right),
\quad
\left(
\begin{matrix}
0 \\
0 \\
0
\end{matrix}
\right)
$$
$w$와 $m$은 임의이다.
\item
\label{models-item-up-up}
$n = 3$이고 $m_i, a_{ij}, w_i, g_i$는 다음과 같이 주어진다.
$$
\left(
\begin{matrix}
2m \\
2m \\
m
\end{matrix}
\right),
\quad
\left(
\begin{matrix}
-2w & 2w & 0 \\
2w & -4w & 4w \\
0 & 4w & -8w
\end{matrix}
\right),
\quad
\left(
\begin{matrix}
w \\
2w \\
4w
\end{matrix}
\right),
\quad
\left(
\begin{matrix}
0 \\
0 \\
0
\end{matrix}
\right)
$$
$w$와 $m$은 임의이다.
\item
\label{models-item-up-down}
$n = 3$이고 $m_i, a_{ij}, w_i, g_i$는 다음과 같이 주어진다.
$$
\left(
\begin{matrix}
m \\
m \\
m
\end{matrix}
\right),
\quad
\left(
\begin{matrix}
-2w & 2w & 0 \\
2w & -4w & 2w \\
0 & 2w & -2w
\end{matrix}
\right),
\quad
\left(
\begin{matrix}
w \\
2w \\
w
\end{matrix}
\right),
\quad
\left(
\begin{matrix}
0 \\
0 \\
0
\end{matrix}
\right)
$$
$w$와 $m$은 임의이다.
\item
\label{models-item-down-up}
$n = 3$이고 $m_i, a_{ij}, w_i, g_i$는 다음과 같이 주어진다.
$$
\left(
\begin{matrix}
m \\
2m \\
m
\end{matrix}
\right),
\quad
\left(
\begin{matrix}
-4w & 2w & 0 \\
2w & -2w & 2w \\
0 & 2w & -4w
\end{matrix}
\right),
\quad
\left(
\begin{matrix}
2w \\
w \\
2w
\end{matrix}
\right),
\quad
\left(
\begin{matrix}
0 \\
0 \\
0
\end{matrix}
\right)
$$
$w$와 $m$은 임의이다.
\item
\label{models-item-four-cycle}
$n = 4$이고 $m_i, a_{ij}, w_i, g_i$는 다음과 같이 주어진다.
$$
\left(
\begin{matrix}
m \\
m \\
m \\
m
\end{matrix}
\right),
\quad
\left(
\begin{matrix}
-2w & w & 0 & w \\
w & -2w & w & 0 \\
0 & w & -2w & w \\
w & 0 & w & -2w
\end{matrix}
\right),
\quad
\left(
\begin{matrix}
w \\
w \\
w \\
w
\end{matrix}
\right),
\quad
\left(
\begin{matrix}
0 \\
0 \\
0 \\
0
\end{matrix}
\right)
$$
$w$와 $m$은 임의이다.
\item
\label{models-item-up-equal-up}
$n = 4$이고 $m_i, a_{ij}, w_i, g_i$는 다음과 같이 주어진다.
$$
\left(
\begin{matrix}
2m \\
2m \\
2m \\
m
\end{matrix}
\right),
\quad
\left(
\begin{matrix}
-2w & 2w & 0 & 0 \\
2w & -4w & 2w & 0 \\
0 & 2w & -4w & 4w \\
0 & 0 & 4w & -8w
\end{matrix}
\right),
\quad
\left(
\begin{matrix}
w \\
2w \\
2w \\
4w
\end{matrix}
\right),
\quad
\left(
\begin{matrix}
0 \\
0 \\
0 \\
0
\end{matrix}
\right)
$$
$w$와 $m$은 임의이다.
\item
\label{models-item-up-equal-down}
$n = 4$이고 $m_i, a_{ij}, w_i, g_i$는 다음과 같이 주어진다.
$$
\left(
\begin{matrix}
m \\
m \\
m \\
m
\end{matrix}
\right),
\quad
\left(
\begin{matrix}
-2w & 2w & 0 & 0 \\
2w & -4w & 2w & 0 \\
0 & 2w & -4w & 2w \\
0 & 0 & 2w & -2w
\end{matrix}
\right),
\quad
\left(
\begin{matrix}
w \\
2w \\
2w \\
w
\end{matrix}
\right),
\quad
\left(
\begin{matrix}
0 \\
0 \\
0 \\
0
\end{matrix}
\right)
$$
$w$와 $m$은 임의이다.
\item
\label{models-item-down-equal-up}
$n = 4$이고 $m_i, a_{ij}, w_i, g_i$는 다음과 같이 주어진다.
$$
\left(
\begin{matrix}
m \\
2m \\
2m \\
m
\end{matrix}
\right),
\quad
\left(
\begin{matrix}
-4w & 2w & 0 & 0 \\
2w & -2w & w & 0 \\
0 & w & -2w & 2w \\
0 & 0 & 2w & -4w
\end{matrix}
\right),
\quad
\left(
\begin{matrix}
2w \\
w \\
w \\
2w
\end{matrix}
\right),
\quad
\left(
\begin{matrix}
0 \\
0 \\
0 \\
0
\end{matrix}
\right)
$$
$w$와 $m$은 임의이다.
\item
\label{models-item-triple-with-up}
$n = 4$이고 $m_i, a_{ij}, w_i, g_i$는 다음과 같이 주어진다.
$$
\left(
\begin{matrix}
2m \\
m \\
m \\
m
\end{matrix}
\right),
\quad
\left(
\begin{matrix}
-2w & w & w & 2w \\
w & -2w & 0 & 0 \\
w & 0 & -2w & 0 \\
2w & 0 & 0 & -4w
\end{matrix}
\right),
\quad
\left(
\begin{matrix}
w \\
w \\
w \\
2w
\end{matrix}
\right),
\quad
\left(
\begin{matrix}
0 \\
0 \\
0 \\
0
\end{matrix}
\right)
$$
$w$와 $m$은 임의이다.
\item
\label{models-item-triple-with-down}
$n = 4$이고 $m_i, a_{ij}, w_i, g_i$는 다음과 같이 주어진다.
$$
\left(
\begin{matrix}
2m \\
m \\
m \\
2m
\end{matrix}
\right),
\quad
\left(
\begin{matrix}
-4w & 2w & 2w & 2w \\
2w & -4w & 0 & 0 \\
2w & 0 & -4w & 0 \\
2w & 0 & 0 & -2w
\end{matrix}
\right),
\quad
\left(
\begin{matrix}
2w \\
2w \\
2w \\
w
\end{matrix}
\right),
\quad
\left(
\begin{matrix}
0 \\
0 \\
0 \\
0
\end{matrix}
\right)
$$
$w$와 $m$은 임의이다.
\item
\label{models-item-five-cycle}
$n = 5$이고 $m_i, a_{ij}, w_i, g_i$는 다음과 같이 주어진다.
$$
\left(
\begin{matrix}
m \\
m \\
m \\
m \\
m
\end{matrix}
\right),
\quad
\left(
\begin{matrix}
-2w & w & 0 & 0 & w \\
w & -2w & w & 0 & 0 \\
0 & w & -2w & w & 0 \\
0 & 0 & w & -2w & w \\
w & 0 & 0 & w & -2w \\
\end{matrix}
\right),
\quad
\left(
\begin{matrix}
w \\
w \\
w \\
w \\
w
\end{matrix}
\right),
\quad
\left(
\begin{matrix}
0 \\
0 \\
0 \\
0 \\
0
\end{matrix}
\right)
$$
$w$와 $m$은 임의이다.
\item
\label{models-item-equal-equal-up-equal}
$n = 5$이고 $m_i, a_{ij}, w_i, g_i$는 다음과 같이 주어진다.
$$
\left(
\begin{matrix}
m \\
2m \\
3m \\
2m \\
m
\end{matrix}
\right),
\quad
\left(
\begin{matrix}
-2w & w & 0 & 0 & 0 \\
w & -2w & w & 0 & 0 \\
0 & w & -2w & 2w & 0 \\
0 & 0 & 2w & -4w & 2w \\
0 & 0 & 0 & 2w & -4w \\
\end{matrix}
\right),
\quad
\left(
\begin{matrix}
w \\
w \\
w \\
2w \\
2w
\end{matrix}
\right),
\quad
\left(
\begin{matrix}
0 \\
0 \\
0 \\
0 \\
0
\end{matrix}
\right)
$$
$w$와 $m$은 임의이다.
\item
\label{models-item-equal-equal-down-equal}
$n = 5$이고 $m_i, a_{ij}, w_i, g_i$는 다음과 같이 주어진다.
$$
\left(
\begin{matrix}
m \\
2m \\
3m \\
4m \\
2m
\end{matrix}
\right),
\quad
\left(
\begin{matrix}
-4w & 2w & 0 & 0 & 0 \\
2w & -4w & 2w & 0 & 0 \\
0 & 2w & -4w & 2w & 0 \\
0 & 0 & 2w & -2w & w \\
0 & 0 & 0 & w & -2w \\
\end{matrix}
\right),
\quad
\left(
\begin{matrix}
2w \\
2w \\
2w \\
w \\
w
\end{matrix}
\right),
\quad
\left(
\begin{matrix}
0 \\
0 \\
0 \\
0 \\
0
\end{matrix}
\right)
$$
$w$와 $m$은 임의이다.
\item
\label{models-item-up-equal-equal-up}
$n = 5$이고 $m_i, a_{ij}, w_i, g_i$는 다음과 같이 주어진다.
$$
\left(
\begin{matrix}
2m \\
2m \\
2m \\
2m \\
m
\end{matrix}
\right),
\quad
\left(
\begin{matrix}
-2w & 2w & 0 & 0 & 0 \\
2w & -4w & 2w & 0 & 0 \\
0 & 2w & -4w & 2w & 0 \\
0 & 0 & 2w & -4w & 4w \\
0 & 0 & 0 & 4w & -8w \\
\end{matrix}
\right),
\quad
\left(
\begin{matrix}
w \\
2w \\
2w \\
2w \\
4w
\end{matrix}
\right),
\quad
\left(
\begin{matrix}
0 \\
0 \\
0 \\
0 \\
0
\end{matrix}
\right)
$$
$w$와 $m$은 임의이다.
\item
\label{models-item-up-equal-equal-down}
$n = 5$이고 $m_i, a_{ij}, w_i, g_i$는 다음과 같이 주어진다.
$$
\left(
\begin{matrix}
m \\
m \\
m \\
m \\
m
\end{matrix}
\right),
\quad
\left(
\begin{matrix}
-2w & 2w & 0 & 0 & 0 \\
2w & -4w & 2w & 0 & 0 \\
0 & 2w & -4w & 2w & 0 \\
0 & 0 & 2w & -4w & 2w \\
0 & 0 & 0 & 2w & -2w \\
\end{matrix}
\right),
\quad
\left(
\begin{matrix}
w \\
2w \\
2w \\
2w \\
w
\end{matrix}
\right),
\quad
\left(
\begin{matrix}
0 \\
0 \\
0 \\
0 \\
0
\end{matrix}
\right)
$$
$w$와 $m$은 임의이다.
\item
\label{models-item-down-equal-equal-up}
$n = 5$이고 $m_i, a_{ij}, w_i, g_i$는 다음과 같이 주어진다.
$$
\left(
\begin{matrix}
m \\
2m \\
2m \\
2m \\
m
\end{matrix}
\right),
\quad
\left(
\begin{matrix}
-4w & 2w & 0 & 0 & 0 \\
2w & -2w & w & 0 & 0 \\
0 & w & -2w & w & 0 \\
0 & 0 & w & -2w & 2w \\
0 & 0 & 0 & 2w & -4w \\
\end{matrix}
\right),
\quad
\left(
\begin{matrix}
2w \\
w \\
w \\
w \\
2w
\end{matrix}
\right),
\quad
\left(
\begin{matrix}
0 \\
0 \\
0 \\
0 \\
0
\end{matrix}
\right)
$$
$w$와 $m$은 임의이다.
\item
\label{models-item-quadruple}
$n = 5$이고 $m_i, a_{ij}, w_i, g_i$는 다음과 같이 주어진다.
$$
\left(
\begin{matrix}
2m \\
m \\
m \\
m \\
m
\end{matrix}
\right),
\quad
\left(
\begin{matrix}
-2w & w & w & w & w \\
w & -2w & 0 & 0 & 0 \\
w & 0 & -2w & 0 & 0 \\
w & 0 & 0 & -2w & 0 \\
w & 0 & 0 & 0 & -2w \\
\end{matrix}
\right),
\quad
\left(
\begin{matrix}
w \\
w \\
w \\
w \\
w
\end{matrix}
\right),
\quad
\left(
\begin{matrix}
0 \\
0 \\
0 \\
0 \\
0
\end{matrix}
\right)
$$
$w$와 $m$은 임의이다.
\item
\label{models-item-triple-extended-up}
$n = 5$이고 $m_i, a_{ij}, w_i, g_i$는 다음과 같이 주어진다.
$$
\left(
\begin{matrix}
m \\
2m \\
2m \\
m \\
m
\end{matrix}
\right),
\quad
\left(
\begin{matrix}
-4w & 2w & 0 & 0 & 0 \\
2w & -2w & w & 0 & 0 \\
0 & w & -2w & w & w \\
0 & 0 & w & -2w & 0 \\
0 & 0 & w & 0 & -2w \\
\end{matrix}
\right),
\quad
\left(
\begin{matrix}
2w \\
w \\
w \\
w \\
w
\end{matrix}
\right),
\quad
\left(
\begin{matrix}
0 \\
0 \\
0 \\
0 \\
0
\end{matrix}
\right)
$$
$w$와 $m$은 임의이다.
\item
\label{models-item-triple-extended-down}
$n = 5$이고 $m_i, a_{ij}, w_i, g_i$는 다음과 같이 주어진다.
$$
\left(
\begin{matrix}
2m \\
2m \\
2m \\
m \\
m
\end{matrix}
\right),
\quad
\left(
\begin{matrix}
-2w & 2w & 0 & 0 & 0 \\
2w & -4w & 2w & 0 & 0 \\
0 & 2w & -4w & 2w & 2w \\
0 & 0 & 2w & -4w & 0 \\
0 & 0 & 2w & 0 & -4w \\
\end{matrix}
\right),
\quad
\left(
\begin{matrix}
w \\
2w \\
2w \\
2w \\
2w
\end{matrix}
\right),
\quad
\left(
\begin{matrix}
0 \\
0 \\
0 \\
0 \\
0
\end{matrix}
\right)
$$
$w$와 $m$은 임의이다.
\item
\label{models-item-n-cycle}
$n \geq 6$이고 (\ref{models-item-five-cycle})를 일반화하는 $n$-순환이 있다.
\begin{enumerate}
\item $m_1 = \ldots = m_n = m$,
\item $a_{12} = \ldots = a_{(n - 1) n} = w$, $a_{1n} = w$,
$i < j$인 나머지 경우에는 $a_{ij} = 0$이고,
\item $w_1 = \ldots = w_n = w$
\end{enumerate}
$w$와 $m$은 임의이다.
\item
\label{models-item-up-chain-equal-up}
$n \geq 6$이고 (\ref{models-item-up-equal-equal-up})을 일반화하는 사슬이 있다.
\begin{enumerate}
\item $m_1 = \ldots = m_{n - 1} = 2m$, $m_n = m$,
\item $a_{12} = \ldots = a_{(n - 2) (n - 1)} = 2w$, $a_{(n - 1) n} = 4w$,
$i < j$인 나머지 경우에는 $a_{ij} = 0$이고,
\item $w_1 = w$, $w_2 = \ldots = w_{n - 1} = 2w$, $w_n = 4w$
\end{enumerate}
$w$와 $m$은 임의이다.
\item
\label{models-item-up-chain-equal-down}
$n \geq 6$이고 (\ref{models-item-up-equal-equal-down})을 일반화하는 사슬이 있다.
\begin{enumerate}
\item $m_1 = \ldots = m_n = m$,
\item $a_{12} = \ldots = a_{(n - 1) n} = w$,
$i < j$인 나머지 경우에는 $a_{ij} = 0$이고,
\item $w_1 = w$, $w_2 = \ldots = w_{n - 1} = 2w$, $w_n = w$
\end{enumerate}
$w$와 $m$은 임의이다.
\item
\label{models-item-down-chain-equal-up}
$n \geq 6$이고 (\ref{models-item-down-equal-equal-up})을 일반화하는 사슬이 있다.
\begin{enumerate}
\item $m_1 = w$, $w_2 = \ldots = m_{n - 1} = 2m$, $m_n = m$,
\item $a_{12} = 2w$, $a_{23} = \ldots = a_{(n - 2) (n - 1)} = w$,
$a_{(n - 1) n} = 2w$이고, $i < j$인 나머지 경우에는 $a_{ij} = 0$이며,
\item $w_1 = 2w$, $w_2 = \ldots = w_{n - 1} = w$, $w_n = 2w$
\end{enumerate}
$w$와 $m$은 임의이다.
\item
\label{models-item-Dn-extended-up}
$n \geq 6$이고 (\ref{models-item-triple-extended-up})을 일반화하는 형이 있다.
\begin{enumerate}
\item $m_1 = m$, $m_2 = \ldots = m_{n - 3} = 2m$, $m_{n - 1} = m_n = m$,
\item $a_{12} = 2w$, $a_{23} = \ldots = a_{(n - 2) (n - 1)} = w$,
$a_{(n - 2) n} = w$이고, $i < j$인 나머지 경우에는 $a_{ij} = 0$이며,
\item $w_1 = 2w$, $w_2 = \ldots = w_n = w$
\end{enumerate}
$w$와 $m$은 임의이다.
\item
\label{models-item-Dn-extended-down}
$n \geq 6$이고 (\ref{models-item-triple-extended-down})을 일반화하는 형이 있다.
\begin{enumerate}
\item $m_1 = \ldots = m_{n - 3} = 2m$, $m_{n - 1} = m_n = m$,
\item $a_{12} = \ldots = a_{(n - 2) (n - 1)} = 2w$,
$a_{(n - 2) n} = 2w$이고, $i < j$인 나머지 경우에는 $a_{ij} = 0$이며,
\item $w_1 = w$, $w_2 = \ldots = w_n = 2w$
\end{enumerate}
$w$와 $m$은 임의이다.
\item
\label{models-item-double-triple}
$n \geq 6$이고 (\ref{models-item-quadruple})을 일반화하는 형이 있다.
\begin{enumerate}
\item $m_1 = m_2 = m$, $m_3 = \ldots = m_{n - 2} = 2m$, $m_{n - 1} = m_n = m$,
\item $a_{13} = w$, $a_{23} = \ldots = a_{(n - 2) (n - 1)} = w$,
$a_{(n - 2) n} = w$이고, $i < j$인 나머지 경우에는 $a_{ij} = 0$이며,
\item $w_1 = \ldots = w_n = w$,
\end{enumerate}
$w$와 $m$은 임의이다.
\item
\label{models-item-E6-completed}
$n = 7$이고 $m_i, a_{ij}, w_i, g_i$는 다음과 같이 주어진다.
$$
\left(
\begin{matrix}
m \\
2m \\
3m \\
m \\
2m \\
m \\
2m
\end{matrix}
\right),
\quad
\left(
\begin{matrix}
-2w & w & 0 & 0 & 0 & 0 & 0 \\
w & -2w & w & 0 & 0 & 0 & 0 \\
0 & w & -2w & 0 & w & 0 & w \\
0 & 0 & 0 & -2w & w & 0 & 0 \\
0 & 0 & w & w & -2w & 0 & 0 \\
0 & 0 & 0 & 0 & 0 & -2w & w \\
0 & 0 & w & 0 & 0 & w & -2w
\end{matrix}
\right),
\quad
\left(
\begin{matrix}
w \\
w \\
w \\
w \\
w \\
w \\
w
\end{matrix}
\right),
\quad
\left(
\begin{matrix}
0 \\
0 \\
0 \\
0 \\
0 \\
0 \\
0
\end{matrix}
\right)
$$
$w$와 $m$은 임의이다.
\item
\label{models-item-E7-completed}
$n = 8$이고 $m_i, a_{ij}, w_i, g_i$는 다음과 같이 주어진다.
$$
\left(
\begin{matrix}
m \\
2m \\
3m \\
4m \\
3m \\
2m \\
m \\
2m
\end{matrix}
\right),
\quad
\left(
\begin{matrix}
-2w & w & 0 & 0 & 0 & 0 & 0 & 0 \\
w & -2w & w & 0 & 0 & 0 & 0 & 0 \\
0 & w & -2w & w & 0 & 0 & 0 & 0 \\
0 & 0 & w & -2w & w & 0 & 0 & w \\
0 & 0 & 0 & w & -2w & w & 0 & 0 \\
0 & 0 & 0 & 0 & w & -2w & w & 0 \\
0 & 0 & 0 & 0 & 0 & w & -2w & 0 \\
0 & 0 & 0 & w & 0 & 0 & 0 & -2w \\
\end{matrix}
\right),
\quad
\left(
\begin{matrix}
w \\
w \\
w \\
w \\
w \\
w \\
w \\
w
\end{matrix}
\right),
\quad
\left(
\begin{matrix}
0 \\
0 \\
0 \\
0 \\
0 \\
0 \\
0 \\
0
\end{matrix}
\right)
$$
$w$와 $m$은 임의이다.
\item
\label{models-item-E8-completed}
$n = 9$이고 $m_i, a_{ij}, w_i, g_i$는 다음과 같이 주어진다.
$$
\left(
\begin{matrix}
m \\
2m \\
3m \\
4m \\
5m \\
6m \\
4m \\
2m \\
3m
\end{matrix}
\right),
\quad
\left(
\begin{matrix}
-2w & w & 0 & 0 & 0 & 0 & 0 & 0 & 0 \\
w & -2w & w & 0 & 0 & 0 & 0 & 0 & 0 \\
0 & w & -2w & w & 0 & 0 & 0 & 0 & 0 \\
0 & 0 & w & -2w & w & 0 & 0 & 0 & 0 \\
0 & 0 & 0 & w & -2w & w & 0 & 0 & 0 \\
0 & 0 & 0 & 0 & w & -2w & w & 0 & w \\
0 & 0 & 0 & 0 & 0 & w & -2w & w & 0 \\
0 & 0 & 0 & 0 & 0 & 0 & w & -2w & 0 \\
0 & 0 & 0 & 0 & 0 & w & 0 & 0 & -2w \\
\end{matrix}
\right),
\quad
\left(
\begin{matrix}
w \\
w \\
w \\
w \\
w \\
w \\
w \\
w \\
w
\end{matrix}
\right),
\quad
\left(
\begin{matrix}
0 \\
0 \\
0 \\
0 \\
0 \\
0 \\
0 \\
0 \\
0
\end{matrix}
\right)
$$
$w$와 $m$은 임의이다.
\end{enumerate}
\end{lemma}

\begin{proof}
이는 절 \ref{models-section-classify-proper-subgraphs}에서 증명하였다. 절
\ref{models-section-classify-proper-subgraphs}의 첫머리 논의를 참조하라.
\end{proof}






\section{수치형 불변량의 한계}
\label{models-section-towards-classification}

\noindent
곡선의 반안정 축약을 증명할 때 종수 $g$인 수치형의 Picard 군에 대한
한계를 사용할 것이며, 이 절에서 그 한계를 증명한다.

\begin{lemma}
\label{models-lemma-bound-neighbours}
$n, m_i, a_{ij}, w_i, g_i$를 종수 $g$인 수치형이라 하자.
$a_{ij} > 0$인 $i, j$에 대해
$m_ia_{ij} \leq m_j|a_{jj}|$이고 $m_iw_i \leq m_j|a_{jj}|$이다.
\end{lemma}

\begin{proof}
모든 지표 $j$에 대해 $m_j a_{jj} + \sum_{i \not = j} m_ia_{ij} = 0$이다.
따라서 $|a_{jj}|$와 $m_j$의 상계가 있으면, 영이 아닌(따라서 양의)
$a_{ij}$ 및 $m_i$의 상계도 얻는다. $w_i$가 $a_{ij}$를 나눈다는 사실을
상기하면 보조정리를 쉽게 확인할 수 있다.
\end{proof}

\begin{lemma}
\label{models-lemma-bound-heart}
$g \geq 2$를 고정하자. $n > 1$이고 종수 $g$인 모든 극소 수치형
$n, m_i, a_{ij}, w_i, g_i$에 대해 다음이 성립한다.
\begin{enumerate}
\item $(-2)$-지표가 아닌 지표들의 집합 $J \subset \{1, \ldots, n\}$의
원소 수는 $2g - 2$ 이하이다.
\item $j \in J$에 대해 $g_j < g$이다.
\item $j \in J$에 대해 $m_j|a_{jj}| \leq 6g - 6$이다.
\item $j \in J$, $i \in \{1, \ldots, n\}$에 대해
$m_ia_{ij} \leq 6g - 6$이다.
\end{enumerate}
\end{lemma}

\begin{proof}
$g = 1 + \sum m_j(w_j(g_j - 1) - \frac{1}{2} a_{jj})$임을 상기하자.
$j \in J$가 종수 $g$에 기여하는 항
$m_j(w_j(g_j - 1) - \frac{1}{2} a_{jj})$는 $> 0$이고 따라서
$\geq 1/2$이다. 여기서는 보조정리 \ref{models-lemma-minus-one}, 정의
\ref{models-definition-type-minus-one}, 정의 \ref{models-definition-type-minimal},
보조정리 \ref{models-lemma-minus-two}, 정의 \ref{models-definition-type-minus-two}를
사용하였다. 아래에서는 이 결과들을 별도 언급 없이 사용한다. 따라서 $J$의
원소 수는 $2(g - 1)$ 이하이다. 이것으로 (1)이 증명된다.

\medskip\noindent
보조정리 \ref{models-lemma-diagonal-negative}에 의해 모든 $i$에 대해
$-a_{ii} > 0$임을 상기하자. 따라서 $j \in J$가 종수 $g$에 기여하는
$m_j(w_j(g_j - 1) - \frac{1}{2} a_{jj})$는
$> m_jw_j(g_j - 1)$이다. 그러므로
$$
g - 1 > m_jw_j(g_j - 1) \Rightarrow g_j < (g - 1)/m_jw_j + 1
$$
이고, 이는 실제로 $g_j < g$를 함의하여 (2)를 증명한다.

\medskip\noindent
$j \in J$에 대해 $g_j > 0$이면 종수 $g$에 기여하는 항
$m_j(w_j(g_j - 1) - \frac{1}{2} a_{jj})$는
$\geq -\frac{1}{2}m_ja_{jj}$이므로 즉시
$m_j|a_{jj}| \leq 2(g - 1)$을 얻는다. 그렇지 않으면 $j \in J$이므로
어떤 정수 $k \geq 3$에 대해 $a_{jj} = -kw_j$이고
$$
m_jw_j(-1 + \frac{k}{2}) \leq g - 1
\Rightarrow
m_jw_j \leq \frac{2(g - 1)}{k - 2}
$$
이를 $a_{jj} = -km_jw_j$에 다시 대입하면
$$
m_j|a_{jj}| \leq 2(g - 1) \frac{k}{k - 2} \leq 6(g - 1)
$$
을 얻는다. 이것으로 (3)이 증명된다.

\medskip\noindent
(4)는 보조정리 \ref{models-lemma-bound-neighbours}과 (3)에서 따른다.
\end{proof}

\begin{lemma}
\label{models-lemma-bound-wm}
$g \geq 2$를 고정하자. 종수 $g$인 모든 극소 수치형
$n, m_i, a_{ij}, w_i, g_i$에 대해 $m_i|a_{ij}| \leq 768g$이다.
\end{lemma}

\begin{proof}
보조정리 \ref{models-lemma-bound-neighbours}에 의해 모든 $i$에 대해
$m_i|a_{ii}| \leq 768g$임을 보이면 충분하다. $J \subset \{1, \ldots, n\}$를
보조정리 \ref{models-lemma-bound-heart}에서와 같이 $(-2)$-지표가 아닌 지표들의
집합이라 하자. $g \geq 2$이므로 $J$는 공집합이 아니다. 또한 그
보조정리에 의해 $j \in J$에 대해 $m_j|a_{jj}| \leq 6g$이다.

\medskip\noindent
$j \in J$와 $(-2)$-지표의 수열 $i_1, \ldots, i_7$이 있어
$a_{ji_1}$과 $a_{i_1i_2}$, $a_{i_2i_3}$, $a_{i_3i_4}$,
$a_{i_4i_5}$, $a_{i_5i_6}$, $a_{i_6i_7}$이 영이 아니라고 하자.
보조정리 \ref{models-lemma-bound-neighbours}에 의해
$m_{i_1}w_{i_1} \leq 6g$이고 $m_{i_1}a_{ji_1} \leq 6g$이다.
$i_1$이 $(-2)$-지표이므로 $a_{i_1i_1} = -2w_{i_1}$이고 따라서
$m_{i_1}|a_{i_1i_1}| \leq 12g$이다. 논증을 반복하면
$m_{i_2}w_{i_2} \leq 12g$이고 $m_{i_2}a_{i_1i_2} \leq 12g$이며,
따라서 $m_{i_2}|a_{i_2i_2}| \leq 24g$인 식으로 계속된다. 마침내
$k = 1, \ldots, 7$에 대해
$m_{i_k}|a_{i_ki_k}| \leq 2^k(6g) \leq 768g$를 얻는다.

\medskip\noindent
$I \subset \{1, \ldots, n\} \setminus J$를 극대 연결 부분집합이라 하자.
즉 $i' \in I'$, $i \in I \setminus I'$일 때 $a_{i'i} = 0$이 되는
공집합이 아닌 진부분집합 $I' \subset I$가 없고, $I$는 이 성질에 관해
극대이다. 특히 수치형은 정의상 연결되어 있으므로 $a_{ij} > 0$인
$j \in J$와 $i \in I$가 존재한다. 명제
\ref{models-proposition-classify-subgraphs}의 이런 $I$에 대한 분류와 앞 문단의
결과를 사용하면, $I$가 보조정리 \ref{models-lemma-long} 또는 보조정리
\ref{models-lemma-Dn}에 기술된 경우가 아닌 한 모든 $i \in I$에 대해
$w_i|a_{ii}| \leq 768g$이다. 따라서 $i, i' \in I$에 대한 $a_{ii'}$의
영 아닌 무늬가 다음 형태
$$
\xymatrix{
\bullet \ar@{-}[r] &
\bullet \ar@{-}[r] &
\bullet \ar@{..}[r] &
\bullet \ar@{-}[r] &
\bullet \ar@{-}[r] &
\bullet
}
$$
(보조정리 \ref{models-lemma-long}에 상세히 적은 세 하위 경우가 있다)이거나
다음 형태라고 가정할 수 있다.
$$
\xymatrix{
\bullet \ar@{-}[r] &
\bullet \ar@{-}[r] &
\bullet \ar@{..}[r] &
\bullet \ar@{-}[r] &
\bullet \ar@{-}[r] \ar@{-}[d] &
\bullet \\
& & & & \bullet
}
$$
보조정리 \ref{models-lemma-long}의 첫째 하위 경우에 한계가 성립함을 증명하고,
나머지 경우는 독자에게 맡긴다. 그 경우들의 논증도 거의 완전히 같다.

\medskip\noindent
번호를 다시 붙여 $I = \{1, \ldots, t\} \subset \{1, \ldots, n\}$라
가정할 수 있고, 어떤 정수 $w$가 있어
$$
w = w_1 = \ldots = w_t =
a_{12} = \ldots = a_{(t - 1) t} =
-\frac{1}{2} a_{i_1i_2} = \ldots = -\frac{1}{2} a_{(t - 1) t}
$$
이다. 등식 $a_{ii}m_i + \sum_{j \not = i} a_{ij}m_j = 0$은
$$
2m_2 \geq m_1 + m_3, \ldots, 2m_{t - 1} \geq m_{t - 2} + m_t
$$
을 함의한다. $2m_i \geq m_{i - 1} + m_{i + 1}$에서 등호가 성립할
필요충분조건은 $i$가 $i - 1$과 $i + 1$ 이외의 지표와 ``만나지'' 않는
것이다. $i$가 다른 지표와 만나면 $I$의 극대성에 의해 그 지표는 $J$에
속한다. 특히 사상 $\{1, \ldots, t\} \to \mathbf{Z}$,
$i \mapsto m_i$는 오목하다.

\medskip\noindent
$m = \max(m_i, i \in \{1, \ldots, t\})$라 하자. 그러면 $i \leq t$에
대해 $m_i|a_{ii}| \leq 2mw$이고, 목표는 $2mw \leq 768g$를 보이는 것이다.
$\{1, \ldots, t\}$에서 $m_s = m = m_{s'}$인 가장 작은 지표와 가장 큰
지표를 각각 $s$, $s'$라 하자. 오목성에 의해 $s \leq i \leq s'$이면
$m_i = m$이다. $s > 1$이면 $2m_s \geq m_{s - 1} + m_{s + 1}$에서
등호가 성립하지 않으므로 $s$는 $J$의 지표 하나와 만난다. 이 경우 증명의
둘째 문단 결과에 의해 $2mw \leq 12g$이다. 마찬가지로 $s' < t$이면
$s'$가 $J$의 지표와 만나고 역시 $2mw \leq 12g$이다. 한편 $s = 1$이고
$s' = t$이면 모든 $j \in J$와 $i \in \{2, \ldots, t - 1\}$에 대해
$a_{ij} = 0$이다. 그러나 $a_{ij} > 0$인 쌍 $(i, j) \in I \times J$가
반드시 있음을 보았으므로, 이는 $i = 1$ 또는 $i = t$일 때 일어난다.
이 경우 $m_1 = m = m_t$이므로 앞과 같은 방식으로 $2mw \leq 12g$를 얻는다.
\end{proof}

\begin{proposition}
\label{models-proposition-bound-picard-group}
$g \geq 2$라 하자. 종수 $g$인 모든 수치형 $T$와 소수
$\ell > 768g$에 대해
$$
\dim_{\mathbf{F}_\ell} \Pic(T)[\ell] \leq g
$$
이다. 여기서 $\Pic(T)$는 정의 \ref{models-definition-picard-group}와 같다.
$T$가 극소이면 더 나아가
$$
\dim_{\mathbf{F}_\ell} \Pic(T)[\ell] \leq g_{top} \leq g
$$
이다. 여기서 $g_{top}$은 정의 \ref{models-definition-top-genus}와 같다.
\end{proposition}

\begin{proof}
$T$가 $n, m_i, a_{ij}, w_i, g_i$로 주어졌다고 하자. $T$가 극소가
아니면 $(-1)$-지표가 존재한다. $T$를 동치인 형으로 바꾸어 $n$이
$(-1)$-지표라고 가정할 수 있다. 보조정리
\ref{models-lemma-contract-picard-group}을 적용하면 $\Pic(T) \subset \Pic(T')$이고,
$T'$은 $n - 1$개의 지표를 갖는 종수 $g$의 수치형이다(보조정리
\ref{models-lemma-contract}). 따라서 극소 수치형에 대해 보조정리를 증명하면
$n$에 대한 귀납법으로 결론을 얻는다.

\medskip\noindent
$T$가 종수 $\geq 2$인 극소 수치형이라고 가정하자. 보조정리
\ref{models-lemma-genus-nonnegative}에 의해 $g_{top} \leq g$이다.
$A = (a_{ij})$라 하면 보조정리 \ref{models-lemma-picard-T-and-A}에 의해
$\Pic(T) \subset \Coker(A)$이므로 $\Coker(A)$에 대해 증명하면 충분하다.
보조정리 \ref{models-lemma-bound-wm}에 의해 모든 $i, j$에 대해
$m_i|a_{ij}| \leq 768g$이다. 이제 보조정리
\ref{models-lemma-recurring-symmetric-integer}에서 결과가 따른다.
\end{proof}







\section{모형}
\label{models-section-models}

\noindent
이 장에서 $R$은 이산 값매김환이고 $K$는 그 분수체이다. 필요하면
균등화 원소를 $\pi \in R$, 잉여체를 $k = R/(\pi)$로 나타낸다.

\medskip\noindent
$V$를 대수적 $K$-스킴이라 하자(Varieties, 정의
\ref{varieties-definition-algebraic-scheme}). $V$의 {\it 모형}이란 동형
$V \to X_K = X \times_{\Spec(R)} \Spec(K)$이 주어진 평탄 유한형
사상\footnote{때로는 $R$ 위에서 국소 유한형인 모형을 허용하는 것이
유용하지만, 필요해질 때 다룬다.} $X \to \Spec(R)$을 뜻한다. 흔히 $V$를
$X$의 일반 올 $X_K$와 동일시하여 단순히 $V = X_K$라 쓴다. 특수 올은
$X_k = X \times_{\Spec(R)} \Spec(k)$이다. $V$의 {\it 모형의 사상
$X \to X'$}이란 $V$ 위에 항등사상을 유도하는 $R$ 위 스킴의 사상
$X \to X'$이다.

\medskip\noindent
$X$가 $V$의 모형이고 구조 사상 $X \to \Spec(R)$이 고유이면
{\it $X$는 $V$의 고유 모형}이라고 한다. 분리 모형, 매끄러운 모형 등도
마찬가지이다. $X$가 $V$의 모형이고 $X$가 정칙 스킴이면
{\it $X$는 $V$의 정칙 모형}이라고 한다. 정규 모형, 기약 모형 등도
마찬가지이다.

\medskip\noindent
$R \subset R'$을 이산 값매김환의 확대라 하자(More on Algebra, 정의
\ref{more-algebra-definition-extension-discrete-valuation-rings}). 이는
분수체의 확대 $K'/K$를 유도한다. $K$ 위의 대수적 스킴 $V$가 주어졌을 때
밑변환 $V \times_{\Spec(K)} \Spec(K')$을 $V'$라 쓰자. 그러면 함자
$$
V\text{의 }R\text{ 위 모형}
\longrightarrow
V'\text{의 }R'\text{ 위 모형}
$$
가 있으며 $X$를 $X \times_{\Spec(R)} \Spec(R')$로 보낸다.

\begin{lemma}
\label{models-lemma-closure-is-model}
$V_1 \to V_2$를 $K$ 위 대수적 스킴의 닫힌 몰입이라 하자. $X_2$가
$V_2$의 모형이면 $V_1 \to X_2$의 스킴론적 상은 $V_1$의 모형이다.
\end{lemma}

\begin{proof}
Morphisms, 보조정리
\ref{morphisms-lemma-quasi-compact-scheme-theoretic-image}와 예
\ref{morphisms-example-scheme-theoretic-image}을 사용하면 다음 대수 명제로
귀착된다. $A_1$을 $R$ 위에 평탄한 유한형 $R$-대수라 하고
$A_1 \otimes_R K \to B_2$를 전사라 하자. 그러면
$A_2 = A_1 / \Ker(A_1 \to B_2)$는 $R$ 위에 평탄한 유한형 $R$-대수이고
$B_2 = A_2 \otimes_R K$이다. 자세한 증명은 생략한다. $A_2$의 평탄성에는
More on Algebra, 보조정리
\ref{more-algebra-lemma-dedekind-torsion-free-flat}을 사용하라.
\end{proof}

\begin{lemma}
\label{models-lemma-normalization}
$X$를 $K$ 위의 기하적으로 정규인 다양체 $V$의 모형이라 하자. 그러면
정규화 $\nu : X^\nu \to X$는 유한하고, $X^\nu$를 완비화 $R^\wedge$로
밑변환한 것은 $X$의 밑변환의 정규화이다. 더욱이 각 $x \in X^\nu$에 대해
$\mathcal{O}_{X^\nu, x}$의 완비화는 정규이다.
\end{lemma}

\begin{proof}
$R^\wedge$는 이산 값매김환이다(More on Algebra, 보조정리
\ref{more-algebra-lemma-completion-dvr}).
$Y = X \times_{\Spec(R)} \Spec(R^\wedge)$라 두자. $R^\wedge$가 이산
값매김환이므로
$$
Y \setminus Y_k =
Y \times_{\Spec(R^\wedge)} \Spec(K^\wedge) =
V \times_{\Spec(K)} \Spec(K^\wedge)
$$
이다. 여기서 $K^\wedge$는 $R^\wedge$의 분수체이다. $V$가 기하적으로
정규이므로 이는 정규 스킴이다. 따라서 보조정리의 첫째 부분은 Resolution
of Surfaces, 보조정리 \ref{resolve-lemma-normalization-completion}에서 따른다.

\medskip\noindent
둘째 부분을 증명할 때 첫째 부분에 의해 $X$와 $Y$가 정규라고 가정할 수
있다. $x$가 일반 올에 있으면 $\mathcal{O}_{X, x} = \mathcal{O}_{V, x}$는
체 위의 본질적 유한형 정규 국소환이다. 이런 환은 우수하다(More on
Algebra, 명제 \ref{more-algebra-proposition-ubiquity-excellent}).
$x$가 특수 올의 점이고 그 상이 $y \in Y$이면 Resolution of Surfaces,
보조정리 \ref{resolve-lemma-iso-completions}에 의해
$\mathcal{O}_{X, x}^\wedge = \mathcal{O}_{Y, y}^\wedge$이다. 이 경우
$R^\wedge$가 우수하므로 앞과 같은 문헌에 의해 $\mathcal{O}_{Y, y}$는
우수한 정규 국소 정역이다. $B$가 우수한 국소 정규 정역이면 완비화
$B^\wedge$는 정규이다. 실제로 $B \to B^\wedge$가 정칙이고 More on
Algebra, 보조정리 \ref{more-algebra-lemma-normal-goes-up}을 적용할 수 있다.
이것으로 증명이 끝난다.
\end{proof}

\begin{lemma}
\label{models-lemma-regular}
$X$를 $K$ 위 매끄러운 곡선 $C$의 모형이라 하자. 그러면 $X$의 특이점
해소가 존재하며, 모든 해소는 $C$의 모형이다.
\end{lemma}

\begin{proof}
Lipman 정리(Resolution of Surfaces, 정리 \ref{resolve-theorem-resolve})의
조건 (4)를 확인한다. $X^\nu$의 특이점이 유한 개라는 주장 외에는
보조정리 \ref{models-lemma-normalization}에서 분명하다. More on Algebra, 명제
\ref{more-algebra-proposition-ubiquity-J-2}에 의해 $R$이 J-2이므로
비특이 자취는 $X^\nu$에서 열린집합이다. $X^\nu$가 차원 $\leq 2$인
정규 스킴이므로 특이점들은 닫힌점이고, 특이 자취가 닫혔다는 사실과 $X$의
준콤팩트성으로부터 그 수가 유한함이 따른다. $X$의 모든 해소는 $X$의
변형이다(Resolution of Surfaces, 정의 \ref{resolve-definition-resolution}).
Varieties, 보조정리
\ref{varieties-lemma-modification-normal-iso-over-codimension-1}에 의해 이는
$X$의 정규 자취 위에서 동형이다. 정규점들의 집합이 $C = X_K$를
포함하므로 모든 해소는 $C$의 모형이다.
\end{proof}

\begin{definition}
\label{models-definition-minimal-model}
$C$를 $H^0(C, \mathcal{O}_C) = K$인 $K$ 위 매끄러운 사영곡선이라 하자.
{\it 극소 모형}이란 $C$의 정칙 고유 모형 $X$로서 $X$가 제1종 예외곡선을
포함하지 않는 것을 말한다(Resolution of Surfaces, 절
\ref{resolve-section-minus-one}).
\end{definition}

\noindent
엄밀히는 이를 극소 정칙 고유 모형, 또는 상대적 극소 정칙 사영 모형이라
불러야 한다. 그러나 이 장에서처럼 이산 값매김환 위의 모형만 다루는 한
혼동은 생기지 않을 것이다.

\medskip\noindent
극소 모형은 항상 존재하고(명제 \ref{models-proposition-exists-minimal-model}),
종수가 $> 0$이면 유일하다(보조정리 \ref{models-lemma-minimal-model-unique}).

\begin{lemma}
\label{models-lemma-pre-exists-minimal-model}
\begin{slogan}
곡선의 정칙 고유 모형은 극소 모형에서 연속적인 부풀리기로 얻어진다.
\end{slogan}
$C$를 $H^0(C, \mathcal{O}_C) = K$인 $K$ 위 매끄러운 사영곡선이라 하자.
$X$가 $C$의 정칙 고유 모형이면 사상들의 수열
$$
X = X_m \to X_{m - 1} \to \ldots \to X_1 \to X_0
$$
이 존재한다. 여기서 각 항은 $C$의 정칙 고유 모형이고, 각 사상은 제1종
예외곡선의 수축이며, $X_0$은 극소 모형이다.
\end{lemma}

\begin{proof}
Resolution of Surfaces, 보조정리
\ref{resolve-lemma-regular-dim-2-projective}에 의해 $X$는 $R$ 위에서
사영이다. 따라서 More on Morphisms, 보조정리
\ref{more-morphisms-lemma-projective}에 의해 $X$에는 풍부한 가역층이 있다.
이 사실은 아래에서 사용한다. $E \subset X$를 제1종 예외곡선이라 하자.
Resolution of Surfaces, 절 \ref{resolve-section-minus-one}을 참조하라.
Resolution of Surfaces, 보조정리 \ref{resolve-lemma-contract-ample}에 의해
$X'$가 정칙이고 $R$ 위에서 사영이 되도록 하는 사상 $X \to X'$로 $E$를
수축할 수 있다. $X'_k$의 기약 성분 수는 분명 $X_k$의 기약 성분 수보다
정확히 하나 작다. 따라서 극소 모형을 얻을 때까지 이런 수축을 유한 번만
수행할 수 있다.
\end{proof}

\begin{proposition}
\label{models-proposition-exists-minimal-model}
$C$를 $H^0(C, \mathcal{O}_C) = K$인 $K$ 위 매끄러운 사영곡선이라 하자.
극소 모형이 존재한다.
\end{proposition}

\begin{proof}
닫힌 몰입 $C \to \mathbf{P}^n_K$를 택하고 $X$를
$C \to \mathbf{P}^n_R$의 스킴론적 상이라 하자. 보조정리
\ref{models-lemma-closure-is-model}에 의해 $X \to \Spec(R)$은 $C$의 사영 모형이다.
보조정리 \ref{models-lemma-regular}에 의해 특이점 해소 $X' \to X$가 존재하고
$X'$은 $C$의 모형이다. $X' \to \Spec(R)$은 고유 사상의 합성이므로
고유하다. 이제 보조정리 \ref{models-lemma-pre-exists-minimal-model}을 적용하여
극소 모형을 얻는다.
\end{proof}





\section{정칙 모형의 기하}
\label{models-section-special-fibre}

\noindent
이 절에서는 $H^0(C, \mathcal{O}_C) = K$인 $K$ 위 매끄러운 사영곡선
$C$의 정칙 고유 모형 $X$의 기하를 기술한다.

\begin{lemma}
\label{models-lemma-divisor-special-fiber}
$X$를 $K$ 위 매끄러운 곡선 $C$의 정칙 모형이라 하자.
\begin{enumerate}
\item 특수 올 $X_k$는 $X$ 위의 유효 Cartier 인자이다.
\item $X_k$의 각 기약 성분 $C_i$는 $X$ 위의 유효 Cartier 인자이다.
\item $X_k = \sum m_i C_i$이고(유효 Cartier 인자의 합), 여기서 $m_i$는
$X_k$에서 $C_i$의 중복도이다.
\item $\mathcal{O}_X(X_k) \cong \mathcal{O}_X$.
\end{enumerate}
\end{lemma}

\begin{proof}
$R$은 균등화 원소가 $\pi$이고 잉여체가 $k = R/(\pi)$인 이산 값매김환이다.
$X \to \Spec(R)$이 평탄하므로 원소 $\pi$는 $X$ 위에서 아핀 국소적으로
영인자가 아니다(More on Algebra, 보조정리
\ref{more-algebra-lemma-dedekind-torsion-free-flat} 참조). 따라서
$U = \Spec(A) \subset X$가 아핀 열린집합이면
$$
X_k \cap U = U_k = \Spec(A \otimes_R k) = \Spec(A/\pi A)
$$
이고 $\pi$는 $A$에서 영인자가 아니다. 따라서 Divisors, 보조정리
\ref{divisors-lemma-characterize-effective-Cartier-divisor}에 의해
$X_k = V(\pi)$는 유효 Cartier 인자이다. 그러므로 (1)이 성립한다.

\medskip\noindent
위 논의에서 쌍 $(\mathcal{O}_X(X_k), 1)$은 $(\mathcal{O}_X, \pi)$와
동형이며, 이것으로 (4)가 증명된다.

\medskip\noindent
Divisors, 보조정리
\ref{divisors-lemma-effective-Cartier-divisor-is-a-sum}에 의해 서로 다른
정수적 유효 Cartier 인자 $D_i \subset X$와 정수 $a_i \geq 0$가 존재하여
$X_k = \sum a_i D_i$이다. $a_i = 0$인 인자 $D_i$는 버릴 수 있다.
그러면 유효 Cartier 인자의 합의 정의에서 집합론적으로
$X_k = \bigcup D_i$임이 분명하다. 따라서 $C_i = D_i$는 $X_k$의 기약
성분이며 (2)가 증명된다. $\xi_i$를 $C_i$의 일반점이라 하자. 그러면
$\mathcal{O}_{X, \xi_i}$는 이산 값매김환이다(Divisors, 보조정리
\ref{divisors-lemma-integral-effective-Cartier-divisor-dvr}). 균등화 원소
$\pi_i \in \mathcal{O}_{X, \xi_i}$는 $C_i$의 국소방정식이고 $\pi$의 상은
$X_k$의 국소방정식이다. $X_k = \sum a_i C_i$이므로 $\pi$와
$\pi_i^{a_i}$는 $\mathcal{O}_{X, \xi_i}$에서 같은 아이디얼을 생성한다.
한편 $X_k$에서 $C_i$의 중복도는
$$
m_i = \text{length}_{\mathcal{O}_{C_i, \xi_i}} \mathcal{O}_{X_k, \xi_i} =
\text{length}_{\mathcal{O}_{C_i, \xi_i}} \mathcal{O}_{X, \xi_i}/(\pi) =
\text{length}_{\mathcal{O}_{C_i, \xi_i}} \mathcal{O}_{X, \xi_i}/(\pi_i^{a_i}) =
a_i
$$
이다. Chow Homology, 정의
\ref{chow-definition-cycle-associated-to-closed-subscheme}를 참조하라.
따라서 $a_i = m_i$이고 (3)이 증명된다.
\end{proof}

\begin{lemma}
\label{models-lemma-gorenstein}
$X$를 $K$ 위 매끄러운 곡선 $C$의 정칙 모형이라 하자. 그러면
\begin{enumerate}
\item $X \to \Spec(R)$은 상대차원 $1$의 Gorenstein 사상이고,
\item $X_k$의 각 기약 성분 $C_i$는 Gorenstein이다.
\end{enumerate}
\end{lemma}

\begin{proof}
$X \to \Spec(R)$이 평탄하므로 (1)을 증명하려면 올들이 Gorenstein임을
보이면 충분하다(Duality for Schemes, 보조정리
\ref{duality-lemma-gorenstein-morphism}). 일반 올은 매끄러운 곡선이므로
정칙이고 따라서 Gorenstein이다(Duality for Schemes, 보조정리
\ref{duality-lemma-regular-gorenstein}). 특수 올 $X_k$는 정칙인, 따라서
Gorenstein인 스킴 위의 유효 Cartier 인자이므로 Gorenstein이다. 예를 들어
Dualizing Complexes, 보조정리
\ref{dualizing-lemma-gorenstein-divide-by-nonzerodivisor}를 사용한다.
같은 논증으로 곡선 $C_i$들도 Gorenstein이다.
\end{proof}

\begin{situation}
\label{models-situation-regular-model}
$R$을 분수체가 $K$, 잉여체가 $k$, 균등화 원소가 $\pi$인 이산 값매김환이라
하자. $C$를 $H^0(C, \mathcal{O}_C) = K$인 $K$ 위 매끄러운 사영곡선이라
하고, $X$를 $C$의 정칙 고유 모형이라 하자. $C_1, \ldots, C_n$을 특수 올
$X_k$의 기약 성분이라 하자. 보조정리 \ref{models-lemma-divisor-special-fiber}처럼
$X_k = \sum m_i C_i$로 쓰자.
\end{situation}

\begin{lemma}
\label{models-lemma-regular-model-connected}
상황 \ref{models-situation-regular-model}에서 특수 올 $X_k$는 연결이다.
\end{lemma}

\begin{proof}
More on Morphisms, 보조정리
\ref{more-morphisms-lemma-geometrically-connected-fibres-towards-normal}의
귀결이다.
\end{proof}

\begin{lemma}
\label{models-lemma-regular-model-pic}
상황 \ref{models-situation-regular-model}에서 완전열
$$
0 \to \mathbf{Z} \to \mathbf{Z}^{\oplus n} \to
\Pic(X) \to \Pic(C) \to 0
$$
이 있다. 첫째 사상은 $1$을 $(m_1, \ldots, m_n)$으로 보내고, 둘째 사상은
$i$번째 기저 벡터를 $\mathcal{O}_X(C_i)$로 보낸다.
\end{lemma}

\begin{proof}
$C \subset X$는 열린 부분스킴이다. 제한 사상 $\Pic(X) \to \Pic(C)$은
Divisors, 보조정리 \ref{divisors-lemma-extend-invertible-module}에 의해
전사이다. 동형 $s : \mathcal{O}_C \to \mathcal{L}|_C$가 존재하는 가역
$\mathcal{O}_X$-가군 $\mathcal{L}$을 택하자. 그러면 $s$는 $\mathcal{L}$의
정칙 유리형 단면이고, 어떤 $a_i \in \mathbf{Z}$에 대해
$\text{div}_\mathcal{L}(s) = \sum a_i C_i$이다(Divisors, 정의
\ref{divisors-definition-divisor-invertible-sheaf}). Divisors, 보조정리
\ref{divisors-lemma-normal-c1-injective}과 $X$의 정규성에 의해
$\mathcal{L} = \mathcal{O}_X(\sum a_iC_i)$이다. 마지막으로
$\mathcal{O}_X(\sum a_i C_i) \cong \mathcal{O}_X$라고 하자. 그러면 $X$의
함수체에 $\text{div}_X(g) = \sum a_i C_i$인 원소 $g$가 존재한다. 특히
유리함수 $g$는 $X$의 일반 올 $C$에서 영점도 극도 갖지 않는다. $C$가
정규 스킴이므로 $g \in H^0(C, \mathcal{O}_C) = K$이다. 따라서 어떤
$a \in \mathbf{Z}$와 $u \in R^*$에 대해 $g = \pi^a u$이다. 그러므로
$\text{div}_X(g) = a \sum m_i C_i$이고 증명이 끝난다.
\end{proof}

\noindent
상황 \ref{models-situation-regular-model}에서 모든 가역 $\mathcal{O}_X$-가군
$\mathcal{L}$과 모든 $i$에 대해 정수
$$
\deg(\mathcal{L}|_{C_i}) =
\chi(C_i, \mathcal{L}|_{C_i}) - \chi(C_i, \mathcal{O}_{C_i})
$$
를 얻는다. 이는 Varieties, 절 \ref{varieties-section-divisors-curves}에서처럼
$\mathcal{L}$을 $C_i$에 제한한 것의 바탕체 $k$에 대한 차수이다.\footnote{체
$\kappa_i = H^0(C_i, \mathcal{O}_{C_i})$가 $k$보다 진정으로 클 수 있다.
이 경우 $C_i$ 위의 모든 가역 가군은 위에서 정의한 차수가
$[\kappa_i : k]$로 나누어진다.}

\begin{lemma}
\label{models-lemma-intersection-pairing}
상황 \ref{models-situation-regular-model}에서 가역 $\mathcal{O}_X$-가군
$\mathcal{L}$과 $a = (a_1, \ldots, a_n) \in \mathbf{Z}^{\oplus n}$에 대해
$$
\langle a, \mathcal{L} \rangle = \sum a_i\deg(\mathcal{L}|_{C_i})
$$
로 정의한다. 그러면 $\langle , \rangle$는 쌍선형이고
$b = (b_1, \ldots, b_n) \in \mathbf{Z}^{\oplus n}$에 대해
$$
\left\langle a, \mathcal{O}_X(\sum b_i C_i) \right\rangle =
\left\langle b, \mathcal{O}_X(\sum a_i C_i) \right\rangle
$$
\end{lemma}

\begin{proof}
쌍선형성은 정의와 Varieties, 보조정리
\ref{varieties-lemma-degree-tensor-product}에서 즉시 따른다. 대칭성을
증명하려면 $a$와 $b$가 $\mathbf{Z}^{\oplus n}$의 표준기저 벡터라고
가정해도 충분하다. 따라서 모든 $1 \leq i, j \leq n$에 대해
$$
\deg(\mathcal{O}_X(C_j)|_{C_i}) = \deg(\mathcal{O}_X(C_i)|_{C_j})
$$
임을 보이면 된다. $i = j$이면 증명할 것이 없다. $i \not = j$이면
$\mathcal{O}_X(C_j)$의 표준 단면 $1$은 $\mathcal{O}_X(C_j)|_{C_i}$의
영이 아닌, 따라서 정칙인 단면으로 제한되고, 그 영점 스킴은 정확히
$C_i \cap C_j$이다(스킴론적 교차). 다시 말해 $C_i \cap C_j$는 $C_i$
위의 유효 Cartier 인자이고
$$
\deg(\mathcal{O}_X(C_j)|_{C_i}) = \deg(C_i \cap C_j)
$$
이다(Varieties, 보조정리
\ref{varieties-lemma-degree-effective-Cartier-divisor}). 대칭성으로 반대쪽에
대해서도 같은 식을 얻으며 증명이 끝난다.
\end{proof}

\noindent
상황 \ref{models-situation-regular-model}에서 $\mathbf{Z}^{\oplus n}$을 집합
$\{C_1, \ldots, C_n\}$ 위의 자유 아벨군으로 생각하면 편리할 때가 많다.
이 군의 원소를 $\sum a_i C_i$로 나타낸다. 이를 형식적 합으로 생각하지만,
동치이게 $X_k$ 위에 지지되는 $X$의 Weil 인자로 생각해도 되고 실제로
그렇게 할 때도 있다. 이제 보조정리 \ref{models-lemma-intersection-pairing}에 의해
이 자유 아벨군 위의 대칭 쌍선형 형식 $(\ \cdot\ )$을 다음 규칙으로 정의한다.
\begin{equation}
\label{models-equation-form}
\left(\sum a_i C_i \cdot \sum b_j C_j\right) =
\left\langle a, \mathcal{O}_X(\sum b_j C_j) \right\rangle =
\left\langle b, \mathcal{O}_X(\sum a_i C_i) \right\rangle
\end{equation}
이 쌍선형 형식의 몇 가지 성질을 증명한다.

\begin{lemma}
\label{models-lemma-properties-form}
상황 \ref{models-situation-regular-model}에서 대칭 쌍선형 형식
(\ref{models-equation-form})은 다음 성질을 갖는다.
\begin{enumerate}
\item $i \not = j$이면 $(C_i \cdot C_j) \geq 0$이고, 등호가 성립할
필요충분조건은 $C_i \cap C_j = \emptyset$이다.
\item $(\sum m_i C_i \cdot C_j) = 0$,
\item $i \in I$, $j \not \in I$일 때 $(C_i \cdot C_j) = 0$이 되는
공집합이 아닌 진부분집합 $I \subset \{1, \ldots, n\}$는 존재하지 않는다.
\item $(\sum a_i C_i \cdot \sum a_i C_i) \leq 0$이고, 등호가 성립할
필요충분조건은 $i = 1, \ldots, n$에 대해 $a_i = qm_i$인
$q \in \mathbf{Q}$가 존재하는 것이다.
\end{enumerate}
\end{lemma}

\begin{proof}
보조정리 \ref{models-lemma-intersection-pairing}의 증명에서 $i \not = j$이면
$(C_i \cdot C_j) = \deg(C_i \cap C_j)$임을 보았다. 이는 $\geq 0$이고,
$> 0 $일 필요충분조건은 $C_i \cap C_j \not = \emptyset$이다. 이로써
(1)이 증명된다.

\medskip\noindent
(2)의 증명. 보조정리 \ref{models-lemma-divisor-special-fiber}에 의해
$\sum m_i C_i$에 결부된 가역층은 자명하고, 자명한 층의 차수는 영이다.

\medskip\noindent
(3)의 증명. 이는 (1)의 증명에서 얻은 교차곱의 기술을 통해 $X_k$가
연결이라는 사실(보조정리 \ref{models-lemma-regular-model-connected})을 나타낸다.

\medskip\noindent
(4)는 (1), (2), (3)과 보조정리 \ref{models-lemma-recurring-symmetric-real}에서
따른다.
\end{proof}

\begin{lemma}
\label{models-lemma-multiple-fibre-normal-bundle}
상황 \ref{models-situation-regular-model}에서 $d = \gcd(m_1, \ldots, m_n)$라 두고
$D = \sum (m_i/d)C_i$를 유효 Cartier 인자로 보자. 그러면
$\mathcal{O}_X(D)$는 $\Pic(X)$에서 위수가 $d$의 약수이고,
$\mathcal{C}_{D/X}$는 $\Pic(D)$에서 위수가 $d$의 약수인 가역
$\mathcal{O}_D$-가군이다.
\end{lemma}

\begin{proof}
다음이 성립한다.
$$
\mathcal{O}_X(D)^{\otimes d} = \mathcal{O}_X(dD) =
\mathcal{O}_X(X_k) = \mathcal{O}_X
$$
이다(보조정리 \ref{models-lemma-divisor-special-fiber}). $\mathcal{C}_{D/X}$는
$\mathcal{O}_X(-D)$의 당김이므로 결론이 따른다.
\end{proof}

\begin{lemma}
\label{models-lemma-regular-model-field}
\begin{reference}
\cite[보조정리 2.6]{Artin-Winters}
\end{reference}
상황 \ref{models-situation-regular-model}에서 $d = \gcd(m_1, \ldots, m_n)$라 하고
$D = \sum (m_i/d) C_i$를 유효 Cartier 인자로 보자. 그러면 유효 Cartier
인자들의 수열
$$
(X_k)_{red} = Z_0 \subset Z_1 \subset \ldots \subset Z_m = D
$$
이 존재한다. $j = 1, \ldots, m$에 대해 어떤
$i_j \in \{1, \ldots, n\}$가 있어 $Z_j = Z_{j - 1} + C_{i_j}$이고,
$j = 0, \ldots m$에 대해 $H^0(Z_j, \mathcal{O}_{Z_j})$는 $k$의 유한
확대체이다.
\end{lemma}

\begin{proof}
기약화 $D_{red} = (X_k)_{red} = \sum C_i$는 연결이고(보조정리
\ref{models-lemma-regular-model-connected}) $k$ 위에서 고유하다. 따라서
Varieties, 보조정리
\ref{varieties-lemma-proper-geometrically-reduced-global-sections}에 의해
$H^0(D_{red}, \mathcal{O})$는 $k$의 유한 확대체이다. 그러므로
$Z_0 = D_{red} = (X_k)_{red}$에 대한 결과가 성립한다. 이미
$$
(X_k)_{red} = Z_0 \subset Z_1 \subset \ldots \subset Z_t \subset D
$$
를 구성했다고 하자. 여기서 $j = 1, \ldots, t$에 대해 어떤
$i_j \in \{1, \ldots, n\}$가 있어 $Z_j = Z_{j - 1} + C_{i_j}$이고,
$j = 0, \ldots, t$에 대해 $H^0(Z_j, \mathcal{O}_{Z_j})$는 $k$의 유한
확대체이다. $1 \leq a_i \leq m_i/d$가 되도록
$Z_t = \sum a_i C_i$로 쓰자. 모든 $i$에 대해 $a_i = m_i/d$이면
$Z_t = D$이고 증명이 끝난다. 그렇지 않으면 어떤 $i$에 대해
$a_i < m_i/d$이고, 보조정리 \ref{models-lemma-properties-form}에 의해
$(Z_t \cdot Z_t) < 0$이다. 또한 그 보조정리에 의해 $(D \cdot Z_t) = 0$이므로
$(D - Z_t \cdot Z_t) > 0$이다. 따라서 $a_i < m_i/d$이고
$(C_i \cdot Z_t) > 0$인 $i$를 찾을 수 있다.
$Z_{t + 1} = Z_t + C_i$, $i_{t + 1} = i$라 두고 짧은 완전열
$$
0 \to \mathcal{O}_X(-Z_t)|_{C_i} \to \mathcal{O}_{Z_{t + 1}} \to
\mathcal{O}_{Z_t} \to 0
$$
을 생각하자(Divisors, 보조정리 \ref{divisors-lemma-ses-add-divisor}).
$i$의 선택에 의해 $\mathcal{O}_X(-Z_t)|_{C_i}$는 고유 곡선 $C_i$ 위의
음의 차수인 가역층이므로 영이 아닌 대역 단면을 갖지 않는다(Varieties,
보조정리 \ref{varieties-lemma-check-invertible-sheaf-trivial}). 따라서
$H^0(\mathcal{O}_{Z_{t + 1}}) \subset H^0(\mathcal{O}_{Z_t})$는 체이며
(이는 분명하지만 Algebra, 보조정리 \ref{algebra-lemma-integral-under-field}에서도
따른다) $k$의 유한 확대이다. 이로써 수열을 연장했다. 예를 들어
$t \leq \sum (m_i/d - 1)$이므로 과정은 반드시 끝나고 증명도 끝난다.
\end{proof}

\begin{lemma}
\label{models-lemma-regular-model-genus}
\begin{reference}
\cite[보조정리 2.6]{Artin-Winters}
\end{reference}
상황 \ref{models-situation-regular-model}에서 $d = \gcd(m_1, \ldots, m_n)$라 하고
$D = \sum (m_i/d) C_i$를 $X$ 위의 유효 Cartier 인자로 보자. 그러면
$$
1 - g_C = d [\kappa : k] (1 - g_D)
$$
이다. 여기서 $g_C$는 $C$의 종수, $g_D$는 $D$의 종수이고
$\kappa = H^0(D, \mathcal{O}_D)$이다.
\end{lemma}

\begin{proof}
보조정리 \ref{models-lemma-regular-model-field}에 의해 $\kappa$는 $k$의 유한
확대체이다. 또한 $H^0(C, \mathcal{O}_C) = K$이므로 $C$와 $D$의 종수가
정의되며(Algebraic Curves, 정의 \ref{curves-definition-genus} 참조),
$g_C = \dim_K H^1(C, \mathcal{O}_C)$ 및
$g_D = \dim_\kappa H^1(D, \mathcal{O}_D)$이다. Derived Categories of
Schemes, 보조정리 \ref{perfect-lemma-chi-locally-constant-geometric}에 의해
$$
1 - g_C = \chi(C, \mathcal{O}_C) =
\chi(X_k, \mathcal{O}_{X_k}) = \dim_k H^0(X_k, \mathcal{O}_{X_k})
- \dim_k H^1(X_k, \mathcal{O}_{X_k})
$$
라고 주장한다.
$$
\chi(X_k, \mathcal{O}_{X_k}) = d \chi(D, \mathcal{O}_D)
$$
이는 다음 등식 때문에 보조정리를 증명한다.
$$
\chi(D, \mathcal{O}_D) =
\dim_k H^0(D, \mathcal{O}_D) - \dim_k H^1(D, \mathcal{O}_D) =
[\kappa : k](1 - g_D)
$$
$X_k = dD$는 유효 Cartier 인자의 등식이다. 주장을 증명하기 위해
$1 \leq r \leq d$에 대한 귀납법으로
$\chi(rD, \mathcal{O}_{rD}) = r \chi(D, \mathcal{O}_D)$임을 보인다.
$r = 1$인 경우는 자명하다. $1 \leq r < d$이면 짧은 완전열
$$
0 \to \mathcal{O}_X(rD)|_D \to \mathcal{O}_{(r + 1)D} \to
\mathcal{O}_{rD} \to 0
$$
을 생각한다(Divisors, 보조정리 \ref{divisors-lemma-ses-add-divisor}).
Euler 지표의 가법성(Varieties, 보조정리
\ref{varieties-lemma-euler-characteristic-additive})에 의해
$\chi(D, \mathcal{O}_X(rD)|_D) = \chi(D, \mathcal{O}_D)$임을 보이면
충분하다. $\mathcal{O}_X(rD)|_D$는 $\Pic(D)$의 꼬임 원소이고(보조정리
\ref{models-lemma-multiple-fibre-normal-bundle}), 선다발의 차수는 가법적이어서
(Varieties, 보조정리 \ref{varieties-lemma-degree-tensor-product}) 꼬임
가역층에서는 영이므로 이 등식이 성립한다.
\end{proof}

\begin{lemma}
\label{models-lemma-exceptional-curves-dont-meet}
상황 \ref{models-situation-regular-model}에서 $C_i$와 $C_j$가 제1종 예외곡선이고
$C_i \cap C_j \not = \emptyset$인 지표의 쌍 $i, j$가 주어지면
$n = 2$, $m_1 = m_2 = 1$, $C_1 \cong \mathbf{P}^1_k$,
$C_2 \cong \mathbf{P}^1_k$이며, $C_1$과 $C_2$는 $k$-유리점에서 만나고
$C$의 종수는 $0$이다.
\end{lemma}

\begin{proof}
동형 $C_i = \mathbf{P}^1_{\kappa_i}$와
$C_j = \mathbf{P}^1_{\kappa_j}$를 택하자. 스킴 $C_i \cap C_j$는 $C_i$와
$C_j$ 모두에서 공집합이 아닌 유효 Cartier 인자이다. 따라서
$$
(C_i \cdot C_j) = \deg(C_i \cap C_j) \geq \max([\kappa_i: k], [\kappa_j : k])
$$
이다. 첫째 등식은 보조정리 \ref{models-lemma-intersection-pairing}의 증명에서
보였다. 한편 자기교차 $(C_i \cdot C_i)$는 $C_i$ 위에서
$\mathcal{O}_X(C_i)$의 차수와 같고, $C_i$가 제1종 예외곡선이므로 이는
$-[\kappa_i : k]$이다. $C_j$에 대해서도 마찬가지이다. 보조정리
\ref{models-lemma-properties-form}에 의해
$$
0 \geq (C_i + C_j)^2 = -[\kappa_i : k] + 2(C_i \cdot C_j) - [\kappa_j : k]
$$
이다. 이는 $[\kappa_i : k] = \deg(C_i \cap C_j) = [\kappa_j : k]$이고
$(C_i + C_j)^2 = 0$임을 함의한다. 보조정리를 다시 보면 $n = 2$,
$\{1, 2\} = \{i, j\}$, $m_1 = m_2$임을 얻는다. 더욱이 스킴론적 교차
$C_i \cap C_j$는 잉여체가 $\kappa$인 한 점 $p$로 이루어지고
$\kappa_i \to \kappa \leftarrow \kappa_j$는 동형이다.
$D = C_1 + C_2$를 $X$ 위의 유효 Cartier 인자로 보자. $D$는 $C_1$과
$C_2$의 스킴론적 합집합이므로(Divisors, 보조정리
\ref{divisors-lemma-sum-effective-Cartier-divisors-union}) 짧은 완전열
$$
0 \to \mathcal{O}_D \to \mathcal{O}_{C_1} \oplus \mathcal{O}_{C_2} \to
\mathcal{O}_p \to 0
$$
이 있다(Morphisms, 보조정리
\ref{morphisms-lemma-scheme-theoretic-union}).
$C_i \cong \mathbf{P}^1_\kappa$의 코호몰로지를 알고 있으므로(Cohomology of
Schemes, 보조정리 \ref{coherent-lemma-cohomology-projective-space-over-ring}),
긴 완전 코호몰로지 열에서 $H^0(D, \mathcal{O}_D) = \kappa$이고
$H^1(D, \mathcal{O}_D) = 0$임을 얻는다. 보조정리
\ref{models-lemma-regular-model-genus}에 의해
$$
1 - g_C = d[\kappa : k](1 - 0)
$$
이다. 여기서 $d = m_1 = m_2$이다. 따라서 $g_C = 0$,
$d = m_1 = m_2 = 1$, $\kappa = k$이다.
\end{proof}





\section{극소 모형의 유일성}
\label{models-section-uniqueness}

\noindent
일반 올의 종수가 양수이면 극소 모형은 유일하고(보조정리
\ref{models-lemma-minimal-model-unique}), 따라서 적절한 사상 성질을 갖는다
(보조정리 \ref{models-lemma-minimal-model-mapping-property}).

\begin{lemma}
\label{models-lemma-minimal-model-unique}
$C$를 $H^0(C, \mathcal{O}_C) = K$이고 종수가 $> 0$인 $K$ 위 매끄러운
사영곡선이라 하자. $C$의 극소 모형은 유일하다.
\end{lemma}

\begin{proof}
극소 모형의 존재라는 보조정리의 어려운 부분은 이미 증명하였다. 그 증명은
곡면 특이점의 해소에 의존한다. 명제 \ref{models-proposition-exists-minimal-model}을
참조하라. 유일성을 증명하기 위해 $X$와 $Y$가 두 극소 모형이라고 하자.
Resolution of Surfaces, 보조정리
\ref{resolve-lemma-birational-regular-surfaces}에 의해 $S$-사상의 그림
$$
X = X_0 \leftarrow X_1 \leftarrow \ldots \leftarrow X_n = Y_m
\to \ldots \to Y_1 \to Y_0 = Y
$$
이 존재하고 각 사상은 닫힌점에서의 부풀리기이다. 사상
$X_n \to X_{n - 1}$의 예외 올은 제1종 예외곡선 $E$이다. 사상
$X_n = Y_m \to Y$ 아래에서 $E$가 한 점으로 수축된다고 주장한다. 이것이
참이면 Resolution of Surfaces, 보조정리
\ref{resolve-lemma-factor-through-contraction}에 의해 $X_n \to Y$는
$X_{n - 1}$을 통해 인수분해된다. 이 경우 Resolution of Surfaces, 보조정리
\ref{resolve-lemma-proper-birational-regular-surfaces}에 의해
$X_{n - 1} \to Y$도 여전히 예외곡선 수축의 수열이다. 따라서 $n$에 대한
귀납법으로 결론을 얻는다. (시작 경우 $n = 0$은 $Y$에서 끝나는 수축의 수열
$X = Y_m \to \ldots \to Y_1 \to Y_0 = Y$
이 있다는 뜻이다. 그러나 $X$는 극소 모형이므로 제1종 예외곡선을 포함하지
않고, 따라서 $m = 0$이고 $X = Y$이다.)

\medskip\noindent
주장의 증명. $m$에 대한 귀납법으로 임의의 제1종 예외곡선
$E \subset Y_m$이 사상 $Y_m \to Y$에 의해 한 점으로 간다는 것을 보인다.
$m = 0$이면 $Y$가 극소 모형이므로 분명하다. $m > 0$이면
$Y_m \to Y_{m - 1}$가 $E$를 수축하여 끝나거나, $Y_m \to Y_{m - 1}$의
예외 올 $E' \subset Y_m$이 또 다른 제1종 예외곡선이다. $E$와 $E'$는
모두 특수 올의 기약 성분이고 가정에 의해 $g_C > 0$이므로 보조정리
\ref{models-lemma-exceptional-curves-dont-meet}에 의해 $E \cap E' = \emptyset$이다.
그러면 $Y_{m - 1}$에서 $E$의 상은 제1종 예외곡선이다. 실제로
$Y_m \to Y_{m - 1}$은 $E$의 근방에서 동형이다. 귀납가정에 의해
$Y_{m - 1} \to Y$가 이 곡선을 수축하므로 증명이 끝난다.
\end{proof}

\begin{lemma}
\label{models-lemma-minimal-model-mapping-property}
$C$를 $H^0(C, \mathcal{O}_C) = K$이고 종수가 $> 0$인 $K$ 위 매끄러운
사영곡선이라 하자. $X$를 $C$의 극소 모형이라 하고(보조정리
\ref{models-lemma-minimal-model-unique}), $Y$를 $C$의 정칙 고유 모형이라 하자.
그러면 모형의 사상 $Y \to X$가 유일하게 존재하며, 이는 제1종 예외곡선
수축의 수열이다.
\end{lemma}

\begin{proof}
사상 $X \to Y$의 존재와 성질은 보조정리
\ref{models-lemma-pre-exists-minimal-model}과 극소 모형의 유일성에서 즉시 따른다.
$C \subset Y$가 스킴론적으로 조밀하고 $X$가 분리이므로 사상
$Y \to X$는 유일하다(Morphisms, 보조정리
\ref{morphisms-lemma-equality-of-morphisms} 참조).
\end{proof}

\begin{example}
\label{models-example-nonunique-in-genus-zero}
$C$의 종수가 $0$이면 극소 모형은 실제로 유일하지 않다. 닫힌 부분스킴
$$
X \subset \mathbf{P}^2_R
$$
을 생각하자. 이는 $T_1T_2 - \pi T_0^2 = 0$으로 정의된다. 더 정확히
$X$는 $\text{Proj}(R[T_0, T_1, T_2]/(T_1T_2 - \pi T_0^2))$로 정의된다.
그러면 특수 올 $X_k$는 둘 다 $\mathbf{P}^1_k$와 동형인 두 예외곡선
$C_1$, $C_2$의 합집합이다(보조정리
\ref{models-lemma-exceptional-curves-dont-meet}와 정확히 같은 상황이다).
$(0 : 1 : 0)$에서의 사영은 $C_2$를 수축하고 $C_1$을
$\mathbf{P}^1_R$의 특수 올과 동형으로 보내는 사상
$X \to \mathbf{P}^1_R$을 정의한다. $(0 : 0 : 1)$에서의 사영은 $C_1$을
수축하고 $C_2$를 $\mathbf{P}^1_R$의 특수 올과 동형으로 보내는 사상
$X \to \mathbf{P}^1_R$을 정의한다. 더 정확히 이 사상들은 등급
$R$-대수 사상
$$
R[T_0, T_1] \longrightarrow
R[T_0, T_1, T_2]/(T_1T_2 - \pi T_0^2) \longleftarrow
R[T_0, T_2]
$$
에 대응한다. 보조정리 \ref{models-lemma-nonuniqueness}에서 이 현상을 연구한다.
\end{example}








\section{종수에 대한 공식}
\label{models-section-genus-formula}

\noindent
고유 정칙 모형에서 유도되는 조합적 구조에는 제약이 하나 더 있다.

\begin{lemma}
\label{models-lemma-add-component}
상황 \ref{models-situation-regular-model}에서 유효 Cartier 인자
$D, D' \subset X$가 주어지고, 어떤 $i \in \{1, \ldots, n\}$에 대하여
$D' = D + C_i$이며 $D' \subset X_k$라고 하자. 그러면
$$
\chi(X_k, \mathcal{O}_{D'}) - \chi(X_k, \mathcal{O}_D) =
\chi(X_k, \mathcal{O}_X(-D)|_{C_i}) =
-(D \cdot C_i) + \chi(C_i, \mathcal{O}_{C_i})
$$
\end{lemma}

\begin{proof}
두 번째 등식은 (\ref{models-equation-form})의 쌍선형 형식
$(\ \cdot\ )$의 정의와 보조정리 \ref{models-lemma-intersection-pairing}에서 따른다.
첫 번째 등식을 보이기 위해 두 경우를 나눈다.
먼저 $C_i \not \subset D$이면 $D'$는 $D$와 $C_i$의 스킴론적 합집합이다
(Divisors, 보조정리
\ref{divisors-lemma-sum-effective-Cartier-divisors-union}). 따라서 짧은 완전열
$$
0 \to \mathcal{O}_{D'} \to
\mathcal{O}_D \times \mathcal{O}_{C_i} \to
\mathcal{O}_{D \cap C_i} \to 0
$$
을 얻는다(Morphisms, 보조정리
\ref{morphisms-lemma-scheme-theoretic-union}). 또한 완전열
$$
0 \to \mathcal{O}_X(-D)|_{C_i} \to
\mathcal{O}_{C_i} \to \mathcal{O}_{D \cap C_i} \to 0
$$
도 성립하므로(Divisors, 주석 \ref{divisors-remark-ses-regular-section}),
Euler 지표의 가법성에 의해 주장이 성립한다
(Varieties, 보조정리 \ref{varieties-lemma-euler-characteristic-additive}).
한편 $C_i \subset D$이면 Divisors, 보조정리
\ref{divisors-lemma-ses-add-divisor}에 의해 완전열
$$
0 \to \mathcal{O}_X(-D)|_{C_i} \to \mathcal{O}_{D'} \to \mathcal{O}_D \to 0
$$
을 얻으며, 이로부터 보조정리가 곧바로 따른다.
\end{proof}

\begin{lemma}
\label{models-lemma-genus-formula}
상황 \ref{models-situation-regular-model}에서
$$
g_C = 1 + \sum\nolimits_{i = 1, \ldots, n}
m_i\left([\kappa_i : k] (g_i - 1) - \frac{1}{2}(C_i \cdot C_i)\right)
$$
이다. 여기서 $\kappa_i = H^0(C_i, \mathcal{O}_{C_i})$이고,
$g_i$는 $C_i$의 종수이며 $g_C$는 $C$의 종수이다.
\end{lemma}

\begin{proof}
기본적으로 사용할 도구는 다음을 주는 Derived Categories of Schemes,
보조정리 \ref{perfect-lemma-chi-locally-constant-geometric}이다.
$$
1 - g_C = \chi(C, \mathcal{O}_C) =
\chi(X_k, \mathcal{O}_{X_k})
$$
유효 Cartier 인자들의 열
$$
X_k = D_m \supset D_{m - 1} \supset \ldots \supset D_1 \supset D_0 = \emptyset
$$
을 각 $j$에 대해 $D_{j + 1} = D_j + C_{i_j}$가 되도록 택한다.
(매 단계마다 $D_{j+1}$의 0이 아닌 중복도 하나를 1씩 줄이면 이러한 열을
택할 수 있음은 분명하다.) $\chi(\mathcal{O}_{D_0}) = 0$에서 시작하여
보조정리 \ref{models-lemma-add-component}를 적용하면
\begin{align*}
1 - g_C
& =
\chi(X_k, \mathcal{O}_{X_k}) \\
& =
\sum\nolimits_j
\left(-(D_j \cdot C_{i_j}) +  \chi(C_{i_j}, \mathcal{O}_{C_{i_j}})\right) \\
& =
- \sum\nolimits_j
(C_{i_1} + C_{i_2} + \ldots + C_{i_{j - 1}} \cdot C_{i_j}) +
\sum\nolimits_j \chi(C_{i_j}, \mathcal{O}_{C_{i_j}}) \\
& =
-\frac{1}{2}\sum\nolimits_{j \not = j'} (C_{i_{j'}} \cdot C_{i_j}) +
\sum m_i \chi(C_i, \mathcal{O}_{C_i}) \\
& =
\frac{1}{2} \sum m_i(C_i \cdot C_i) + \sum m_i \chi(C_i, \mathcal{O}_{C_i})
\end{align*}
마지막 등식은 설명이 필요할 수 있다. 실제로
$\sum_j C_{i_j} = \sum m_i C_i$이므로 보조정리
\ref{models-lemma-properties-form}에 의해
$(\sum_j C_{i_j} \cdot \sum_j C_{i_j}) = 0$이다. 따라서
$$
0 = \sum\nolimits_{j \not = j'} (C_{i_{j'}} \cdot C_{i_j}) +
\sum m_i(C_i \cdot C_i)
$$
이는 이 곱을 ``비대각'' 항과 ``대각'' 항으로 나누어 얻는다.
$\kappa_i$는 Varieties, 보조정리
\ref{varieties-lemma-regular-functions-proper-variety}에 의해
$k$의 유한 확대체이다. 따라서 $C_i$의 종수가 정의되며
$\chi(C_i, \mathcal{O}_{C_i}) = [\kappa_i : k](1 - g_i)$이다.
이들을 모두 합치고 항을 정리하면
$$
g_C = - \frac{1}{2}\sum m_i(C_i \cdot C_i) +
\sum m_i[\kappa_i : k](g_i - 1) + 1
$$
을 얻으며, 이는 바로 보조정리의 주장이다.
\end{proof}

\begin{lemma}
\label{models-lemma-numerical-type-of-model}
상황 \ref{models-situation-regular-model}에서
$\kappa_i = H^0(C_i, \mathcal{O}_{C_i})$라 하고 $g_i$를 $C_i$의 종수라 하자.
그러면 자료
$$
n, m_i, (C_i \cdot C_j), [\kappa_i : k], g_i
$$
는 종수가 $C$의 종수와 같은 수치형이다.
\end{lemma}

\begin{proof}
(보조정리 \ref{models-lemma-genus-formula}의 증명에서 이 보조정리의 명제에
쓰인 양들이 잘 정의됨을 이미 보였다.) 정의 \ref{models-definition-type}의
조건 (1)--(5)를 확인해야 한다.

\medskip\noindent
조건 (1)은 자명하다.

\medskip\noindent
조건 (2). 행렬 $(C_i \cdot C_j)$의 대칭성은 식 (\ref{models-equation-form})과
보조정리 \ref{models-lemma-intersection-pairing}에서 따른다.
$i \not = j$일 때 $(C_i \cdot C_j)$의 비음성은
보조정리 \ref{models-lemma-properties-form}의 (1)이다.

\medskip\noindent
조건 (3)은 보조정리 \ref{models-lemma-properties-form}의 (3)이다.

\medskip\noindent
조건 (4)는 보조정리 \ref{models-lemma-properties-form}의 (2)이다.

\medskip\noindent
조건 (5)는 $(C_i \cdot C_j)$가 $C_i$ 위의 가역 가군의 차수이고
$[\kappa_i : k]$로 나누어떨어진다는 사실에서 따른다. Varieties,
보조정리 \ref{varieties-lemma-divisible}를 보라.

\medskip\noindent
보조정리 \ref{models-lemma-genus-formula}에서 증명한 종수 공식은 이 수치형의
종수가 명시한 값임을 보여 준다. 정의 \ref{models-definition-genus}를 보라.
\end{proof}

\begin{definition}
\label{models-definition-numerical-type-model}
상황 \ref{models-situation-regular-model}에서 {\it $X$에 결부된 수치형}이란
보조정리 \ref{models-lemma-numerical-type-of-model}에서 설명한 수치형을 말한다.
\end{definition}

\noindent
이제 모형의 극소성과 수치형의 극소성을 대응시킨다.

\begin{lemma}
\label{models-lemma-numerical-type-minimal-model}
상황 \ref{models-situation-regular-model}에서 다음은 서로 동치이다.
\begin{enumerate}
\item $X$는 극소 모형이다.
\item $X$에 결부된 수치형은 극소이다.
\end{enumerate}
\end{lemma}

\begin{proof}
수치형이 극소이면 $g_i = 0$이고
$(C_i \cdot C_i) = -[\kappa_i: k]$인 $i$는 존재하지 않는다.
정의 \ref{models-definition-type-minimal}을 보라. 따라서 어떤 곡선 $C_i$도
제1종 예외곡선이 아님은 분명하다.

\medskip\noindent
역으로 수치형이 극소가 아니라고 하자. 그러면
$g_i = 0$이고 $(C_i \cdot C_i) = -[\kappa_i: k]$인 $i$가 존재한다.
이는 $C_i$가 제1종 예외곡선임을 함의한다고 주장한다. 실제로 가역층
$\mathcal{O}_X(-C_i)|_{C_i}$는 $C_i$를 $k$ 위의 고유곡선으로 보면 차수가
$-(C_i \cdot C_i) = [\kappa_i : k]$이고, 따라서 $C_i$를 $\kappa_i$ 위의
고유곡선으로 보면 차수가 $1$이다. Algebraic Curves, 명제
\ref{curves-proposition-projective-line}을 적용하면 $\kappa_i$ 위의 스킴으로서
$C_i \cong \mathbf{P}^1_{\kappa_i}$임을 얻는다. 체 위의 $\mathbf{P}^1$의
Picard 군은 $\mathbf{Z}$이므로, $X$ 안에서 $C_i$의 법선층은
$\mathcal{O}_{\mathbf{P}_{\kappa_i}}(-1)$과 동형이다. 이로써 증명이 끝난다.
\end{proof}

\begin{remark}
\label{models-remark-numerical-type-not-from-model}
종수가 음수인 수치형이 존재한다는 단순한 이유만으로도 모든 수치형이
모형에서 나오는 것은 아니다. 또한 양의 종수를 가진 극소 수치형 중에서도
(어떤 이산 값매김환의 분수체 위의) 매끄럽고 사영적이며 기하적으로
기약인 곡선의 (그 이산 값매김환 위의) 모형에 결부된 수치형이 아닌 것이
존재한다. 간단한 예는
$n = 1$, $m_1 = 1$, $a_{11} = 0$, $w_1 = 6$, $g_1 = 1$.
이다. 실제로 이 경우 특수 올 $X_k$는 $k$의 차수 $6$인 확대체 $\kappa$
위에 놓이므로 기하적으로 연결일 수 없다. 이는 일반 올이 기하적으로
연결이라는 사실과 모순이다(More on Morphisms, 보조정리
\ref{more-morphisms-lemma-geometrically-connected-fibres-towards-normal}).
마찬가지 이유로 $n = 2$, $m_1 = m_2 = 1$,
$-a_{11} = -a_{22} = a_{12} = a_{21} = 6$,
$w_1 = w_2 = 6$, $g_1 = g_2 = 1$도 한 예가 된다(세부사항은 생략한다).
그러나 $w_i$들의 최대공약수가 $1$인 예는 알려져 있지 않다.
\end{remark}

\begin{lemma}
\label{models-lemma-numerical-type-rational-point}
상황 \ref{models-situation-regular-model}에서 $C$가 $K$-유리점을 가진다고 하자.
그러면
\begin{enumerate}
\item $X_k$는 $k$ 위에서 $X_k$의 매끄러운 점인 $k$-유리점 $x$를 가진다.
\item $x \in C_i$이면 $H^0(C_i, \mathcal{O}_{C_i}) = k$이고 $m_i = 1$이다.
\item $H^0(X_k, \mathcal{O}_{X_k}) = k$이고, $X_k$의 종수는 $C$의 종수와 같다.
\end{enumerate}
\end{lemma}

\begin{proof}
사상 $X \to \Spec(R)$가 고유이므로, 고유성의 값매김 판정법에 의해
$K$-유리점은 사상 $a : \Spec(R) \to X$로 연장된다
(Morphisms, 보조정리 \ref{morphisms-lemma-characterize-proper}).
$x \in X$를 $\Spec(R)$의 닫힌점의 $a$에 의한 상이라 하자.
그러면 $a$는 $R$-대수 준동형
$\psi : \mathcal{O}_{X, x} \to R$
에 대응한다(Schemes, 절 \ref{schemes-section-points}).
$R$에서 $\pi$의 상은 $\mathfrak m_R^2$에 속하지 않으므로
$\pi \not \in \mathfrak m_x^2$이다. 따라서
$\mathcal{O}_{X_k, x} = \mathcal{O}_{X, x}/\pi \mathcal{O}_{X, x}$는
정칙이다(Algebra, 보조정리 \ref{algebra-lemma-regular-ring-CM}).
그러므로 Algebra, 보조정리 \ref{algebra-lemma-separable-smooth}에 의해
$X_k \to \Spec(k)$는 $x$에서 매끄럽다. 이에 따라 $x$는 $X_k$의 유일한
기약성분 $C_i$에 속하고, $\mathcal{O}_{C_i, x} = \mathcal{O}_{X_k, x}$이며
$m_i = 1$이다. $C_i$가 $k$-유리점을 가진다는 사실은 체
$\kappa_i = H^0(C_i, \mathcal{O}_{C_i})$
(Varieties, 보조정리
\ref{varieties-lemma-regular-functions-proper-variety})가 $k$와 같음을
함의한다. 이로써 (1)을 증명했다. 또한
$H^0(X_k, \mathcal{O}_{X_k}) = k$
이다. 왜냐하면 $H^0(X_k, \mathcal{O}_{X_k})$는 $k$의 확대체이고
(보조정리 \ref{models-lemma-regular-model-field}), 이것이
$H^0(C_i, \mathcal{O}_{C_i}) = k$로 사상되기 때문이다.
종수의 등식은 보조정리 \ref{models-lemma-regular-model-genus}에서 따른다.
\end{proof}

\begin{lemma}
\label{models-lemma-genus-reduction-smaller}
상황 \ref{models-situation-regular-model}에서 $X$가 극소 모형이고,
$\gcd(m_1, \ldots, m_n) = 1$이며
$H^0((X_k)_{red}, \mathcal{O}) = k$라고 하자. 그러면 사상
$$
H^1(X_k, \mathcal{O}_{X_k}) \to H^1((X_k)_{red}, \mathcal{O}_{(X_k)_{red}})
$$
은 전사이며, $(X_k)_{red} \not = X_k$이면 그 핵은 자명하지 않다.
\end{lemma}

\begin{proof}
$X_k$ 위에서 차수 $\geq 2$인 코호몰로지가 소멸하므로
(Cohomology, 명제 \ref{cohomology-proposition-vanishing-Noetherian}),
$X_k$ 위의 아벨 층의 모든 전사는 $H^1$ 위의 전사를 유도한다.
보조정리 \ref{models-lemma-regular-model-field}의 열
$$
(X_k)_{red} = Z_0 \subset Z_1 \subset \ldots \subset Z_m = X_k
$$
을 생각하자. 체의 사상
$H^0(Z_j, \mathcal{O}_{Z_j}) \to
H^0((X_k)_{red}, \mathcal{O}_{(X_k)_{red}}) = k$
들이 단사이므로 $j = 0, \ldots, m$에 대하여
$H^0(Z_j, \mathcal{O}_{Z_j}) = k$임을 얻는다. 따라서
$H^0(X_k, \mathcal{O}_{X_k}) \to H^0(Z_{m - 1}, \mathcal{O}_{Z_{m - 1}})$
는 전사이다. $C = C_{i_m}$라 하자. 그러면 $X_k = Z_{m - 1} + C$이다.
$\mathcal{L} = \mathcal{O}_X(-Z_{m - 1})|_C$라 하면,
$\mathcal{L}$은 가역 $\mathcal{O}_C$-가군이다.
보조정리 \ref{models-lemma-regular-model-field}의 증명에서와 같이 완전열
$$
0 \to \mathcal{L} \to \mathcal{O}_{X_k} \to \mathcal{O}_{Z_{m - 1}} \to 0
$$
이 $X_k$ 위의 연접층의 완전열이다. 따라서 짧은 완전열
$$
0 \to
H^1(C, \mathcal{L}) \to H^1(X_k, \mathcal{O}_{X_k}) \to
H^1(Z_{m - 1}, \mathcal{O}_{Z_{m - 1}}) \to 0
$$
$k$ 위의 $C$에서 $\mathcal{L}$의 차수는
$$
(C \cdot -Z_{m - 1}) = (C \cdot C - X_k) = (C \cdot C)
$$
$\kappa = H^0(C, \mathcal{O}_C)$, $w = [\kappa : k]$라 하자.
가역층의 차수 정의에 의해
$$
\chi(C, \mathcal{L}) =
\chi(C, \mathcal{O}_C) + (C \cdot C) =
w(1 - g_C) + (C \cdot C)
$$
이다. 여기서 $g_C$는 $C$의 종수이다. $X$가 극소이므로 $C$는 제1종
예외곡선이 아니며(보조정리 \ref{models-lemma-numerical-type-minimal-model}의
증명을 보라), 따라서 이 식은 $<0$이다. 그러므로
$\dim_k H^1(C, \mathcal{L}) > 0$이고 증명이 끝난다.
\end{proof}

\begin{lemma}
\label{models-lemma-genus-reduction-bigger-than}
상황 \ref{models-situation-regular-model}에서 $X_k$가 $X_k \to \Spec(k)$의
매끄러운 점인 $k$-유리점 $x$를 가진다고 하자. 그러면
$$
\dim_k H^1((X_k)_{red}, \mathcal{O}_{(X_k)_{red}}) \geq
g_{top} + g_{geom}(X_k/k)
$$
이다. 여기서 $g_{geom}$은 Algebraic Curves, 절
\ref{curves-section-genus-geometric-genus}에서 정의한 것이고,
$g_{top}$은 $X_k$에 결부된 수치형
(정의 \ref{models-definition-numerical-type-model})의 위상적 종수
(정의 \ref{models-definition-top-genus})이다.
\end{lemma}

\begin{proof}
연결이고 축소된 모든 유효 Cartier 인자
$D \subset (X_k)_{red}$ 중 $x$를 포함하는 것에 대하여, $D$의 기약성분
개수에 대한 귀납법으로 부등식
$$
\dim_k H^1(D, \mathcal{O}_D) \geq g_{top}(D) + g_{geom}(D/k)
$$
을 증명하겠다. 여기서 $g_{top}(D) = 1 - m + e$이고, $m$은 $D$의
기약성분의 개수이며 $e$는 서로 만나는 $D$의 성분들의 순서 없는 쌍의
개수이다.

\medskip\noindent
기초 단계: $D$가 하나의 기약성분을 가진다고 하자. 그러면 $D = C_i$는
$x$를 포함하는 유일한 기약성분이다. 이 경우
$\dim_k H^1(D, \mathcal{O}_D) = g_i$이고 $g_{top}(D) = 0$이다.
$C_i$는 $k$-유리 매끄러운 점을 가지므로 기하적으로 정역이다
(Varieties, 보조정리
\ref{varieties-lemma-variety-with-smooth-rational-point}).
따라서 $g_i$는 $C_{i, \overline{k}}$의 종수이다
(Algebraic Curves, 보조정리 \ref{curves-lemma-genus-base-change}).
또한 $g_{geom}(D/k)$는 $C_{i, \overline{k}}$의 정규화
$C_{i, \overline{k}}^\nu$의 종수이다. 정규화 사상
$C_{i, \overline{k}}^\nu \to C_{i, \overline{k}}$에 Algebraic Curves,
보조정리 \ref{curves-lemma-genus-normalization}을 적용하면
\begin{equation}
\label{models-equation-genus-change-special-component}
\text{종수 } C_{i, \overline{k}} \geq
\text{종수 } C_{i, \overline{k}}^\nu
\end{equation}
위 사실들을 합치면 이 경우
$\dim_k H^1(D, \mathcal{O}_D) \geq g_{top}(D) + g_{geom}(D/k)$임을 얻는다.

\medskip\noindent
귀납 단계. $D$가 $1$개보다 많은 기약성분을 가진다고 하자. 그러면
$x \in D'$이고 $D'$도 연결인 채로 $D = C_i + D'$로 쓸 수 있다.
이는 독자에게 맡기는 그래프 이론의 연습문제이다(힌트: $C_i$를 $D$의
성분 중 $x$에서 가장 먼 것으로 택하라). 불변량들이 어떻게 변하는지
계산하자. $x \in D'$이므로
$H^0(D, \mathcal{O}_D) = H^0(D', \mathcal{O}_{D'}) = k$이다.
층의 짧은 완전열
$$
0 \to \mathcal{O}_D \to \mathcal{O}_{C_i} \oplus \mathcal{O}_{D'}
\to \mathcal{O}_{C_i \cap D'} \to 0
$$
을 살펴보고(Morphisms, 보조정리
\ref{morphisms-lemma-scheme-theoretic-union}) Euler 지표의 가법성을
사용하면
\begin{align*}
\dim_k H^1(D, \mathcal{O}_D) - \dim_k H^1(D', \mathcal{O}_{D'})
& =
-\chi(\mathcal{O}_{C_i}) + \chi(\mathcal{O}_{C_i \cap D'}) \\
& =
w_i(g_i - 1) + \sum\nolimits_{C_j \subset D'} a_{ij}
\end{align*}
을 얻는다. 보조정리 \ref{models-lemma-numerical-type-of-model}에서와 같이
$w_i = [\kappa_i : k]$, $\kappa_i = H^0(C_i, \mathcal{O}_{C_i})$라 하고,
$g_i$를 $C_i$의 종수, $a_{ij} = (C_i \cdot C_j)$라 하였다. 또한
$$
g_{top}(D) - g_{top}(D') = -1 +
\sum\nolimits_{C_j \subset D'\text{이고 }C_i\text{와 만나는}} 1
$$
이고
$$
g_{geom}(D/k) - g_{geom}(D'/k) = g_{geom}(C_i/k)
$$
이다(Algebraic Curves, 보조정리
\ref{curves-lemma-bound-geometric-genus}). 이를 귀납가정과 합치면
$$
w_i g_i - g_{geom}(C_i/k) +
\sum\nolimits_{C_j \subset D'\text{이고 } C_i\text{와 만나는}} (a_{ij} - 1) - (w_i - 1)
$$
가 비음임을 보이면 충분하다. 실제로
\begin{equation}
\label{models-equation-genus-change}
w_i g_i \geq [\kappa_i : k]_s g_i \geq g_{geom}(C_i/k)
\end{equation}
이며, 두 번째 부등식은 Algebraic Curves, 보조정리
\ref{curves-lemma-bound-geometric-genus-curve}에서 따른다.
한편 $w_i$가 $a_{ij}$를 나누므로
(Varieties, 보조정리 \ref{varieties-lemma-divisible}),
\begin{equation}
\label{models-equation-change-intersections}
\sum\nolimits_{C_j \subset D'\text{이고 } C_i\text{와 만나는}} (a_{ij} - 1) - (w_i - 1)
\geq 0
\end{equation}
이다. 실제로 $C_i$와 만나는 $C_j \subset D'$가 적어도 하나 존재한다.
\end{proof}

\begin{lemma}
\label{models-lemma-equality-genus-reduction-bigger-than}
보조정리 \ref{models-lemma-genus-reduction-bigger-than}에서 등호가 성립하면
\begin{enumerate}
\item $x$를 포함하는 $X_k$의 유일한 기약성분은 $k$ 위에서 매끄럽고
사영적이며 기하적으로 기약인 곡선이다.
\item $C \subset X_k$가 다른 기약성분이면
$\kappa = H^0(C, \mathcal{O}_C)$는 $k$의 유한 분리 확대체이고,
$C$는 $\kappa$-유리점을 가지며 $C$는 $\kappa$ 위에서 매끄럽다.
\end{enumerate}
\end{lemma}

\begin{proof}
보조정리 \ref{models-lemma-genus-reduction-bigger-than}의 증명을 살펴보면,
등호가 성립하기 위해서는 부등식
(\ref{models-equation-genus-change-special-component}),
(\ref{models-equation-genus-change}),
(\ref{models-equation-change-intersections})
모두에서 등호가 성립해야 함을 알 수 있다.

\medskip\noindent
$C_i$를 $x$를 포함하는 기약성분이라 하자.
(\ref{models-equation-genus-change-special-component})의 등호와
Algebraic Curves, 보조정리 \ref{curves-lemma-genus-normalization}에 의해
$C_{i, \overline{k}}^\nu \to C_{i, \overline{k}}$는 동형이다.
따라서 $C_{i, \overline{k}}$는 매끄럽고 (1)이 성립한다.

\medskip\noindent
다음으로 $C_i \subset X_k$가 다른 기약성분이라고 하자. 보조정리
\ref{models-lemma-genus-reduction-bigger-than}의 증명에서 귀납 단계와 같이
$D = D' + C_i$라고 해도 된다. (\ref{models-equation-genus-change})의 등호는
$\kappa_i/k$가 유한 분리 확대임을 곧바로 함의한다.
(\ref{models-equation-change-intersections})의 등호는 어떤 $j$에 대해
$a_{ij} = 1$이거나, $C_i$와 만나는 유일한 $C_j \subset D'$가 존재하고
$a_{ij} = w_i$임을 함의한다. 어느 경우든 $C_i$는 $\kappa_i$-유리점
$c$를 가지며, 스킴론적으로 $c = C_i \cap C_j$이다.
$\mathcal{O}_{X, c}$가 정칙 국소환이므로, 이는 $C_i$와 $C_j$의 국소
방정식들이 국소환 $\mathcal{O}_{X, c}$의 정칙 매개변수계를 이룸을
뜻한다. 따라서 $\mathcal{O}_{C_i, c}$는 정칙이다
(Algebra, 보조정리 \ref{algebra-lemma-regular-ring-CM}).
그러므로 $C_i \to \Spec(\kappa_i)$는 $c$에서 매끄럽다
(Algebra, 보조정리 \ref{algebra-lemma-separable-smooth}).
이에 따라 $C_i$는 $\kappa_i$ 위에서 기하적으로 정역이다
(Varieties, 보조정리
\ref{varieties-lemma-variety-with-smooth-rational-point}).
끝으로 $C_i$가 $\kappa_i$ 위에서 매끄러움을 보여야 한다. 다음을 관찰하자.
$$
C_{i, \overline{k}} = C_i \times_{\Spec(k)} \Spec(\overline{k})
= \coprod\nolimits_{\kappa_i \to \overline{k}}
C_i \times_{\Spec(\kappa_i)} \Spec(\overline{k})
$$
여기에는 $[\kappa_i : k]$개의 성분이 있다. 따라서 $C_i$가 $\kappa_i$
위에서 매끄럽지 않다면 이 곡선들은 각각 매끄럽지 않고, 따라서 정규가
아니어서 정규화 사상이 종수를 낮춘다
(Algebraic Curves, 보조정리 \ref{curves-lemma-genus-normalization}).
그러면 $C_i/k$의 기하적 종수도 낮아져
$[\kappa_i : k] g_i = g_{geom}(C_i/k)$와 모순이므로, 이런 일은 불가능하다.
\end{proof}








\section{예외곡선의 수축}
\label{models-section-contract}

\noindent
다음 보조정리는 고유 정칙 모형에서 제1종 예외곡선을 수축할 때 교차수가
어떻게 변하는지 설명한다. 주된 목적은 이를 보조정리 \ref{models-lemma-contract}에서
도입한 수치적 수축과 비교하는 데 있다. 기하적 수축과 수치적 수축은
주석 \ref{models-remark-compare-contractions}에서 비교한다.

\begin{lemma}
\label{models-lemma-blowdown-regular-model}
상황 \ref{models-situation-regular-model}에서 $C_n$이 제1종 예외곡선이라고 하자.
$f : X \to X'$를 $C_n$의 수축이라 하고 $C'_i = f(C_i)$라 하자.
$X'_k = \sum m'_i C'_i$로 쓰자. 그러면 $X'$, $C'_i$,
$i = 1, \ldots, n' = n - 1$, $m'_i = m_i$는 상황
\ref{models-situation-regular-model}과 같으며, 다음이 성립한다.
\begin{enumerate}
\item $i, j < n$에 대하여
$(C'_i \cdot C'_j) =
(C_i \cdot C_j) - (C_i \cdot C_n) (C_j \cdot C_n) /(C_n \cdot C_n)$,
\item $i < n$이고 $C_i \cap C_n \not = \emptyset$이면 사상
$\kappa_i \leftarrow \kappa'_i \rightarrow \kappa_n$이 존재한다.
\end{enumerate}
여기서 $\kappa_i = H^0(C_i, \mathcal{O}_{C_i})$이고
$\kappa'_i = H^0(C'_i, \mathcal{O}_{C'_i})$이다.
\end{lemma}

\begin{proof}
Resolution of Surfaces, 보조정리 \ref{resolve-lemma-contract-ample}에 의해
$X'$가 정칙이고 $R$ 위에서 사영이 되도록 하는 사상
$f : X \to X'$로 $C_n$을 수축할 수 있다. 따라서 $X'$는 상황
\ref{models-situation-regular-model}과 같다. $x \in X'$를 $C_n$의 상이라 하자.
$f$는 동형 $X \setminus C_n \to X' \setminus \{x\}$를 정의하므로
$i < n$에 대해 $m'_i = m_i$임은 분명하다.

\medskip\noindent
보조정리의 (2)는 사상 $C_i \to C'_i$와 $C_n \to x \to C'_i$가
존재한다는 사실에서 곧바로 따른다.

\medskip\noindent
Divisors, 보조정리
\ref{divisors-lemma-blow-up-pullback-effective-Cartier}에 의해 역상
$f^{-1}C'_i$가 정의된다. Divisors, 보조정리
\ref{divisors-lemma-effective-Cartier-divisor-is-a-sum}에 의해 어떤
$e_i \geq 0$에 대하여 $f^{-1}C'_i = C_i + e_i C_n$이다.
$\mathcal{O}_X(C_i + e_i C_n) = \mathcal{O}_X(f^{-1}C'_i) =
f^*\mathcal{O}_{X'}(C'_i)$
이고(Divisors, 보조정리
\ref{divisors-lemma-pullback-effective-Cartier-divisors}), 가역층의
당김은 $C_n$에 제한하면 자명한 가역층이 되므로
$$
0 = \deg_{C_n}(\mathcal{O}_X(C_i + e_i C_n)) =
(C_i + e_i C_n \cdot C_n) = (C_i \cdot C_n) + e_i(C_n \cdot C_n)
$$
$f_j = f|_{C_j} : C_j \to C_j$는 $k$ 위의 고유곡선 사이의 고유
쌍유리 사상이므로,
$\deg_{C'_j}(\mathcal{O}_{X'}(C'_i)|_{C'_j})$는
$\deg_{C_j}(f_j^*\mathcal{O}_{X'}(C'_i)|_{C'_j})$와 같다
(Varieties, 보조정리 \ref{varieties-lemma-degree-birational-pullback}).
가환 그림
$$
\xymatrix{
C_j \ar[r] \ar[d]_{f_j} & X \ar[d]^f \\
C'_j \ar[r] & X'
}
$$
을 살펴보고 Divisors, 보조정리
\ref{divisors-lemma-pullback-effective-Cartier-divisors}를 사용하면
$$
(C'_i \cdot C'_j) = \deg_{C'_j}(\mathcal{O}_{X'}(C'_i)|_{C'_j}) =
\deg_{C_j}(\mathcal{O}_X(C_i + e_i C_n)) = (C_i + e_i C_n \cdot C_j)
$$
앞에서 구한 $e_i$의 공식을 대입하면 (1)이 성립한다.
\end{proof}

\begin{remark}
\label{models-remark-genus-change}
보조정리 \ref{models-lemma-blowdown-regular-model}의 상황에서는 $C_i$의 종수
$g_i$와 $C'_i$의 종수 $g'_i$ 사이의 관계도 정확히 기술할 수 있다.
그 공식은
$$
g'_i = \frac{w_i}{w'_i}(g_i - 1) + 1 +
\frac{(C_i \cdot C_n)^2 - w_n(C_i \cdot C_n)}{2w'_iw_n}
$$
이다. 여기서 $w_i = [\kappa_i : k]$, $w_n = [\kappa_n : k]$,
$w'_i = [\kappa'_i : k]$이다. 이를 증명하기 위해 짧은 완전열
$$
0 \to \mathcal{O}_{X'}(-C'_i) \to \mathcal{O}_{X'} \to
\mathcal{O}_{C'_i} \to 0
$$
과 그 $X$로의 당김
$$
0 \to \mathcal{O}_X(-C'_i - e_iC_n) \to \mathcal{O}_X \to
\mathcal{O}_{C_i + e_i C_n} \to 0
$$
을 생각하자. 여기서 $e_i$는 보조정리
\ref{models-lemma-blowdown-regular-model}의 증명에서와 같다.
$X'$ 위의 모든 가역 가군 $\mathcal{L}$에 대해
$Rf_*f^*\mathcal{L} = \mathcal{L}$이므로(세부사항은 생략한다),
$$
Rf_*\mathcal{O}_{C_i + e_i C_n} = \mathcal{O}_{C'_i}
$$
임을 얻는다. 이는 $X'_k$ 위의 연접층의 복합체로서의 등식이다.
따라서 양변의 Euler 지표는 같고, 이는 $X_k$ 위의
$\mathcal{O}_{C_i + e_i C_n}$의 Euler 지표와 일치한다. 완전열
$$
0 \to \mathcal{O}_{C_i + e_i C_n} \to
\mathcal{O}_{C_i} \oplus \mathcal{O}_{e_iC_n} \to
\mathcal{O}_{C_i \cap e_iC_n} \to 0
$$
을 사용하고 $\mathcal{O}_{e_iC_n}$을 다시 여과하면(세부사항은 생략한다)
$$
\chi(\mathcal{O}_{C'_i}) =
\chi(\mathcal{O}_{C_i}) - {e_i + 1 \choose 2}(C_n \cdot C_n)
- e_i(C_i \cdot C_n)
$$
$e_i = -(C_i \cdot C_n)/(C_n \cdot C_n)$이고
$(C_n \cdot C_n) = -w_n$이므로, 이는 이 주석의 첫머리에 제시한
공식을 준다. 이 결과가 필요해지면 이를 보조정리로 정식화하고
상세한 증명을 제시할 것이다.
\end{remark}

\begin{remark}
\label{models-remark-compare-contractions}
$f : X \to X'$가 보조정리 \ref{models-lemma-blowdown-regular-model}에서와
같다고 하자. $n, m_i, a_{ij}, w_i, g_i$를 $X$에 결부된 수치형이라 하고,
$n', m'_i, a'_{ij}, w'_i, g'_i$를 $X'$에 결부된 수치형이라 하자.
보조정리 \ref{models-lemma-blowdown-regular-model}과 주석
\ref{models-remark-genus-change}에 의해, 이는 $w'_i$의 값을 제외하면
보조정리 \ref{models-lemma-contract}의 수치형 수축과 일치한다.
기하적 상황에서 $w'_i$는 $w_i$와 $w_n$을 모두 나누는 어떤 양의
정수이다. 수치적 경우에는 $w_i$를 나누면서 공식으로 주어지는 $g'_i$가
정수가 되게 하는 가장 큰 정수로 $w'_i$를 택했다. 이 선택은 수치적
상황에서 수치형의 Picard 군을 비교하는 데 유용하다. 보조정리
\ref{models-lemma-contract-picard-group}을 보라. 다만 이후에는 단사성만
사용하며, 이 단사성은 더 작은 $w'_i$에 대해서도 성립한다.
\end{remark}

\begin{lemma}
\label{models-lemma-nonuniqueness}
$C$를 $H^0(C, \mathcal{O}_C) = K$이고 종수가 $0$인 $K$ 위의 매끄러운
사영곡선이라 하자. $C$의 극소 모형이 둘 이상이면, 모든 극소 모형의
특수 올은 $\mathbf{P}^1_k$와 동형이다.
\end{lemma}

\noindent
이 보조정리는 서로 동형이 아닌 두 극소 모형 사이의 쌍유리 변환을
예 \ref{models-example-nonunique-in-genus-zero}에서와 같은 초등 변환들의 열로
분해할 수 있다는 명제로 강화할 수 있다. 이 결과가 필요해지면 여기에서
정확히 정식화하고 증명할 것이다.

\begin{proof}
$X$를 $C$의 한 극소 모형이라 하자. $X$에 결부된 수치형은 종수가
$0$이고 극소이다(정의 \ref{models-definition-numerical-type-model} 및 보조정리
\ref{models-lemma-numerical-type-minimal-model}). 따라서 보조정리
\ref{models-lemma-genus-zero}에 의해 $X_k$는 축소되고 기약이며,
$H^0(X_k, \mathcal{O}_{X_k}) = k$이고 종수가 $0$이다.
$Y$를 $X$와 동형이 아닌 $C$의 두 번째 극소 모형이라 하자.
Resolution of Surfaces, 보조정리
\ref{resolve-lemma-birational-regular-surfaces}에 의해 $S$-사상들의 그림
$$
X = X_0 \leftarrow X_1 \leftarrow \ldots \leftarrow X_n = Y_m
\to \ldots \to Y_1 \to Y_0 = Y
$$
이 존재하며, 각 사상은 닫힌점에서의 부풀리기이다. $m$에 대한 귀납법으로
보조정리를 증명하자. 기초 단계는 $m = 0$이다. $Y$가 극소라고 가정했으므로
이 경우에는 $n = 0$이어야 하지만, $X$와 $Y$는 동형이 아니어서 이러한
경우는 일어나지 않는다. 따라서 확인할 것이 없다.

\medskip\noindent
계속하기 전에, $X_k$와 $Y_k$가 모두 기약이므로 $n + 1 = m + 1$은
$X_n = Y_m$의 특수 올의 기약성분 개수와 같다는 점에 유의하자.
또 아래에서 사용할 관찰은 다음과 같다. $X' \to X''$가 $C$의 고유 정칙
모형 사이의 사상이면, $X' \to X''$는 $X''$의 어떤 열린집합 위에서
동형이고 그 여집합은 $X''$의 특수 올의 닫힌점들로 이루어진 유한집합이다.
Varieties, 보조정리
\ref{varieties-lemma-modification-normal-iso-over-codimension-1}을 보라.
실제로 이러한 $X' \to X''$는 닫힌점들에서의 부풀리기들의 합성이고
(Resolution of Surfaces, 보조정리
\ref{resolve-lemma-proper-birational-regular-surfaces}), 부풀리기의 개수는
$X'$와 $X''$의 특수 올의 기약성분 개수의 차이이다.

\medskip\noindent
$E_i \subset Y_i$ ($m \geq i \geq 1$)를 사상 $Y_i \to Y_{i - 1}$에
의해 수축되는 곡선이라 하자. $E_i$가 $Y_i$의 특수 올에서 중복도 $>1$을
가지는 가장 큰 첨자 $i$를 택하자. 그러면 그 이후의 부풀리기
$Y_m \to \ldots \to Y_{i + 1} \to Y_i$는 $E_i$ 위에서 동형이다.
그렇지 않으면 어떤 $j > i$에 대해 $E_j$의 중복도가 $>1$이기 때문이다.
$E \subset Y_m$을 $E_i$의 역상이라 하자. 방금 말한 바에 의해
$E \subset Y_m$은 제1종 예외곡선이다. $Y_m \to Y'$를 $E$의 수축이라
하자. 이는 Resolution of Surfaces, 보조정리
\ref{resolve-lemma-contract-when-quasi-projective}에 의해 존재한다.
$X_k$가 축소되어 있으므로 사상 $Y_m \to X$는 $E$를 수축해야 한다.
따라서 Resolution of Surfaces, 보조정리
\ref{resolve-lemma-factor-through-contraction}에 의해 사상 $Y' \to Y$와
$Y' \to X$가 존재하며, 이들은 많아야 $n - 1 = m - 1$개의 예외곡선
수축의 합성이다(위 논의를 보라). 이제 $m$에 대한 귀납가정을 적용한다.
결론적으로 모든 $X_i$와 $Y_i$의 특수 올이 축소되어 있다고 가정해도 된다.

\medskip\noindent
$X_i$와 $Y_i$의 올이 축소되어 있으므로 부풀리기
$X_i \to X_{i - 1}$와 $Y_i \to Y_{i - 1}$는 특수 올의 정칙점인
닫힌점에서 일어나야 한다. 실제로 $X''$가 $C$의 정칙 모형이고
$x \in X''$가 특수 올의 닫힌점이며 $\pi \in \mathfrak m_x^2$이면,
$x$에서의 부풀리기 $X' \to X''$의 예외 올 $E$는 $X'$의 특수 올에서
적어도 $2$의 중복도를 가진다(국소 계산은 생략한다). 따라서 주장한 대로
$\mathcal{O}_{X''_k, x} = \mathcal{O}_{X'', x}/\pi$는 정칙이다
(Algebra, 보조정리 \ref{algebra-lemma-regular-ring-CM}).
특히 $x$는 자신을 포함하는 $X''_k$의 유일한 기약성분 $Z'$ 위의
Cartier 인자이다(Varieties, 보조정리
\ref{varieties-lemma-regular-point-on-curve}). 따라서 $Z'$의 고유변환
$Z \subset X'$는 $Z'$에 동형으로 사상된다(Divisors, 보조정리
\ref{divisors-lemma-strict-transform} 및
\ref{divisors-lemma-blow-up-effective-Cartier-divisor}를 사용하라).
달리 말하면 $X_i$의 기약성분 $Z$가 사상 $X_i \to X_j$ ($i > j$)에서
수축되지 않으면 그 상에 동형으로 사상된다.

\medskip\noindent
이제 보조정리를 증명할 준비가 되었다. $E \subset Y_m$을 사상
$Y_m \to Y_{m - 1}$에 의해 수축되는 제1종 예외곡선이라 하자.
$E$가 사상 $Y_m = X_n \to X$에 의해 수축되면
$Y_{m - 1} \to X$로의 분해가 존재한다(Resolution of Surfaces,
보조정리 \ref{resolve-lemma-factor-through-contraction}).
더욱이 $Y_{m - 1} \to X$는 닫힌점들에서의 부풀리기들의 합성이다
(Resolution of Surfaces, 보조정리
\ref{resolve-lemma-proper-birational-regular-surfaces}).
이 경우 $m$을 줄이고 귀납가정을 적용한다. 끝으로 $E$가 사상
$Y_m \to X$에 의해 수축되지 않는다고 하자. $X_k$가 기약이므로
$E \to X_k$는 전사이고, 위에서 본 바에 의해 이는 동형이다.
따라서 원하는 대로 $X_k$는 사영직선과 동형이다.
\end{proof}





\section{모형의 Picard 군}
\label{models-section-picard-groups-models}

\noindent
$R, K, k, \pi, C, X, n, C_1, \ldots, C_n, m_1, \ldots, m_n$이
상황 \ref{models-situation-regular-model}에서와 같다고 하자.
보조정리 \ref{models-lemma-regular-model-pic}에서 완전열
$$
0 \to \mathbf{Z} \to \mathbf{Z}^{\oplus n} \to
\Pic(X) \to \Pic(C) \to 0
$$
을 얻었다. 이 열을 사용하여 적절한 소수 $\ell$에 대한 Picard 군의
$\ell$-꼬임을 연구하고자 한다.

\begin{lemma}
\label{models-lemma-characterize-trivial}
상황 \ref{models-situation-regular-model}에서
$d = \gcd(m_1, \ldots, m_n)$라 하자. 가역 $\mathcal{O}_X$-가군
$\mathcal{L}$이 다음을 만족한다고 하자.
\begin{enumerate}
\item $C$에 제한하면 자명한 가역 가군이다.
\item 각 $C_i$ 위에서 차수가 $0$이다.
\end{enumerate}
그러면 $\mathcal{L}^{\otimes d} \cong \mathcal{O}_X$이다.
\end{lemma}

\begin{proof}
보조정리 \ref{models-lemma-regular-model-pic}에 의해 어떤
$a_i \in \mathbf{Z}$에 대하여
$\mathcal{L} \cong \mathcal{O}_X(\sum a_i C_i)$이다.
$\mathcal{L}|_{C_j}$의 차수는 $\sum_j a_j(C_i \cdot C_j)$이다. 특히
$(\sum a_i C_i \cdot \sum a_i C_i) = 0$.
따라서 보조정리 \ref{models-lemma-properties-form}에 의해 어떤
$q \in \mathbf{Q}$에 대하여
$(a_1, \ldots, a_n) = q(m_1, \ldots, m_n)$이다.
그러므로 어떤 $l \in \mathbf{Z}$에 대하여
$\mathcal{L} = \mathcal{O}_X(lD)$이며, 여기서
$D = \sum (m_i/d) C_i$는 보조정리
\ref{models-lemma-multiple-fibre-normal-bundle}에서와 같다. 이제 결론이 따른다.
\end{proof}

\begin{lemma}
\label{models-lemma-canonical-map-of-pic}
상황 \ref{models-situation-regular-model}에서 $T$를 $X$에 결부된 수치형이라 하자.
표준 사상
$$
\Pic(C) \to \Pic(T)
$$
이 존재한다. 그 핵은 정확히 다음과 같은 $C$ 위의 가역 가군들로
이루어진다. 즉, $i = 1, \ldots, n$에 대해
$\deg_{C_i}(\mathcal{L}|_{C_i}) = 0$인 $X$ 위의 가역 가군
$\mathcal{L}$의 제한으로 얻어지는 것들이다.
\end{lemma}

\begin{proof}
$\kappa_i = H^0(C_i, \mathcal{O}_{C_i)})$일 때
$w_i = [\kappa_i : k]$임을 상기하자. 또한 $C_i$ 위의 모든 가역 가군의
차수는 $w_i$로 나누어떨어진다(Varieties, 보조정리
\ref{varieties-lemma-divisible}). 따라서 사상
$$
\frac{\deg}{w} : \Pic(X) \to \mathbf{Z}^{\oplus n}, \quad
\mathcal{L} \mapsto
(\frac{\deg(\mathcal{L}|_{C_1})}{w_1}, \ldots,
\frac{\deg(\mathcal{L}|_{C_n})}{w_n})
$$
을 생각할 수 있다. 이 사상에 의한 $\mathcal{O}_X(C_j)$의 상은
$$
((C_j \cdot C_1)/w_1, \ldots, (C_j \cdot C_n)/w_n) =
(a_{1j}/w_1, \ldots, a_{nj}/w_n)
$$
이다. 이는 $T$의 Picard 군을 정의하는 사상
$(a_{ij}/w_i) : \mathbf{Z}^{\oplus n} \to \mathbf{Z}^{\oplus n}$에
의한 $j$번째 기저벡터의 상과 정확히 같다. 정의
\ref{models-definition-picard-group}을 보라. 따라서 보조정리의 표준 사상은
가환 그림
$$
\xymatrix{
\mathbf{Z}^{\oplus n} \ar[r] \ar[d]_{\text{id}} &
\Pic(X) \ar[r] \ar[d]^{\frac{\deg}{w}} &
\Pic(C) \ar[r] \ar[d] & 0 \\
\mathbf{Z}^{\oplus n} \ar[r]^{(a_{ij}/w_i)} &
\mathbf{Z}^{\oplus n} \ar[r] &
\Pic(T) \ar[r] & 0
}
$$
에서 나온다. 두 행은 완전하다(위 행은 보조정리
\ref{models-lemma-regular-model-pic}에 의함). 핵에 대한 기술은 분명하다.
\end{proof}

\begin{lemma}
\label{models-lemma-sequence-torsion}
상황 \ref{models-situation-regular-model}에서
$d = \gcd(m_1, \ldots, m_n)$라 하고, $T$를 $X$에 결부된 수치형이라 하자.
$h \geq 1$을 $d$와 서로소인 정수라 하면 완전열
$$
0 \to \Pic(X)[h] \to \Pic(C)[h] \to \Pic(T)[h]
$$
이 존재한다.
\end{lemma}

\begin{proof}
보조정리 \ref{models-lemma-regular-model-pic}의 완전열에서 $h$-꼬임을 취하면,
$h$와 $d$가 서로소이므로
$0 \to \Pic(X)[h] \to \Pic(C)[h]$의 완전성을 얻는다.
보조정리 \ref{models-lemma-canonical-map-of-pic}의 사상을 사용하면
$\Pic(X)[h]$의 원소들을 소멸시키는 사상
$\Pic(C)[h] \to \Pic(T)[h]$를 얻는다.
역으로 $\xi \in \Pic(C)[h]$의 $\Pic(T)[h]$에서의 상이 0이면,
모든 $i$에 대해 $\deg(\mathcal{L}|_{C_i}) = 0$이고 $C$로의 제한이
$\xi$인 가역 $\mathcal{O}_X$-가군 $\mathcal{L}$을 찾을 수 있다.
보조정리 \ref{models-lemma-characterize-trivial}에 의해
$\mathcal{L}^{\otimes h}$는 $d$-꼬임이다.
$dd' \equiv 1 \bmod h$를 만족하는 정수 $d'$를 택하자. $h$와 $d$가
서로소이므로 이러한 정수가 존재한다. 그러면
$\mathcal{L}^{\otimes dd'}$는 $X$ 위의 $h$-꼬임 가역층이고,
$C$로의 제한은 $\xi$이다.
\end{proof}

\begin{lemma}
\label{models-lemma-torsion-embeds}
상황 \ref{models-situation-regular-model}에서 $h$를 $k$의 표수와 서로소인
정수라 하자. 그러면 사상
$$
\Pic(X)[h] \longrightarrow \Pic((X_k)_{red})[h]
$$
은 단사이다.
\end{lemma}

\begin{proof}
$X \times_{\Spec(R)} \Spec(R/\pi^n)$은 $(X_k)_{red}$의 유한 차수
두꺼워짐임을 관찰하자. 이는 예컨대 Cohomology of Schemes, 보조정리
\ref{coherent-lemma-power-ideal-kills-sheaf}에서 따른다. 따라서 표준 사상
$\Pic(X \times_{\Spec(R)} \Spec(R/\pi^n)) \to \Pic((X_k)_{red})$
은 More on Morphisms, 보조정리
\ref{more-morphisms-lemma-torsion-pic-thickening}과 $h$에 대한 가정에
의해 $h$-꼬임을 동형으로 대응시킨다. 그러므로 $\mathcal{L}$이 $X$ 위의
$h$-꼬임 가역층이고 $(X_k)_{red}$에 제한하면 자명한 층이면,
모든 $n$에 대해 $\mathcal{L}$을
$X \times_{\Spec(R)} \Spec(R/\pi^n)$에 제한한 것도 자명하다. 따라서
\begin{align*}
H^0(X, \mathcal{L})^\wedge
& =
\lim H^0(X \times_{\Spec(R)} \Spec(R/\pi^n),
\mathcal{L}|_{X \times_{\Spec(R)} \Spec(R/\pi^n)}) \\
& \cong
\lim H^0(X \times_{\Spec(R)} \Spec(R/\pi^n),
\mathcal{O}_{X \times_{\Spec(R)} \Spec(R/\pi^n)}) \\
& =
R^\wedge
\end{align*}
을 얻는다. 첫 번째와 마지막 등식에는 형식 함수 정리
(Cohomology of Schemes, 정리 \ref{coherent-theorem-formal-functions})를,
가운데 동형에는 예컨대 More on Algebra, 보조정리
\ref{more-algebra-lemma-isomorphic-completions}를 사용하였다.
$H^0(X, \mathcal{L})$은 유한 $R$-가군이고 $R$은 이산 값매김환이므로,
$H^0(X, \mathcal{L})$은 계수 $1$인 자유 $R$-가군이다.
$s \in H^0(X, \mathcal{L})$를 한 기저원소라 하자. 위의 동형들을 거슬러
추적하면 모든 $n$에 대해
$s|_{X \times_{\Spec(R)} \Spec(R/\pi^n)}$이 자명화를 준다.
$s$의 소멸 자취는 $X$에서 닫혀 있고 $X \to \Spec(R)$는 고유이므로,
원하는 대로 $s$의 소멸 자취는 공집합이다.
\end{proof}






\section{반안정 축약}
\label{models-section-semistable-reduction}

\noindent
이 절에서는 반안정 축약의 의미를 엄밀히 정의한다.

\begin{example}
\label{models-example-blowup}
$R$을 균등화 원소 $\pi$를 갖는 이산 값매김환이라 하자.
$n \geq 0$이 주어졌을 때 환 준동형
$$
R \longrightarrow A = R[x, y]/(xy - \pi^n)
$$
을 생각하자. $X = \Spec(A)$, $S = \Spec(R)$라 하자.
$n = 0$이면 $X \to S$는 매끄럽다. 모든 $n$에 대하여 사상
$X \to S$는 Algebraic Curves, 절
\ref{curves-section-families-nodal}에서 정의한 상대차원 $1$의
마디 이하 특이성을 갖는다. $n = 1$이면 $X$는 정칙이지만,
$n > 1$이면 $(x, y)$가 극대 아이디얼
$\mathfrak m = (\pi, x, y)$을 더 이상 생성하지 않으므로 $X$는
정칙이 아니다. $n > 1$인 경우를 개선하기 위해 $\mathfrak m$에서의
$X$의 부풀리기 $b : X' \to X$를 생각하자. Divisors, 절
\ref{divisors-section-blowing-up}을 보라. 구성상 $X'$는 부풀리기 대수
$A[\frac{\mathfrak m}{\pi}]$, $A[\frac{\mathfrak m}{x}]$,
$A[\frac{\mathfrak m}{y}]$에 대응하는 세 아핀 조각으로 덮인다.

\medskip\noindent
대수 $A[\frac{\mathfrak m}{\pi}]$는 생성원 $x' = x/\pi$와
$y' = y/\pi$를 가지며 $x'y' = \pi^{n - 2}$이다.
따라서 $X'$의 이 부분은 $R[x', y'](x'y' - \pi^{n - 2})$의 스펙트럼이다.

\medskip\noindent
대수 $A[\frac{\mathfrak m}{x}]$는 관계식 $xu - \pi$를 만족하는
생성원 $x$, $u = \pi/x$를 가진다. 이 환은
$y/x = \pi^n/x^2 = u^2\pi^{n - 2}$를 포함한다. 따라서 $X'$의 이 부분은
정칙이다.

\medskip\noindent
대칭성에 의해 대수 $A[\frac{\mathfrak m}{y}]$의 경우도
$A[\frac{\mathfrak m}{x}]$의 경우와 같다.

\medskip\noindent
따라서 $X' \to S$는 상대차원 $1$의 마디 이하 특이성을 가지며,
$X'$는 한 점을 제외하면 정칙이다. 이 한 점은 위와 정확히 같은 아핀
열린 근방을 가지되 $n$이 $n - 2$로 바뀐다. $n$에 대한 귀납법을 쓰면
닫힌점들에서의 부풀리기들의 열
$$
X_{\lfloor n/2 \rfloor} \to \ldots \to X_1 \to X_0 = X
$$
이 존재하여 $X_{\lfloor n/2 \rfloor} \to S$는 상대차원 $1$의
마디 이하 특이성을 가지고 $X_{\lfloor n/2 \rfloor}$은 정칙이다.
\end{example}

\begin{lemma}
\label{models-lemma-etale-local-at-worst-nodal}
$R$을 이산 값매김환이라 하고, $X$를 $R$ 위에서 상대차원 $1$의
마디 이하 특이성을 갖는 스킴이라 하자. $x \in X$를 $R$ 위의 $X$의
특수 올의 한 점이라 하자. 그러면 가환 그림
$$
\xymatrix{
X \ar[d] &
U \ar[r] \ar[d] \ar[l] &
\Spec(A) \ar[dl] \\
\Spec(R) &
\Spec(R') \ar[l]
}
$$
이 존재한다. 여기서 $R \subset R'$은 이산 값매김환의 \'etale 확대이고,
사상 $U \to X$와 $U \to \Spec(A)$는 \'etale이며, $x$로 사상되는
점 $x' \in U$가 존재하고,
$$
A = R'[u, v]/(uv)
\quad\text{또는}\quad
A = R'[u, v]/(uv - \pi^n)
$$
이다. 여기서 $n \geq 0$이고 $\pi \in R'$은 균등화 원소이다.
\end{lemma}

\begin{proof}
이 보조정리는 훨씬 더 일반적인 형태로 이미 증명했다. Algebraic Curves,
보조정리 \ref{curves-lemma-etale-local-structure-nodal-family}를 보라.
여기서는 그곳의 명제를 위의 형태로 옮기기만 하면 된다.

\medskip\noindent
먼저 사상 $X \to \Spec(R)$가 $x$에서 매끄러우면, $x$의 어떤 아핀
열린 근방 $U \subset X$에 대하여 \'etale 사상
$U \to \mathbf{A}^1_R = \Spec(R[u])$를 찾을 수 있다. 이는 Morphisms,
보조정리 \ref{morphisms-lemma-smooth-etale-over-affine-space}이다.
필요하면 좌표 $u$를 $u + 1$로 바꾸어, $x$가 표준 열린집합
$D(u) \subset \mathbf{A}^1_R$의 한 점으로 사상된다고 가정할 수 있다.
그러면 $A = R[u, v]/(uv - 1)$에 대하여 $D(u) = \Spec(A)$이므로
이 경우 결과가 성립한다.

\medskip\noindent
다음으로 $x$가 올의 특이점이라고 하자. Algebraic Curves, 보조정리
\ref{curves-lemma-etale-local-structure-nodal-family}을 적용하면 그림
$$
\xymatrix{
X \ar[d] &
U \ar[rr] \ar[l] \ar[rd] & &
W \ar[r] \ar[ld] &
\Spec(\mathbf{Z}[u, v, a]/(uv - a)) \ar[d] \\
\Spec(R) & &
V \ar[ll] \ar[rr] & & \Spec(\mathbf{Z}[a])
}
$$
을 얻으며, 이는 인용한 보조정리의 명제에 언급된 모든 성질을 가진다.
$x' \in U$를 그 보조정리가 보장하는, $x$로 사상되는 점이라 하자.
먼저 $V$를 $x'$의 상의 아핀 근방으로 줄인다. $V = \Spec(R')$라 하자.
그러면 $R \to R'$은 \'etale이다. $R$이 이산 값매김환이므로 $R'$은
준국소 Dedekind 정역들의 유한 곱이다(More on Algebra, 보조정리
\ref{more-algebra-lemma-Dedekind-etale-extension}을 사용하라).
따라서 예컨대 소 아이디얼 회피를 사용하면, $x'$의 상을 포함하고
$R'_f$가 이산 값매김환이 되는 표준 열린집합
$D(f) \subset V = \Spec(R')$를 찾을 수 있다. $R'$을 $R'_f$로 바꾸면
$V = \Spec(R')$이고 $R \subset R'$이 이산 값매김환의 \'etale 확대인
상황에 이른다. 이산 값매김환의 확대는 More on Algebra, 정의
\ref{more-algebra-definition-extension-discrete-valuation-rings}에서
정의하였다.

\medskip\noindent
사상 $V \to \Spec(\mathbf{Z}[a])$는 $R'$에서 $a$의 상 $h$로 결정된다.
그러면 $W = \Spec(R'[u, v]/(uv - h))$이다. 따라서
$A = R'[u, v]/(uv - h)$로 두면 보조정리가 성립한다.
$h = 0$이면 보조정리에 언급된 첫 번째 경우를 얻음은 분명하다.
$h \not = 0$이면 $R'$의 어떤 단원 $\epsilon$과 $n \geq 0$에 대하여
$h = \epsilon \pi^n$으로 쓸 수 있다. 좌표를
$u_{new} = \epsilon u$, $v_{new} = v$로 바꾸면 보조정리에 열거된
$A$의 두 번째 동형형을 얻는다.
\end{proof}

\begin{lemma}
\label{models-lemma-blowup-at-worst-nodal}
$R$을 이산 값매김환이라 하자. $X$를 $R$ 위에서 상대차원 $1$의
마디 이하 특이성을 가지며 일반 올이 매끄러운 준콤팩트 스킴이라 하자.
그러면 $m \geq 0$과 열
$$
X_m \to \ldots \to X_1 \to X_0 = X
$$
이 존재하여 다음을 만족한다.
\begin{enumerate}
\item $X_{i + 1} \to X_i$는 $X_i$의 특이한 닫힌점 $x_i$에서의 부풀리기이다.
\item $X_i \to \Spec(R)$는 상대차원 $1$의 마디 이하 특이성을 가진다.
\item $X_m$은 정칙이다.
\end{enumerate}
\end{lemma}

\noindent
조금 더 강한 참인 명제는 다음과 같다. 특이점들을 어떤 순서로
부풀리더라도 결국 특이점 해소에 이르며, 중간의 모든 부풀리기는
$R$ 위에서 상대차원 $1$의 마디 이하 특이성을 가진다.

\begin{proof}
$X$가 준콤팩트이므로 특수 올 $X_k$도 준콤팩트이다. $X_k$의 특이점들은
마디 이하이므로 $X_k$는 유한 개의 마디를 가지며 그 밖에서는 $k$ 위에서
매끄럽다. $X \to \Spec(R)$가 평탄하고 일반 올이 매끄러우므로,
$X$는 $X_k$의 유한 개 마디를 제외하면 $R$ 위에서 매끄럽다
(Morphisms, 보조정리 \ref{morphisms-lemma-smooth-at-point}을 사용하라).
따라서 $X$는 특수 올의 마디들을 제외한 모든 점에서 정칙이다
(Algebra, 보조정리 \ref{algebra-lemma-regular-goes-up}을 보라).
$x \in X$를 그러한 마디라 하고, 보조정리
\ref{models-lemma-etale-local-at-worst-nodal}에서와 같은 그림
$$
\xymatrix{
X \ar[d] &
U \ar[r] \ar[d] \ar[l] &
\Spec(A) \ar[dl] \\
\Spec(R) &
\Spec(R') \ar[l]
}
$$
을 택하자. $A = R'[u, v]/(uv)$인 경우는 일어날 수 없다. 이는 $X/R$의
일반 올이 특이하다는 뜻이 되기 때문이다(사소한 세부사항은 생략한다).
따라서 어떤 $n \geq 0$에 대하여 $A = R'[u, v]/(uv - \pi^n)$이다.
$x$가 특이점이므로 $n \geq 2$이다. 예 \ref{models-example-blowup}의 논의를 보라.

\medskip\noindent
$U$를 줄인 뒤에는 $x$로 사상되는 유일한 점 $u \in U$가 있다고 가정할
수 있다. $w \in \Spec(A)$를 $u$의 상이라 하자. 또한 $u$가 $w$로
사상되는 $U$의 유일한 점이라고 가정할 수 있다. 두 수평 화살표가
\'etale이므로, $U$의 닫힌 부분스킴으로 본 $u$는 $x \in X$의 스킴론적
역상이면서 $w \in \Spec(A)$의 스킴론적 역상이다. 부풀리기는 평탄한
밑변환과 가환하므로(Divisors, 보조정리
\ref{divisors-lemma-flat-base-change-blowing-up}), 가환 그림
$$
\xymatrix{
X' \ar[d] &
U' \ar[l] \ar[d] \ar[r] &
W' \ar[d] \\
X & U \ar[l] \ar[r] & \Spec(A)
}
$$
을 얻는다. 두 정사각형은 카르테시안이고, 수직 화살표들은 각각
$X, U, \Spec(A)$에서 $x,u,w$의 부풀리기이다. 스킴 $W'$는 예
\ref{models-example-blowup}에서 기술하였다. 그곳에서 $W'$가 $R'$ 위에서
상대차원 $1$의 마디 이하 특이성을 가짐을 보였다. 따라서 $W'$는 $R$
위에서도 상대차원 $1$의 마디 이하 특이성을 가진다
(Algebraic Curves, 보조정리
\ref{curves-lemma-nodal-family-postcompose-etale}). 그러므로 $U'$도
$R$ 위에서 상대차원 $1$의 마디 이하 특이성을 가진다
(Algebraic Curves, 보조정리
\ref{curves-lemma-nodal-family-etale-local-source}를 보라).
$X' \to X$가 $x$의 여집합 위에서 동형이므로, 보조정리
\ref{curves-lemma-nodal-family-etale-local-source}를 다시 적용하면
$X'/R$도 그러함을 얻는다.

\medskip\noindent
끝으로 이 부풀리기를 유한 번 수행하면 정칙 모형 $X_m$에 이른다는 것을
보여야 한다. 위에서 기술한 부풀리기 아래에서 ``불변량'' $n$이 $2$씩
감소하므로 이는 상당히 분명하다. 예 \ref{models-example-blowup}의 계산을 보라.
그러나 이 불변량을 엄밀히 정의하고 성질을 확립하는 일을 피하기 위해
다음과 같이 논증한다. $n = 2$이면 $W'$는 정칙이고, 따라서 $X'$는
$x$ 위에 놓인 모든 점에서 정칙이어서 $X$의 특이점 개수가 $1$ 줄었다.
$n > 2$이면 $w$ 위에 놓인 $W'$의 유일한 특이점 $w'$는
$\kappa(w) = \kappa(w')$를 만족한다. 따라서 $U'$에는 $u$ 위에 놓인
유일한 특이점 $u'$가 있고 $\kappa(u) = \kappa(u')$이다.
이는 $X'$에 $x$ 위에 놓인 유일한 특이점 $x'$가 있음을 뜻하며,
그 점은 바로 $u'$의 상이다. 그러므로 위와 정확히 같은 논증으로 가환 그림
$$
\xymatrix{
X'' \ar[d] &
U'' \ar[l] \ar[d] \ar[r] &
W'' \ar[d] \\
X' & U' \ar[l] \ar[r] & W'
}
$$
을 얻는다. 두 정사각형은 카르테시안이고, 수직 화살표들은 각각
$X',U',W'$에서 $x',u',w'$의 부풀리기이다. 이를 계속하면
$\lfloor n/2 \rfloor$단계 후에 끝나는 서로 양립하는 부풀리기들의 열을
얻는다. 이 과정이 끝나면 스킴 $X^{(\lfloor n/2 \rfloor)}$은 $X$보다
특이점이 하나 적다. 특이점의 개수에 대한 귀납법으로 증명이 끝난다.
\end{proof}

\begin{lemma}
\label{models-lemma-blowdown-at-worst-nodal}
$R$을 분수체 $K$와 잉여체 $k$를 갖는 이산 값매김환이라 하자.
$X \to \Spec(R)$가 $R$ 위에서 상대차원 $1$의 마디 이하 특이성을
가진다고 하자. $X \to X'$가 제1종 예외곡선 $E \subset X$의
수축이라 하자. 그러면 $X'$는 $R$ 위에서 상대차원 $1$의 마디 이하
특이성을 가진다.
\end{lemma}

\begin{proof}
$x' \in X'$를 $E$의 상이라 하자. 확인할 것은
$X' \to \Spec(R)$가 $x'$의 근방에서 상대차원 $1$의 마디 이하
특이성을 가진다는 것뿐이다. $X \to \Spec(R)$의 닫힌 올은 축소되어
있으므로 $\pi \in R$은 $E$ 위에서 차수 $1$로 소멸한다. 이는 곧
$\pi$가 $\mathfrak m_{x'} \subset \mathcal{O}_{X', x'}$의 원소이지만
$\mathfrak m_{x'}^2$에는 속하지 않음을 뜻한다.
$\mathcal{O}_{X', x'}$는 차원 $2$의 정칙환이므로(Resolution of Surfaces,
절 \ref{resolve-section-minus-one}의 수축 정의에 의해),
$\mathcal{O}_{X'_k, x'}$는 차원 $1$의 정칙환이다
(Algebra, 보조정리 \ref{algebra-lemma-regular-ring-CM}).
한편 곡선 $E$는 닫힌 올 $X_k$의 다른 성분 하나 이상과 만나야 한다.
그 성분을 $C$라 하고 $x \in E \cap C$라 하자. 그러면 $x$는 특수 올
$X_k$의 마디이고, 따라서 $\kappa(x)/k$는 유한 분리 확대이다
(Algebraic Curves, 보조정리 \ref{curves-lemma-nodal}).
$x \mapsto x'$이므로 $\kappa(x')/k$도 유한 분리 확대이다.
Algebra, 보조정리 \ref{algebra-lemma-separable-smooth}에 의해
$X'_k \to \Spec(k)$는 $x'$의 어떤 열린 근방에서 매끄럽다.
평탄성과 합치면 $X' \to \Spec(R)$도 $x'$의 근방에서 매끄럽다
(Morphisms, 보조정리 \ref{morphisms-lemma-smooth-at-point}).
상대차원 $1$의 매끄러운 사상은 상대차원 $1$의 마디 이하 특이성을
가지므로 증명이 끝난다(Algebraic Curves, 보조정리
\ref{curves-lemma-smooth-relative-dimension-1}).
\end{proof}

\begin{lemma}
\label{models-lemma-semistable}
$R$을 분수체 $K$를 갖는 이산 값매김환이라 하고,
$C$를 $H^0(C, \mathcal{O}_C) = K$인 $K$ 위의 매끄러운 사영곡선이라 하자.
다음은 서로 동치이다.
\begin{enumerate}
\item $R$ 위에서 상대차원 $1$의 마디 이하 특이성을 갖는 $C$의
고유 모형이 존재한다.
\item $R$ 위에서 상대차원 $1$의 마디 이하 특이성을 갖는 $C$의
극소 모형이 존재한다.
\item $C$의 모든 극소 모형은 $R$ 위에서 상대차원 $1$의 마디 이하
특이성을 갖는다.
\end{enumerate}
\end{lemma}

\begin{proof}
이 명제를 이해하기 위해 다음을 상기하자. 극소 모형은 제1종 예외곡선을
갖지 않는 고유 정칙 모형으로 정의되고(정의
\ref{models-definition-minimal-model}), 극소 모형은 존재하며(명제
\ref{models-proposition-exists-minimal-model}), $C$의 종수가 $>0$이면
극소 모형은 유일하다(보조정리 \ref{models-lemma-minimal-model-unique}).
이를 염두에 두면 (2) $\Rightarrow$ (1)과 (3) $\Rightarrow$ (2)는
분명하다.

\medskip\noindent
(1)을 가정하자. $X$를 $R$ 위에서 상대차원 $1$의 마디 이하 특이성을
갖는 $C$의 고유 모형이라 하자. 보조정리
\ref{models-lemma-blowup-at-worst-nodal}을 적용하면 $X$가 정칙이라고도
가정할 수 있다. 보조정리 \ref{models-lemma-pre-exists-minimal-model}에서와 같이
$$
X = X_m \to X_{m - 1} \to \ldots \to X_1 \to X_0
$$
라 하자. 보조정리 \ref{models-lemma-blowdown-at-worst-nodal}과 귀납법에 의해
$X_0$는 $R$ 위에서 상대차원 $1$의 마디 이하 특이성을 가진다.

\medskip\noindent
증명을 끝내려면 (2)가 (3)을 함의함을 보여야 한다. $C$의 종수가
$>0$이면 극소 모형이 유일하므로 이는 분명하다(위 논의를 보라).
한편 극소 모형이 유일하지 않으면, 보조정리 \ref{models-lemma-nonuniqueness}에
의해 모든 극소 모형의 특수 올은 $\mathbf{P}^1_k$와 동형이다.
따라서 모든 극소 모형에 대해 사상 $X \to \Spec(R)$는 매끄럽다.
\end{proof}

\begin{definition}
\label{models-definition-semistable}
$R$을 분수체 $K$를 갖는 이산 값매김환이라 하고,
$C$를 $H^0(C, \mathcal{O}_C) = K$인 $K$ 위의 매끄러운 사영곡선이라 하자.
보조정리 \ref{models-lemma-semistable}의 동치 조건들이 성립할 때
$C$가 {\it 반안정 축약을 가진다}고 한다.
\end{definition}

\begin{lemma}
\label{models-lemma-good}
$R$을 분수체 $K$를 갖는 이산 값매김환이라 하고,
$C$를 $H^0(C, \mathcal{O}_C) = K$인 $K$ 위의 매끄러운 사영곡선이라 하자.
다음은 서로 동치이다.
\begin{enumerate}
\item $C$의 매끄러운 고유 모형이 존재한다.
\item $R$ 위에서 매끄러운 $C$의 극소 모형이 존재한다.
\item 모든 극소 모형은 $R$ 위에서 매끄럽다.
\end{enumerate}
\end{lemma}

\begin{proof}
$X$가 매끄러운 고유 모형이면 특수 올은 연결이고
(보조정리 \ref{models-lemma-regular-model-connected}) 매끄러우므로 기약이다.
이는 곧 그 모형이 극소임을 뜻한다. 따라서 (1)은 (2)를 함의한다.
증명을 끝내려면 (2)가 (3)을 함의함을 보여야 한다.
$C$의 종수가 $>0$이면 극소 모형이 유일하므로 이는 분명하다
(보조정리 \ref{models-lemma-minimal-model-unique}). 한편 극소 모형이 유일하지
않으면 보조정리 \ref{models-lemma-nonuniqueness}에 의해 특수 올이
$\mathbf{P}^1_k$와 동형이므로, 모든 극소 모형에 대해
$X \to \Spec(R)$는 매끄럽다.
\end{proof}

\begin{definition}
\label{models-definition-good}
$R$을 분수체 $K$를 갖는 이산 값매김환이라 하고,
$C$를 $H^0(C, \mathcal{O}_C) = K$인 $K$ 위의 매끄러운 사영곡선이라 하자.
보조정리 \ref{models-lemma-good}의 동치 조건들이 성립할 때
$C$가 {\it 양호한 축약을 가진다}고 한다.
\end{definition}








\section{종수 0에서의 반안정 축약}
\label{models-section-semistable-reduction-genus-zero}

\noindent
이 절에서는 종수 0인 곡선에 대한 반안정 축약 정리
(정리 \ref{models-theorem-semistable-reduction})를 증명한다.

\medskip\noindent
$R$을 분수체 $K$를 갖는 이산 값매김환이라 하고,
$C$를 $H^0(C, \mathcal{O}_C) = K$인 $K$ 위의 매끄러운 사영곡선이라 하자.
$C$의 종수가 $0$이면 $C$는 이차곡선과 동형이다
(Algebraic Curves, 보조정리 \ref{curves-lemma-genus-zero}).
따라서 $C(K') \not = \emptyset$이 되는 차수 $2$ 이하의 유한 분리 확대
$K'/K$가 존재한다. Algebraic Curves, 보조정리
\ref{curves-lemma-smooth-plane-curve-point-over-separable}를 보라.
$R' \subset K'$을 $R$의 정수적 폐포라 하자. More on Algebra, 주석
\ref{more-algebra-remark-finite-separable-extension}의 논의를 보라.
$R'$의 각 극대 아이디얼 $\mathfrak m$에 대해 $C_{K'}$가
$R'_{\mathfrak m}$ 위에서 반안정 축약을 가짐을 보이겠다
(물론 현재의 경우 그러한 아이디얼은 많아야 두 개이다).
$R$을 $R'_{\mathfrak m}$으로, $C$를 $C_{K'}$로 바꾸면 다음 문단에서
논의할 경우로 환원된다.

\medskip\noindent
이 문단에서 $R$은 분수체 $K$를 갖는 이산 값매김환이고,
$C$는 $H^0(C, \mathcal{O}_C) = K$이며 종수가 $0$인 $K$ 위의 매끄러운
사영곡선이고, $C$는 $K$-유리점을 가진다. 이 경우 Algebraic Curves,
명제 \ref{curves-proposition-projective-line}에 의해
$C \cong \mathbf{P}^1_K$이다. 따라서 $\mathbf{P}^1_R$을 모형으로 쓸 수
있고, $C$는 양호한 축약과 반안정 축약을 모두 가진다.

\begin{example}
\label{models-example-extension-necessary-genus-zero}
$R = \mathbf{R}[[\pi]]$라 하고 스킴
$$
X = V(T_1^2 + T_2^2 - \pi T_0^2) \subset \mathbf{P}^2_R
$$
을 생각하자. $\mathbf{C}$ 위에서
$T_1^2 + T_2^2 = (T_1 + iT_2)(T_1 - iT_2)$라는 분해가 있으므로,
$X$를 $\mathbf{C}[[\pi]]$로 밑변환한 것은 예
\ref{models-example-nonunique-in-genus-zero}에서 정의한 스킴과 동형이다.
따라서 $X$는 정칙이고 그 특수 올은 기약이지만 특이하다.
그러므로 $X$는 그 일반 올의 유일한 극소 모형이다
(보조정리 \ref{models-lemma-nonuniqueness}을 사용하라).
따라서 종수 $0$에서도 확대가 필요하다.
\end{example}



\section{종수 1에서의 반안정 축약}
\label{models-section-semistable-reduction-genus-one}

\noindent
이 절에서는 종수 1인 곡선에 대한 반안정 축약 정리
(정리 \ref{models-theorem-semistable-reduction})를 증명한다.
독자는 먼저 종수 $\geq 2$인 경우의 증명
(절 \ref{models-section-semistable-reduction-genus-at-least-two})을 읽는 것이 좋다.
보조정리 \ref{models-lemma-genus-one}에서 주어진 종수 $1$의 극소 수치형의
분류를 최대한 활용할 것이다.

\medskip\noindent
분수체가 $K$인 이산 값매김환 $R$과,
$H^0(C, \mathcal{O}_C) = K$인 $K$ 위의 매끄러운 사영곡선 $C$를 택하자.
$C$의 종수가 $1$이라고 가정하자. $k$의 표수와 다른 소수
$\ell \geq 7$을 택한다. $C(K') \not = \emptyset$이고 다음을 만족하는
유한 분리 확대 $K'/K$를 택한다.
$\Pic(C_{K'})[\ell] \cong (\mathbf{Z}/\ell \mathbf{Z})^{\oplus 2}$.
Algebraic Curves, 보조정리
\ref{curves-lemma-torsion-picard-becomes-visible}를 보라.
$R' \subset K'$을 $R$의 정수적 폐포라 하자. More on Algebra, 주석
\ref{more-algebra-remark-finite-separable-extension}의 논의를 보라.
$R$을 $R'$의 어떤 극대 아이디얼 $\mathfrak m$에 대한
$R'_{\mathfrak m}$으로, $C$를 $C_{K'}$로 바꿀 수 있다.
이로써 다음 문단에서 다룰 경우로 환원된다.

\medskip\noindent
이 절의 나머지 부분에서 $R$은 분수체 $K$를 갖는 이산 값매김환이고,
$C$는 $H^0(C, \mathcal{O}_C) = K$이며 종수가 $1$인 $K$ 위의 매끄러운
사영곡선으로서 $K$-유리점을 가지고,
$\Pic(C)[\ell] \cong (\mathbf{Z}/\ell \mathbf{Z})^{\oplus 2}$
을 만족한다고 하자. 여기서 $\ell \geq 7$은 $k$의 표수와 다른 소수이다.
$C$가 반안정 축약을 가짐을 증명하겠다.

\medskip\noindent
$X$를 $C$의 극소 모형이라 하자. 명제
\ref{models-proposition-exists-minimal-model}을 보라.
$T = (n, m_i, (a_{ij}), w_i, g_i)$를 $X$에 결부된 수치형이라 하자
(정의 \ref{models-definition-numerical-type-model}). 그러면 $T$는 극소
수치형이다(보조정리 \ref{models-lemma-numerical-type-minimal-model}).
$C$가 유리점을 가지므로 보조정리
\ref{models-lemma-numerical-type-rational-point}에 의해
$m_i = w_i = 1$인 $i$가 존재한다. 보조정리 \ref{models-lemma-genus-one}의
종수 $1$인 극소 수치형의 분류를 살펴보면 $m = w = 1$이고, 경우
(\ref{models-item-up4}),
(\ref{models-item-equal-down3}),
(\ref{models-item-up-up}),
(\ref{models-item-down-up}),
(\ref{models-item-up-equal-up}),
(\ref{models-item-down-equal-up}),
(\ref{models-item-triple-with-down}),
(\ref{models-item-equal-equal-down-equal}),
(\ref{models-item-up-equal-equal-up}),
(\ref{models-item-down-equal-equal-up}),
(\ref{models-item-triple-extended-down}),
(\ref{models-item-up-chain-equal-up}),
(\ref{models-item-down-chain-equal-up}),
(\ref{models-item-Dn-extended-down})는 불가능하다. 그 경우에는
$w_i$와 $m_i$가 모두 $1$인 첨자가 없기 때문이다.
$e$를 $i < j$이고 $a_{ij} > 0$인 쌍 $(i,j)$의 개수라 하자.
나머지 경우에는 다음이 성립한다.
\begin{enumerate}
\item[(A)] 다음 경우에는 $e = n - 1$이다.
(\ref{models-item-one}),
(\ref{models-item-two-cycle}),
(\ref{models-item-equal-up3}),
(\ref{models-item-up-down}),
(\ref{models-item-up-equal-down}),
(\ref{models-item-triple-with-up}),
(\ref{models-item-equal-equal-up-equal}),
(\ref{models-item-up-equal-equal-down}),
(\ref{models-item-quadruple}),
(\ref{models-item-triple-extended-up}),
(\ref{models-item-up-chain-equal-down}),
(\ref{models-item-Dn-extended-up}),
(\ref{models-item-double-triple}),
(\ref{models-item-E6-completed}),
(\ref{models-item-E7-completed}),
(\ref{models-item-E8-completed}).
\item[(B)] 다음 경우에는 $e = n$이다.
(\ref{models-item-three-cycle}),
(\ref{models-item-four-cycle}),
(\ref{models-item-five-cycle}),
(\ref{models-item-n-cycle}).
\end{enumerate}
이 두 경우를 따로 논증한다.

\medskip\noindent
경우 (A). 이 경우 $\Pic(T)[\ell]$은 자명하다. 수치형의 Picard 군은
절 \ref{models-section-picard-group}에서 정의하였다.
$\Pic(T) \subset \Coker(A)$이고(보조정리 \ref{models-lemma-picard-T-and-A}),
$\ell$을 $a_{ij}$ 및 $m_i$와 서로소로 택했으므로 보조정리
\ref{models-lemma-recurring-symmetric-integer}에 의해 $\Coker(A)[\ell] = 0$이다.
따라서 위의 소멸이 성립한다. 보조정리 \ref{models-lemma-sequence-torsion}과
\ref{models-lemma-torsion-embeds}에 의해 매장이 존재한다.
$$
(\mathbf{Z}/\ell \mathbf{Z})^{\oplus 2} \subset \Pic((X_k)_{red})[\ell].
$$
Algebraic Curves, 보조정리
\ref{curves-lemma-bound-torsion-simple}에 의해
$$
2 \leq \dim_k H^1((X_k)_{red}, \mathcal{O}_{(X_k)_{red}}) +
g_{geom}((X_k)_{red}/k)
$$
을 얻는다. Algebraic Curves, 보조정리
\ref{curves-lemma-bound-geometric-genus}과
\ref{curves-lemma-bound-geometric-genus-curve}에 의해
$g_{geom}((X_k)_{red}/k) \leq \sum w_ig_i$이다.
보조정리 \ref{models-lemma-numerical-type-rational-point}에 의해 보조정리
\ref{models-lemma-genus-reduction-smaller}의 가정이 성립하므로
$\dim_k H^1((X_k)_{red}, \mathcal{O}_{(X_k)_{red}}) \leq g = 1$이다.
이들을 합치면
$$
2 \leq 1 + \sum w_i g_i
$$
이다. 목록을 살펴보면 수치형은
$n = 1$, $w_1 = m_1 = g_1 = 1$로 주어짐을 얻는다.
모든 곳에서 등호가 성립하므로 $g_{geom}(C_1/k) = 1$이다.
한편 $C_1$은 $C_1 \to \Spec(k)$가 $x$에서 매끄럽게 되는
$k$-유리점 $x$를 가진다. 따라서 $C_1$은 기하적으로 정역이다
(Varieties, 보조정리
\ref{varieties-lemma-variety-with-smooth-rational-point}).
그러므로 $g_{geom}(C_1/k) = 1$은 $C_{1,\overline{k}}$의 정규화의
종수와 $C_{1,\overline{k}}$의 종수 모두와 같다. 따라서 정규화 사상
$C_{1, \overline{k}}^\nu \to C_{1, \overline{k}}$는 동형이다
(Algebraic Curves, 보조정리 \ref{curves-lemma-genus-normalization}).
원하는 대로 $C_1$은 $k$ 위에서 매끄럽다.

\medskip\noindent
경우 (B). 여기서는 매장
$$
\mathbf{Z}/\ell \mathbf{Z} \subset \Pic(X_k)[\ell]
$$
만을 얻는다. 수치형의 분류에서 각 $i$에 대하여
$m_i = w_i = 1$이고 $g_i = 0$임을 알 수 있다. 따라서 각 $C_i$는
$k$ 위의 종수 0인 곡선이다. 더욱이 각 $i$에 대해
$C_i \cap C_j$가 $k$-유리점이 되는 $j$가 존재한다. 따라서
Algebraic Curves, 명제 \ref{curves-proposition-projective-line}에 의해
$C_i \cong \mathbf{P}^1_k$이다. 특히 $X_k$는 $C_i$들의 스킴론적
합집합이므로, $X_{\overline{k}}$는 $C_{i,\overline{k}}$들의 스킴론적
합집합이다. 따라서 $X_{\overline{k}}$는 $\overline{k}$ 위의 차원 $1$인
축소되고 연결된 고유 스킴이며
$\dim_{\overline{k}} H^1(X_{\overline{k}},
\mathcal{O}_{X_{\overline{k}}}) = 1$이다. 또한 Varieties, 보조정리
\ref{varieties-lemma-change-fields-pic}와 위 사실에 의해 여전히
% P07-KO55-ERR-002: frozen source omits [\ell]) in this dimension expression; target preserves canonical bytes pending canon adjudication.
$$
\dim_{\mathbf{F}_\ell}(\Pic(X_{\overline{k}}) \geq 1
$$
이다. Algebraic Curves, 명제
\ref{curves-proposition-torsion-picard-reduced-proper}에 의해
$X_{\overline{k}}$는 다중교차 특이점 이하만 가진다. 그런데 $X_k$는
Gorenstein이므로(보조정리 \ref{models-lemma-gorenstein}),
$X_{\overline{k}}$도 Gorenstein이다(Duality for Schemes, 보조정리
\ref{duality-lemma-gorenstein-base-change}). 따라서 Algebraic Curves,
보조정리 \ref{curves-lemma-multicross-gorenstein-is-nodal}에 의해
$X_{\overline{k}}$는 마디 이하 특이성을 가진다. 이로써 경우 (B)의
증명이 끝난다.

\begin{example}
\label{models-example-extension-necessary-genus-one}
$k$를 대수적으로 닫힌 체라 하고, $Z$를 $k$ 위의 양의 종수 $g$인
매끄러운 사영곡선이라 하자. $n \geq 1$을 $k$의 표수와 서로소인
정수라 하자. $\mathcal{L}$을 위수가 $n$인 가역
$\mathcal{O}_Z$-가군이라 하자. Algebraic Curves, 보조정리
\ref{curves-lemma-torsion-picard-smooth-projective}를 보라.
동형 $\varphi : \mathcal{L}^{\otimes n} \to \mathcal{O}_Z$를 택하자.
$R = k[[\pi]]$라 하고 그 분수체를 $K = k((\pi))$라 하자.
$Z_R$로 $Z$의 $R$로의 밑변환을 나타내고, $\mathcal{L}_R$을
$\mathcal{L}$의 $Z_R$로의 당김이라 하자. 유한 평탄 사상
$$
p : X \longrightarrow Z_R
$$
중 다음을 만족하는 것을 생각하자.
$$
p_*\mathcal{O}_X =
\text{Sym}^*_{\mathcal{O}_{Z_R}}(\mathcal{L}_R)/(\varphi - \pi) =
\mathcal{O}_{Z_R} \oplus \mathcal{L}_R \oplus
\mathcal{L}_R^{\otimes 2} \oplus \ldots \oplus \mathcal{L}_R^{\otimes n - 1}
$$
더 정확히 말해 $U = \Spec(A) \subset Z$가 아핀 열린집합이고,
$\varphi(s^{\otimes n}) = f$인 절단 $s$가 $\mathcal{L}|_U$를
자명화한다고 하자. 여기서 $f$는 단원이다. 그러면
$$
p^{-1}(U_R) = \Spec\left(
(A \otimes_R R[[\pi]])[x]/(x^n - \pi f)
\right)
$$
독자는 일반 올 사이의 사상 $X_K \to Z_K$가 유한 \'etale임을 확인할
수 있다. 구조층의 기술에서 $H^0(X, \mathcal{O}_X) = R$이고
$H^0(X_K, \mathcal{O}_{X_K}) = K$임을 알 수 있다.
Riemann--Hurwitz 공식(Algebraic Curves, 보조정리
\ref{curves-lemma-rhe})에 의해 $X_K$의 종수는 $n(g - 1) + 1$이다.
특히 $Z$의 종수가 $1$이면 $X_K$의 종수도 $1$이다.
한편 위의 국소 방정식에 의해 스킴 $X$는 정칙이고, 유효 Cartier
인자로서 특수 올 $X_k$는 그 축소 특수 올의 $n$배이다.
따라서 $X_K$가 반안정 축약을 갖게 되는 모든 유한 확대 $K'/K$는
분기지수 적어도 $n$으로 분기해야 한다(일부 세부사항은 생략한다).
그러므로 종수 $1$인 곡선이 반안정 축약을 갖게 되는 확대의 차수에는
보편적인 상계가 존재하지 않는다.
\end{example}






\section{종수 2 이상에서의 반안정 축약}
\label{models-section-semistable-reduction-genus-at-least-two}

\noindent
이 절에서는 종수 $\geq 2$인 곡선에 대한 반안정 축약 정리
(정리 \ref{models-theorem-semistable-reduction})를 증명한다.
$g \geq 2$를 고정하자.

\medskip\noindent
분수체가 $K$인 이산 값매김환 $R$과,
$H^0(C, \mathcal{O}_C) = K$인 $K$ 위의 매끄러운 사영곡선 $C$를 택하자.
$C$의 종수가 $g$라고 가정하자. $k$의 표수와 다른 소수
$\ell > 768g$를 택한다. $C(K') \not = \emptyset$이고 다음을 만족하는
유한 분리 확대 $K'/K$를 택한다.
$\Pic(C_{K'})[\ell] \cong (\mathbf{Z}/\ell \mathbf{Z})^{\oplus 2g}$.
Algebraic Curves, 보조정리
\ref{curves-lemma-torsion-picard-becomes-visible}를 보라.
$R' \subset K'$을 $R$의 정수적 폐포라 하자. More on Algebra, 주석
\ref{more-algebra-remark-finite-separable-extension}의 논의를 보라.
$R$을 $R'$의 어떤 극대 아이디얼 $\mathfrak m$에 대한
$R'_{\mathfrak m}$으로, $C$를 $C_{K'}$로 바꿀 수 있다.
이로써 다음 문단에서 다룰 경우로 환원된다.

\medskip\noindent
이 절의 나머지 부분에서 $R$은 분수체 $K$를 갖는 이산 값매김환이고,
$C$는 $H^0(C, \mathcal{O}_C) = K$이며 종수가 $g$인 $K$ 위의 매끄러운
사영곡선으로서 $K$-유리점을 가지고,
$\Pic(C)[\ell] \cong (\mathbf{Z}/\ell \mathbf{Z})^{\oplus 2g}$
을 만족한다고 하자. 여기서 $\ell \geq 768g$는 $k$의 표수와 다른
소수이다. $C$가 반안정 축약을 가짐을 증명하겠다.

\medskip\noindent
이 절의 나머지에서는 보조정리
\ref{models-lemma-numerical-type-rational-point}의 결론들을 별도의 언급 없이
사용한다.

\medskip\noindent
$X$를 $C$의 극소 모형이라 하자. 명제
\ref{models-proposition-exists-minimal-model}을 보라.
$T = (n, m_i, (a_{ij}), w_i, g_i)$를 $X$에 결부된 수치형이라 하자
(정의 \ref{models-definition-numerical-type-model}). 그러면 $T$는 종수 $g$의
극소 수치형이다(보조정리 \ref{models-lemma-numerical-type-minimal-model}).
명제 \ref{models-proposition-bound-picard-group}에 의해
$$
\dim_{\mathbf{F}_\ell} \Pic(T)[\ell] \leq g_{top}
$$
이다. 보조정리 \ref{models-lemma-sequence-torsion}과
\ref{models-lemma-torsion-embeds}에 의해 매장
$$
(\mathbf{Z}/\ell \mathbf{Z})^{\oplus 2g - g_{top}}
\subset \Pic((X_k)_{red})[\ell].
$$
이 존재한다. Algebraic Curves, 보조정리
\ref{curves-lemma-bound-torsion-simple}에 의해
$$
2g - g_{top} \leq
\dim_k H^1((X_k)_{red}, \mathcal{O}_{(X_k)_{red}}) + g_{geom}(X_k/k)
$$
을 얻는다. 보조정리 \ref{models-lemma-genus-reduction-smaller}와
\ref{models-lemma-genus-reduction-bigger-than}에 의해
$$
g \geq \dim_k H^1((X_k)_{red}, \mathcal{O}_{(X_k)_{red}}) \geq
g_{top} + g_{geom}(X_k/k)
$$
이다. 초등 정수론에 따르면 이 $3$개의 부등식이 동시에 성립할 수 있는
유일한 경우는 모두 등식인 경우이다. 보조정리
\ref{models-lemma-genus-reduction-smaller}를 살펴보면 모든 $i$에 대해
$m_i = 1$임을 얻는다. 보조정리
\ref{models-lemma-equality-genus-reduction-bigger-than}을 살펴보면 $X_k$의
모든 기약성분이 $k$ 위에서 매끄러움을 얻는다.

\medskip\noindent
특히 $X_k$는 그 기약성분 $C_i$들의 스킴론적 합집합이므로,
$X_{\overline{k}}$는 $C_{i,\overline{k}}$들의 스킴론적 합집합이다.
따라서 $X_{\overline{k}}$는 $\overline{k}$ 위의 차원 $1$인 축소되고
연결된 고유 스킴이며
$\dim_{\overline{k}}H^1(X_{\overline{k}},
\mathcal{O}_{X_{\overline{k}}}) = g$이다.
또한 Varieties, 보조정리 \ref{varieties-lemma-change-fields-pic}와
위 사실에 의해 여전히
$$
\dim_{\mathbf{F}_\ell}(\Pic(X_{\overline{k}})[\ell])
\geq 2g - g_{top} =
\dim_{\overline{k}} H^1(X_{\overline{k}}, \mathcal{O}_{X_{\overline{k}}})
+ g_{geom}(X_{\overline{k}})
$$
이다. Algebraic Curves, 명제
\ref{curves-proposition-torsion-picard-reduced-proper}에 의해
$X_{\overline{k}}$는 다중교차 특이점 이하만 가진다. 그런데 $X_k$는
Gorenstein이므로(보조정리 \ref{models-lemma-gorenstein}),
$X_{\overline{k}}$도 Gorenstein이다(Duality for Schemes, 보조정리
\ref{duality-lemma-gorenstein-base-change}). 따라서 Algebraic Curves,
보조정리 \ref{curves-lemma-multicross-gorenstein-is-nodal}에 의해
$X_{\overline{k}}$는 마디 이하 특이성을 가진다. 이로써 증명이 끝난다.










\section{곡선의 반안정 축약}
\label{models-section-semistable-reduction-theorem}

\noindent
이 절에서는 정리의 증명을 마무리한다. $g \geq 2$에 대하여
$768g < \ell' < \ell$을 $> 768g$인 처음 두 소수라 하고
\begin{equation}
\label{models-equation-bound}
B_g = (2g - 2)(\ell^{2g})!
\end{equation}
라 하자. $B_g$의 정확한 형태는 중요하지 않다. 핵심은 그것이 오직
$g$에만 의존한다는 것이다.

\begin{theorem}
\label{models-theorem-semistable-reduction}
\begin{reference}
\cite[따름정리 2.7]{DM}
\end{reference}
$R$을 분수체 $K$를 갖는 이산 값매김환이라 하고, $C$를
$H^0(C, \mathcal{O}_C) = K$인 $K$ 위의 매끄러운 사영곡선이라 하자.
그러면 분수체의 유한 분리 확대 $K'/K$를 유도하는 이산 값매김환의 확대
$R \subset R'$이 존재하여 $C_{K'}$는 반안정 축약을 가진다.
더 정확히는 다음이 성립한다.
\begin{enumerate}
\item $C$의 종수가 0이면, $C_{K'} \cong \mathbf{P}^1_{K'}$이 되는
차수 $2$의 분리 확대 $K'/K$가 존재한다. 따라서 $C_{K'}$는 $K'$ 안에서
$R$의 정수적 폐포 $R'$ 위의 매끄러운 사영 스킴
$\mathbf{P}^1_{R'}$의 일반 올과 동형이다.
\item $C$의 종수가 1이면, 유한 분리 확대 $K'/K$가 존재하여
$K'$ 안에서 $R$의 정수적 폐포 $R'$의 모든 극대 아이디얼
$\mathfrak m$에 대해 $C_{K'}$는 $R'_\mathfrak m$ 위에서 반안정
축약을 가진다. 더욱이 $R'_\mathfrak m$ 위에서 $C_{K'}$의
(유일한) 극소 모형의 특수 올은 종수 1의 매끄러운 곡선이거나
유리곡선들의 순환이다.
\item $C$의 종수 $g$가 1보다 크면, 차수가 많아야
$B_g$인 유한 분리 확대 $K'/K$가 존재하여
(\ref{models-equation-bound}), $K'$ 안에서 $R$의 정수적 폐포 $R'$의
모든 극대 아이디얼 $\mathfrak m$에 대해 $C_{K'}$는
$R'_\mathfrak m$ 위에서 반안정 축약을 가진다.
\end{enumerate}
\end{theorem}

\begin{proof}
종수 0인 경우는 절 \ref{models-section-semistable-reduction-genus-zero}을,
종수 1인 경우는 절 \ref{models-section-semistable-reduction-genus-one}을,
종수 1보다 큰 경우는 절
\ref{models-section-semistable-reduction-genus-at-least-two}을 보라.
차수 $[K' : K]$의 상계를 얻기 위해서는 모든 $\ell$-꼬임 또는
$\ell'$-꼬임이 드러나게 하는 데 필요한 확대의 차수에 대한 상계를
사용할 수 있다. 이는 Algebraic Curves, 보조정리
\ref{curves-lemma-torsion-picard-becomes-visible}에서 증명하였다.
($\ell$과 $\ell'$을 모두 사용하는 이유는 잉여체 $k$의 표수를
피해야 하기 때문이다.)
\end{proof}

\begin{remark}[상계의 개선]
\label{models-remark-improving-bound}
문헌의 결과들은 정리 \ref{models-theorem-semistable-reduction}의 명제에 제시된
상계를 개선할 수 있음을 시사한다. 예를 들어 \cite{DM}에서는 잉여체가
완전체이면 $C$의 반안정 축약과 그 Jacobian의 반안정 축약이 동치임을
보였으며, 이는 일반 잉여체에 대해서도 성립할 것으로 보인다.
아벨 다양체의 경우, 잉여표수와 다른 임의의 $\ell \geq 3$에 대해
$\ell$-꼬임 위의 Galois 작용이 자명이면 반안정 축약을 가진다.
따라서 $B_g$의 공식 (\ref{models-equation-bound})에서 $\ell = 5$를 택할 수
있을 것으로 보인다. 다만 증명에는 훨씬 더 많은 작업이 필요하다.
이 결과가 필요해지면 정확한 명제를 세우고 여기에서 증명할 것이다.
\end{remark}
